% =============================================================
% New GDP series for interwar and Soviet Estonia:
% a Finnish miracle that never was
%
% Authors: Adomas Klimantas (ISM University of Management and
%          Economics) and Martin Klesment (University of Tartu)
% Target journal: Economic History Review
% Citation style: Oxford footnote (biblatex-chicago notes)
%
% Build (from Paper/estonia/):  xelatex main; biber main; xelatex main; xelatex main
% All tables and the number macros are generated by scripts/R/tables.R
% =============================================================

\documentclass[12pt,a4paper]{article}

% ---- Core packages ----
\usepackage{fontspec}
\usepackage{microtype}
\emergencystretch=2em
\usepackage{geometry}
\geometry{a4paper, margin=2.5cm}

% ---- Spacing (EHR submission: double-spaced) ----
\usepackage{setspace}
\doublespacing

% ---- Fonts ----
\setmainfont{Times New Roman}

% ---- Language ----
\usepackage[main=british]{babel}
\usepackage[autostyle=true]{csquotes}   % follows babel british: single outer quotes (EHR house style)

% ---- Mathematics ----
\usepackage{amsmath}

% ---- Tables ----
\usepackage{booktabs}
\usepackage{threeparttable}
\usepackage{longtable}
\usepackage{multirow}
\usepackage{array}
\usepackage[justification=centering]{caption}

% ---- Footnote formatting ----
\usepackage[bottom]{footmisc}
\usepackage{ragged2e}

% ---- Bibliography: Chicago notes (Oxford footnote style) ----
\usepackage[notes,
            backend=biber,
            footmarkoff,
            isbn=false,
            url=false,
            doi=false]{biblatex-chicago}
\addbibresource{bibliography.bib}
\addbibresource{bibliography_additions.bib}

\ExecuteBibliographyOptions{
    maxbibnames=99,
    maxcitenames=2,
    uniquetitle=true
}

% ---- PDF metadata and links ----
\usepackage{hyperref}
\hypersetup{
    colorlinks=true,
    linkcolor=black,
    citecolor=black,
    urlcolor=black,
    pdftitle={New GDP series for interwar and Soviet Estonia: a Finnish miracle that never was},
    pdfauthor={Adomas Klimantas and Martin Klesment}
}

% ---- Custom commands ----
\newcommand{\intlnineteen}{Int\$\textsubscript{1990}}   % Int$_1990
\newcommand{\intltwenty}{Int\$\textsubscript{2011}}     % Int$_2011
\newcommand{\eek}{kroon}                                % Estonian kroon (table headers)

% ---- Numbers quoted in the prose: generated by scripts/R/tables.R ----
% Auto-generated by scripts/R/tables.R -- do not edit by hand
\newcommand{\ee}[1]{\csname eeGDP#1\endcsname}
\newcommand{\eeb}[1]{\csname eeGDPb#1\endcsname}
\newcommand{\pop}[1]{\csname eePOP#1\endcsname}
\newcommand{\gr}[2]{\csname eeGR#1x#2\endcsname}
\newcommand{\grx}[3]{\csname GR#1x#2x#3\endcsname}
\newcommand{\pctfin}[1]{\csname eePCTfin#1\endcsname}
\newcommand{\pctswe}[1]{\csname eePCTswe#1\endcsname}
\newcommand{\pctwc}[1]{\csname eePCTwc#1\endcsname}
\newcommand{\pctussr}[1]{\csname eePCTussr#1\endcsname}
\newcommand{\pctlt}[1]{\csname eePCTlt#1\endcsname}
\newcommand{\pctlv}[1]{\csname eePCTlv#1\endcsname}
\newcommand{\cum}[2]{\csname eeCUM#1x#2\endcsname}
\expandafter\def\csname eeGDPbench1990\endcsname{8,425}
\expandafter\def\csname eeGDP1913\endcsname{2,058}
\expandafter\def\csname eeGDPb1913\endcsname{3,280}
\expandafter\def\csname eePOP1913\endcsname{1,164.0}
\expandafter\def\csname eePCTfin1913\endcsname{97}
\expandafter\def\csname eePCTswe1913\endcsname{72}
\expandafter\def\csname eePCTwc1913\endcsname{51}
\expandafter\def\csname eePCTussr1913\endcsname{146}
\expandafter\def\csname eePCTlt1913\endcsname{126}
\expandafter\def\csname eePCTlv1913\endcsname{87}
\expandafter\def\csname eeGDP1919\endcsname{1,600}
\expandafter\def\csname eeGDPb1919\endcsname{2,550}
\expandafter\def\csname eePOP1919\endcsname{1,064.0}
\expandafter\def\csname eePCTfin1919\endcsname{96}
\expandafter\def\csname eePCTswe1919\endcsname{57}
\expandafter\def\csname eePCTwc1919\endcsname{48}
\expandafter\def\csname eePCTussr1919\endcsname{279}
\expandafter\def\csname eePCTlt1919\endcsname{157}
\expandafter\def\csname eePCTlv1919\endcsname{108}
\expandafter\def\csname eeGDP1920\endcsname{1,475}
\expandafter\def\csname eeGDPb1920\endcsname{2,351}
\expandafter\def\csname eePOP1920\endcsname{1,067.8}
\expandafter\def\csname eePCTfin1920\endcsname{80}
\expandafter\def\csname eePCTswe1920\endcsname{49}
\expandafter\def\csname eePCTwc1920\endcsname{43}
\expandafter\def\csname eePCTussr1920\endcsname{256}
\expandafter\def\csname eePCTlt1920\endcsname{124}
\expandafter\def\csname eePCTlv1920\endcsname{97}
\expandafter\def\csname eeGDP1921\endcsname{1,673}
\expandafter\def\csname eeGDPb1921\endcsname{2,666}
\expandafter\def\csname eePOP1921\endcsname{1,087.1}
\expandafter\def\csname eePCTfin1921\endcsname{89}
\expandafter\def\csname eePCTswe1921\endcsname{59}
\expandafter\def\csname eePCTwc1921\endcsname{48}
\expandafter\def\csname eePCTussr1921\endcsname{318}
\expandafter\def\csname eePCTlt1921\endcsname{118}
\expandafter\def\csname eePCTlv1921\endcsname{100}
\expandafter\def\csname eeGDP1922\endcsname{1,756}
\expandafter\def\csname eeGDPb1922\endcsname{2,800}
\expandafter\def\csname eePOP1922\endcsname{1,102.4}
\expandafter\def\csname eePCTfin1922\endcsname{85}
\expandafter\def\csname eePCTswe1922\endcsname{58}
\expandafter\def\csname eePCTwc1922\endcsname{46}
\expandafter\def\csname eePCTussr1922\endcsname{287}
\expandafter\def\csname eePCTlt1922\endcsname{115}
\expandafter\def\csname eePCTlv1922\endcsname{106}
\expandafter\def\csname eeGDP1923\endcsname{1,913}
\expandafter\def\csname eeGDPb1923\endcsname{3,050}
\expandafter\def\csname eePOP1923\endcsname{1,110.8}
\expandafter\def\csname eePCTfin1923\endcsname{87}
\expandafter\def\csname eePCTswe1923\endcsname{61}
\expandafter\def\csname eePCTwc1923\endcsname{52}
\expandafter\def\csname eePCTussr1923\endcsname{272}
\expandafter\def\csname eePCTlt1923\endcsname{109}
\expandafter\def\csname eePCTlv1923\endcsname{101}
\expandafter\def\csname eeGDP1924\endcsname{1,895}
\expandafter\def\csname eeGDPb1924\endcsname{3,020}
\expandafter\def\csname eePOP1924\endcsname{1,115.6}
\expandafter\def\csname eePCTfin1924\endcsname{85}
\expandafter\def\csname eePCTswe1924\endcsname{56}
\expandafter\def\csname eePCTwc1924\endcsname{47}
\expandafter\def\csname eePCTussr1924\endcsname{213}
\expandafter\def\csname eePCTlt1924\endcsname{115}
\expandafter\def\csname eePCTlv1924\endcsname{98}
\expandafter\def\csname eeGDP1925\endcsname{2,038}
\expandafter\def\csname eeGDPb1925\endcsname{3,249}
\expandafter\def\csname eePOP1925\endcsname{1,117.0}
\expandafter\def\csname eePCTfin1925\endcsname{88}
\expandafter\def\csname eePCTswe1925\endcsname{59}
\expandafter\def\csname eePCTwc1925\endcsname{49}
\expandafter\def\csname eePCTussr1925\endcsname{183}
\expandafter\def\csname eePCTlt1925\endcsname{125}
\expandafter\def\csname eePCTlv1925\endcsname{98}
\expandafter\def\csname eeGDP1926\endcsname{2,030}
\expandafter\def\csname eeGDPb1926\endcsname{3,235}
\expandafter\def\csname eePOP1926\endcsname{1,116.8}
\expandafter\def\csname eePCTfin1926\endcsname{85}
\expandafter\def\csname eePCTswe1926\endcsname{56}
\expandafter\def\csname eePCTwc1926\endcsname{48}
\expandafter\def\csname eePCTussr1926\endcsname{162}
\expandafter\def\csname eePCTlt1926\endcsname{136}
\expandafter\def\csname eePCTlv1926\endcsname{97}
\expandafter\def\csname eeGDP1927\endcsname{2,043}
\expandafter\def\csname eeGDPb1927\endcsname{3,256}
\expandafter\def\csname eePOP1927\endcsname{1,115.6}
\expandafter\def\csname eePCTfin1927\endcsname{80}
\expandafter\def\csname eePCTswe1927\endcsname{55}
\expandafter\def\csname eePCTwc1927\endcsname{46}
\expandafter\def\csname eePCTussr1927\endcsname{157}
\expandafter\def\csname eePCTlt1927\endcsname{136}
\expandafter\def\csname eePCTlv1927\endcsname{95}
\expandafter\def\csname eeGDP1928\endcsname{2,046}
\expandafter\def\csname eeGDPb1928\endcsname{3,262}
\expandafter\def\csname eePOP1928\endcsname{1,115.7}
\expandafter\def\csname eePCTfin1928\endcsname{76}
\expandafter\def\csname eePCTswe1928\endcsname{53}
\expandafter\def\csname eePCTwc1928\endcsname{45}
\expandafter\def\csname eePCTussr1928\endcsname{149}
\expandafter\def\csname eePCTlt1928\endcsname{137}
\expandafter\def\csname eePCTlv1928\endcsname{95}
\expandafter\def\csname eeGDP1929\endcsname{2,113}
\expandafter\def\csname eeGDPb1929\endcsname{3,368}
\expandafter\def\csname eePOP1929\endcsname{1,115.7}
\expandafter\def\csname eePCTfin1929\endcsname{78}
\expandafter\def\csname eePCTswe1929\endcsname{52}
\expandafter\def\csname eePCTwc1929\endcsname{45}
\expandafter\def\csname eePCTussr1929\endcsname{152}
\expandafter\def\csname eePCTlt1929\endcsname{134}
\expandafter\def\csname eePCTlv1929\endcsname{90}
\expandafter\def\csname eeGDP1930\endcsname{2,282}
\expandafter\def\csname eeGDPb1930\endcsname{3,638}
\expandafter\def\csname eePOP1930\endcsname{1,116.1}
\expandafter\def\csname eePCTfin1930\endcsname{86}
\expandafter\def\csname eePCTswe1930\endcsname{54}
\expandafter\def\csname eePCTwc1930\endcsname{49}
\expandafter\def\csname eePCTussr1930\endcsname{158}
\expandafter\def\csname eePCTlt1930\endcsname{137}
\expandafter\def\csname eePCTlv1930\endcsname{90}
\expandafter\def\csname eeGDP1931\endcsname{2,176}
\expandafter\def\csname eeGDPb1931\endcsname{3,469}
\expandafter\def\csname eePOP1931\endcsname{1,118.4}
\expandafter\def\csname eePCTfin1931\endcsname{84}
\expandafter\def\csname eePCTswe1931\endcsname{53}
\expandafter\def\csname eePCTwc1931\endcsname{50}
\expandafter\def\csname eePCTussr1931\endcsname{149}
\expandafter\def\csname eePCTlt1931\endcsname{130}
\expandafter\def\csname eePCTlv1931\endcsname{94}
\expandafter\def\csname eeGDP1932\endcsname{2,105}
\expandafter\def\csname eeGDPb1932\endcsname{3,356}
\expandafter\def\csname eePOP1932\endcsname{1,121.5}
\expandafter\def\csname eePCTfin1932\endcsname{83}
\expandafter\def\csname eePCTswe1932\endcsname{53}
\expandafter\def\csname eePCTwc1932\endcsname{51}
\expandafter\def\csname eePCTussr1932\endcsname{146}
\expandafter\def\csname eePCTlt1932\endcsname{124}
\expandafter\def\csname eePCTlv1932\endcsname{92}
\expandafter\def\csname eeGDP1933\endcsname{2,186}
\expandafter\def\csname eeGDPb1933\endcsname{3,485}
\expandafter\def\csname eePOP1933\endcsname{1,124.3}
\expandafter\def\csname eePCTfin1933\endcsname{81}
\expandafter\def\csname eePCTswe1933\endcsname{54}
\expandafter\def\csname eePCTwc1933\endcsname{51}
\expandafter\def\csname eePCTussr1933\endcsname{146}
\expandafter\def\csname eePCTlt1933\endcsname{127}
\expandafter\def\csname eePCTlv1933\endcsname{88}
\expandafter\def\csname eeGDP1934\endcsname{2,353}
\expandafter\def\csname eeGDPb1934\endcsname{3,751}
\expandafter\def\csname eePOP1934\endcsname{1,126.4}
\expandafter\def\csname eePCTfin1934\endcsname{79}
\expandafter\def\csname eePCTswe1934\endcsname{55}
\expandafter\def\csname eePCTwc1934\endcsname{53}
\expandafter\def\csname eePCTussr1934\endcsname{144}
\expandafter\def\csname eePCTlt1934\endcsname{131}
\expandafter\def\csname eePCTlv1934\endcsname{86}
\expandafter\def\csname eeGDP1935\endcsname{2,355}
\expandafter\def\csname eeGDPb1935\endcsname{3,754}
\expandafter\def\csname eePOP1935\endcsname{1,128.9}
\expandafter\def\csname eePCTfin1935\endcsname{76}
\expandafter\def\csname eePCTswe1935\endcsname{52}
\expandafter\def\csname eePCTwc1935\endcsname{51}
\expandafter\def\csname eePCTussr1935\endcsname{126}
\expandafter\def\csname eePCTlt1935\endcsname{124}
\expandafter\def\csname eePCTlv1935\endcsname{85}
\expandafter\def\csname eeGDP1936\endcsname{2,492}
\expandafter\def\csname eeGDPb1936\endcsname{3,972}
\expandafter\def\csname eePOP1936\endcsname{1,130.0}
\expandafter\def\csname eePCTfin1936\endcsname{76}
\expandafter\def\csname eePCTswe1936\endcsname{54}
\expandafter\def\csname eePCTwc1936\endcsname{52}
\expandafter\def\csname eePCTussr1936\endcsname{125}
\expandafter\def\csname eePCTlt1936\endcsname{129}
\expandafter\def\csname eePCTlv1936\endcsname{89}
\expandafter\def\csname eeGDP1937\endcsname{2,722}
\expandafter\def\csname eeGDPb1937\endcsname{4,338}
\expandafter\def\csname eePOP1937\endcsname{1,130.7}
\expandafter\def\csname eePCTfin1937\endcsname{79}
\expandafter\def\csname eePCTswe1937\endcsname{57}
\expandafter\def\csname eePCTwc1937\endcsname{54}
\expandafter\def\csname eePCTussr1937\endcsname{126}
\expandafter\def\csname eePCTlt1937\endcsname{135}
\expandafter\def\csname eePCTlv1937\endcsname{90}
\expandafter\def\csname eeGDP1938\endcsname{2,685}
\expandafter\def\csname eeGDPb1938\endcsname{4,280}
\expandafter\def\csname eePOP1938\endcsname{1,132.5}
\expandafter\def\csname eePCTfin1938\endcsname{75}
\expandafter\def\csname eePCTswe1938\endcsname{55}
\expandafter\def\csname eePCTwc1938\endcsname{52}
\expandafter\def\csname eePCTussr1938\endcsname{125}
\expandafter\def\csname eePCTlt1938\endcsname{130}
\expandafter\def\csname eePCTlv1938\endcsname{87}
\expandafter\def\csname eeGDP1939\endcsname{2,736}
\expandafter\def\csname eeGDPb1939\endcsname{4,361}
\expandafter\def\csname eePOP1939\endcsname{1,140.4}
\expandafter\def\csname eePCTfin1939\endcsname{80}
\expandafter\def\csname eePCTswe1939\endcsname{52}
\expandafter\def\csname eePCTwc1939\endcsname{50}
\expandafter\def\csname eePCTussr1939\endcsname{122}
\expandafter\def\csname eePCTlt1939\endcsname{141}
\expandafter\def\csname eePCTlv1939\endcsname{---}
\expandafter\def\csname eeGDP1940\endcsname{2,771}
\expandafter\def\csname eeGDPb1940\endcsname{4,418}
\expandafter\def\csname eePOP1940\endcsname{1,122.7}
\expandafter\def\csname eePCTfin1940\endcsname{86}
\expandafter\def\csname eePCTswe1940\endcsname{57}
\expandafter\def\csname eePCTwc1940\endcsname{52}
\expandafter\def\csname eePCTussr1940\endcsname{129}
\expandafter\def\csname eePCTlt1940\endcsname{159}
\expandafter\def\csname eePCTlv1940\endcsname{---}
\expandafter\def\csname eeGDP1941\endcsname{2,390}
\expandafter\def\csname eeGDPb1941\endcsname{3,809}
\expandafter\def\csname eePOP1941\endcsname{1,049.1}
\expandafter\def\csname eePCTfin1941\endcsname{72}
\expandafter\def\csname eePCTswe1941\endcsname{51}
\expandafter\def\csname eePCTwc1941\endcsname{44}
\expandafter\def\csname eePCTussr1941\endcsname{113}
\expandafter\def\csname eePCTlt1941\endcsname{147}
\expandafter\def\csname eePCTlv1941\endcsname{---}
\expandafter\def\csname eeGDP1942\endcsname{2,382}
\expandafter\def\csname eeGDPb1942\endcsname{3,797}
\expandafter\def\csname eePOP1942\endcsname{1,007.0}
\expandafter\def\csname eePCTfin1942\endcsname{72}
\expandafter\def\csname eePCTswe1942\endcsname{50}
\expandafter\def\csname eePCTwc1942\endcsname{44}
\expandafter\def\csname eePCTussr1942\endcsname{115}
\expandafter\def\csname eePCTlt1942\endcsname{161}
\expandafter\def\csname eePCTlv1942\endcsname{---}
\expandafter\def\csname eeGDP1943\endcsname{2,328}
\expandafter\def\csname eeGDPb1943\endcsname{3,710}
\expandafter\def\csname eePOP1943\endcsname{981.9}
\expandafter\def\csname eePCTfin1943\endcsname{63}
\expandafter\def\csname eePCTswe1943\endcsname{48}
\expandafter\def\csname eePCTwc1943\endcsname{43}
\expandafter\def\csname eePCTussr1943\endcsname{114}
\expandafter\def\csname eePCTlt1943\endcsname{181}
\expandafter\def\csname eePCTlv1943\endcsname{---}
\expandafter\def\csname eeGDP1944\endcsname{1,909}
\expandafter\def\csname eeGDPb1944\endcsname{3,044}
\expandafter\def\csname eePOP1944\endcsname{917.8}
\expandafter\def\csname eePCTfin1944\endcsname{52}
\expandafter\def\csname eePCTswe1944\endcsname{38}
\expandafter\def\csname eePCTwc1944\endcsname{36}
\expandafter\def\csname eePCTussr1944\endcsname{95}
\expandafter\def\csname eePCTlt1944\endcsname{179}
\expandafter\def\csname eePCTlv1944\endcsname{---}
\expandafter\def\csname eeGDP1945\endcsname{2,656}
\expandafter\def\csname eeGDPb1945\endcsname{4,234}
\expandafter\def\csname eePOP1945\endcsname{903.1}
\expandafter\def\csname eePCTfin1945\endcsname{77}
\expandafter\def\csname eePCTswe1945\endcsname{52}
\expandafter\def\csname eePCTwc1945\endcsname{58}
\expandafter\def\csname eePCTussr1945\endcsname{134}
\expandafter\def\csname eePCTlt1945\endcsname{201}
\expandafter\def\csname eePCTlv1945\endcsname{---}
\expandafter\def\csname eeGDP1946\endcsname{2,833}
\expandafter\def\csname eeGDPb1946\endcsname{4,516}
\expandafter\def\csname eePOP1946\endcsname{938.7}
\expandafter\def\csname eePCTfin1946\endcsname{77}
\expandafter\def\csname eePCTswe1946\endcsname{50}
\expandafter\def\csname eePCTwc1946\endcsname{67}
\expandafter\def\csname eePCTussr1946\endcsname{148}
\expandafter\def\csname eePCTlt1946\endcsname{178}
\expandafter\def\csname eePCTlv1946\endcsname{---}
\expandafter\def\csname eeGDP1947\endcsname{3,123}
\expandafter\def\csname eeGDPb1947\endcsname{4,978}
\expandafter\def\csname eePOP1947\endcsname{990.2}
\expandafter\def\csname eePCTfin1947\endcsname{84}
\expandafter\def\csname eePCTswe1947\endcsname{51}
\expandafter\def\csname eePCTwc1947\endcsname{71}
\expandafter\def\csname eePCTussr1947\endcsname{147}
\expandafter\def\csname eePCTlt1947\endcsname{161}
\expandafter\def\csname eePCTlv1947\endcsname{---}
\expandafter\def\csname eeGDP1948\endcsname{3,323}
\expandafter\def\csname eeGDPb1948\endcsname{5,297}
\expandafter\def\csname eePOP1948\endcsname{1,026.3}
\expandafter\def\csname eePCTfin1948\endcsname{84}
\expandafter\def\csname eePCTswe1948\endcsname{54}
\expandafter\def\csname eePCTwc1948\endcsname{72}
\expandafter\def\csname eePCTussr1948\endcsname{138}
\expandafter\def\csname eePCTlt1948\endcsname{160}
\expandafter\def\csname eePCTlv1948\endcsname{---}
\expandafter\def\csname eeGDP1949\endcsname{3,474}
\expandafter\def\csname eeGDPb1949\endcsname{5,538}
\expandafter\def\csname eePOP1949\endcsname{1,026.1}
\expandafter\def\csname eePCTfin1949\endcsname{84}
\expandafter\def\csname eePCTswe1949\endcsname{55}
\expandafter\def\csname eePCTwc1949\endcsname{69}
\expandafter\def\csname eePCTussr1949\endcsname{132}
\expandafter\def\csname eePCTlt1949\endcsname{158}
\expandafter\def\csname eePCTlv1949\endcsname{---}
\expandafter\def\csname eeGDP1950\endcsname{3,667}
\expandafter\def\csname eeGDPb1950\endcsname{5,846}
\expandafter\def\csname eePOP1950\endcsname{1,036.4}
\expandafter\def\csname eePCTfin1950\endcsname{86}
\expandafter\def\csname eePCTswe1950\endcsname{54}
\expandafter\def\csname eePCTwc1950\endcsname{66}
\expandafter\def\csname eePCTussr1950\endcsname{129}
\expandafter\def\csname eePCTlt1950\endcsname{160}
\expandafter\def\csname eePCTlv1950\endcsname{---}
\expandafter\def\csname eeGDP1951\endcsname{3,709}
\expandafter\def\csname eeGDPb1951\endcsname{5,912}
\expandafter\def\csname eePOP1951\endcsname{1,061.6}
\expandafter\def\csname eePCTfin1951\endcsname{81}
\expandafter\def\csname eePCTswe1951\endcsname{53}
\expandafter\def\csname eePCTwc1951\endcsname{64}
\expandafter\def\csname eePCTussr1951\endcsname{132}
\expandafter\def\csname eePCTlt1951\endcsname{157}
\expandafter\def\csname eePCTlv1951\endcsname{---}
\expandafter\def\csname eeGDP1952\endcsname{3,758}
\expandafter\def\csname eeGDPb1952\endcsname{5,990}
\expandafter\def\csname eePOP1952\endcsname{1,083.1}
\expandafter\def\csname eePCTfin1952\endcsname{80}
\expandafter\def\csname eePCTswe1952\endcsname{54}
\expandafter\def\csname eePCTwc1952\endcsname{63}
\expandafter\def\csname eePCTussr1952\endcsname{128}
\expandafter\def\csname eePCTlt1952\endcsname{154}
\expandafter\def\csname eePCTlv1952\endcsname{---}
\expandafter\def\csname eeGDP1953\endcsname{3,747}
\expandafter\def\csname eeGDPb1953\endcsname{5,973}
\expandafter\def\csname eePOP1953\endcsname{1,106.5}
\expandafter\def\csname eePCTfin1953\endcsname{81}
\expandafter\def\csname eePCTswe1953\endcsname{52}
\expandafter\def\csname eePCTwc1953\endcsname{60}
\expandafter\def\csname eePCTussr1953\endcsname{124}
\expandafter\def\csname eePCTlt1953\endcsname{141}
\expandafter\def\csname eePCTlv1953\endcsname{---}
\expandafter\def\csname eeGDP1954\endcsname{3,787}
\expandafter\def\csname eeGDPb1954\endcsname{6,036}
\expandafter\def\csname eePOP1954\endcsname{1,128.9}
\expandafter\def\csname eePCTfin1954\endcsname{76}
\expandafter\def\csname eePCTswe1954\endcsname{51}
\expandafter\def\csname eePCTwc1954\endcsname{58}
\expandafter\def\csname eePCTussr1954\endcsname{122}
\expandafter\def\csname eePCTlt1954\endcsname{142}
\expandafter\def\csname eePCTlv1954\endcsname{---}
\expandafter\def\csname eeGDP1955\endcsname{3,951}
\expandafter\def\csname eeGDPb1955\endcsname{6,298}
\expandafter\def\csname eePOP1955\endcsname{1,144.2}
\expandafter\def\csname eePCTfin1955\endcsname{76}
\expandafter\def\csname eePCTswe1955\endcsname{52}
\expandafter\def\csname eePCTwc1955\endcsname{58}
\expandafter\def\csname eePCTussr1955\endcsname{119}
\expandafter\def\csname eePCTlt1955\endcsname{151}
\expandafter\def\csname eePCTlv1955\endcsname{---}
\expandafter\def\csname eeGDP1956\endcsname{4,243}
\expandafter\def\csname eeGDPb1956\endcsname{6,763}
\expandafter\def\csname eePOP1956\endcsname{1,157.9}
\expandafter\def\csname eePCTfin1956\endcsname{80}
\expandafter\def\csname eePCTswe1956\endcsname{54}
\expandafter\def\csname eePCTwc1956\endcsname{60}
\expandafter\def\csname eePCTussr1956\endcsname{119}
\expandafter\def\csname eePCTlt1956\endcsname{150}
\expandafter\def\csname eePCTlv1956\endcsname{---}
\expandafter\def\csname eeGDP1957\endcsname{4,460}
\expandafter\def\csname eeGDPb1957\endcsname{7,109}
\expandafter\def\csname eePOP1957\endcsname{1,171.9}
\expandafter\def\csname eePCTfin1957\endcsname{81}
\expandafter\def\csname eePCTswe1957\endcsname{55}
\expandafter\def\csname eePCTwc1957\endcsname{61}
\expandafter\def\csname eePCTussr1957\endcsname{125}
\expandafter\def\csname eePCTlt1957\endcsname{142}
\expandafter\def\csname eePCTlv1957\endcsname{---}
\expandafter\def\csname eeGDP1958\endcsname{4,727}
\expandafter\def\csname eeGDPb1958\endcsname{7,535}
\expandafter\def\csname eePOP1958\endcsname{1,187.8}
\expandafter\def\csname eePCTfin1958\endcsname{86}
\expandafter\def\csname eePCTswe1958\endcsname{58}
\expandafter\def\csname eePCTwc1958\endcsname{64}
\expandafter\def\csname eePCTussr1958\endcsname{125}
\expandafter\def\csname eePCTlt1958\endcsname{142}
\expandafter\def\csname eePCTlv1958\endcsname{---}
\expandafter\def\csname eeGDP1959\endcsname{5,010}
\expandafter\def\csname eeGDPb1959\endcsname{7,986}
\expandafter\def\csname eePOP1959\endcsname{1,201.6}
\expandafter\def\csname eePCTfin1959\endcsname{87}
\expandafter\def\csname eePCTswe1959\endcsname{60}
\expandafter\def\csname eePCTwc1959\endcsname{65}
\expandafter\def\csname eePCTussr1959\endcsname{137}
\expandafter\def\csname eePCTlt1959\endcsname{140}
\expandafter\def\csname eePCTlv1959\endcsname{---}
\expandafter\def\csname eeGDP1960\endcsname{5,319}
\expandafter\def\csname eeGDPb1960\endcsname{8,478}
\expandafter\def\csname eePOP1960\endcsname{1,211.5}
\expandafter\def\csname eePCTfin1960\endcsname{85}
\expandafter\def\csname eePCTswe1960\endcsname{61}
\expandafter\def\csname eePCTwc1960\endcsname{65}
\expandafter\def\csname eePCTussr1960\endcsname{135}
\expandafter\def\csname eePCTlt1960\endcsname{144}
\expandafter\def\csname eePCTlv1960\endcsname{---}
\expandafter\def\csname eeGDP1961\endcsname{5,513}
\expandafter\def\csname eeGDPb1961\endcsname{8,788}
\expandafter\def\csname eePOP1961\endcsname{1,225.1}
\expandafter\def\csname eePCTfin1961\endcsname{83}
\expandafter\def\csname eePCTswe1961\endcsname{60}
\expandafter\def\csname eePCTwc1961\endcsname{65}
\expandafter\def\csname eePCTussr1961\endcsname{135}
\expandafter\def\csname eePCTlt1961\endcsname{147}
\expandafter\def\csname eePCTlv1961\endcsname{---}
\expandafter\def\csname eeGDP1962\endcsname{5,547}
\expandafter\def\csname eeGDPb1962\endcsname{8,841}
\expandafter\def\csname eePOP1962\endcsname{1,241.6}
\expandafter\def\csname eePCTfin1962\endcsname{81}
\expandafter\def\csname eePCTswe1962\endcsname{59}
\expandafter\def\csname eePCTwc1962\endcsname{64}
\expandafter\def\csname eePCTussr1962\endcsname{134}
\expandafter\def\csname eePCTlt1962\endcsname{148}
\expandafter\def\csname eePCTlv1962\endcsname{---}
\expandafter\def\csname eeGDP1963\endcsname{5,840}
\expandafter\def\csname eeGDPb1963\endcsname{9,309}
\expandafter\def\csname eePOP1963\endcsname{1,258.9}
\expandafter\def\csname eePCTfin1963\endcsname{84}
\expandafter\def\csname eePCTswe1963\endcsname{59}
\expandafter\def\csname eePCTwc1963\endcsname{65}
\expandafter\def\csname eePCTussr1963\endcsname{147}
\expandafter\def\csname eePCTlt1963\endcsname{147}
\expandafter\def\csname eePCTlv1963\endcsname{---}
\expandafter\def\csname eeGDP1964\endcsname{6,150}
\expandafter\def\csname eeGDPb1964\endcsname{9,803}
\expandafter\def\csname eePOP1964\endcsname{1,277.1}
\expandafter\def\csname eePCTfin1964\endcsname{84}
\expandafter\def\csname eePCTswe1964\endcsname{58}
\expandafter\def\csname eePCTwc1964\endcsname{65}
\expandafter\def\csname eePCTussr1964\endcsname{139}
\expandafter\def\csname eePCTlt1964\endcsname{145}
\expandafter\def\csname eePCTlv1964\endcsname{---}
\expandafter\def\csname eeGDP1965\endcsname{6,466}
\expandafter\def\csname eeGDPb1965\endcsname{10,308}
\expandafter\def\csname eePOP1965\endcsname{1,294.6}
\expandafter\def\csname eePCTfin1965\endcsname{84}
\expandafter\def\csname eePCTswe1965\endcsname{60}
\expandafter\def\csname eePCTwc1965\endcsname{67}
\expandafter\def\csname eePCTussr1965\endcsname{140}
\expandafter\def\csname eePCTlt1965\endcsname{144}
\expandafter\def\csname eePCTlv1965\endcsname{---}
\expandafter\def\csname eeGDP1966\endcsname{6,613}
\expandafter\def\csname eeGDPb1966\endcsname{10,541}
\expandafter\def\csname eePOP1966\endcsname{1,308.6}
\expandafter\def\csname eePCTfin1966\endcsname{85}
\expandafter\def\csname eePCTswe1966\endcsname{60}
\expandafter\def\csname eePCTwc1966\endcsname{67}
\expandafter\def\csname eePCTussr1966\endcsname{138}
\expandafter\def\csname eePCTlt1966\endcsname{137}
\expandafter\def\csname eePCTlv1966\endcsname{---}
\expandafter\def\csname eeGDP1967\endcsname{7,013}
\expandafter\def\csname eeGDPb1967\endcsname{11,178}
\expandafter\def\csname eePOP1967\endcsname{1,318.9}
\expandafter\def\csname eePCTfin1967\endcsname{88}
\expandafter\def\csname eePCTswe1967\endcsname{63}
\expandafter\def\csname eePCTwc1967\endcsname{69}
\expandafter\def\csname eePCTussr1967\endcsname{141}
\expandafter\def\csname eePCTlt1967\endcsname{135}
\expandafter\def\csname eePCTlv1967\endcsname{---}
\expandafter\def\csname eeGDP1968\endcsname{7,348}
\expandafter\def\csname eeGDPb1968\endcsname{11,712}
\expandafter\def\csname eePOP1968\endcsname{1,331.2}
\expandafter\def\csname eePCTfin1968\endcsname{91}
\expandafter\def\csname eePCTswe1968\endcsname{64}
\expandafter\def\csname eePCTwc1968\endcsname{70}
\expandafter\def\csname eePCTussr1968\endcsname{141}
\expandafter\def\csname eePCTlt1968\endcsname{136}
\expandafter\def\csname eePCTlv1968\endcsname{---}
\expandafter\def\csname eeGDP1969\endcsname{7,442}
\expandafter\def\csname eeGDPb1969\endcsname{11,863}
\expandafter\def\csname eePOP1969\endcsname{1,345.2}
\expandafter\def\csname eePCTfin1969\endcsname{84}
\expandafter\def\csname eePCTswe1969\endcsname{62}
\expandafter\def\csname eePCTwc1969\endcsname{67}
\expandafter\def\csname eePCTussr1969\endcsname{142}
\expandafter\def\csname eePCTlt1969\endcsname{132}
\expandafter\def\csname eePCTlv1969\endcsname{---}
\expandafter\def\csname eeGDP1970\endcsname{7,688}
\expandafter\def\csname eeGDPb1970\endcsname{12,255}
\expandafter\def\csname eePOP1970\endcsname{1,360.1}
\expandafter\def\csname eePCTfin1970\endcsname{80}
\expandafter\def\csname eePCTswe1970\endcsname{60}
\expandafter\def\csname eePCTwc1970\endcsname{67}
\expandafter\def\csname eePCTussr1970\endcsname{138}
\expandafter\def\csname eePCTlt1970\endcsname{131}
\expandafter\def\csname eePCTlv1970\endcsname{---}
\expandafter\def\csname eeGDP1971\endcsname{8,002}
\expandafter\def\csname eeGDPb1971\endcsname{12,755}
\expandafter\def\csname eePOP1971\endcsname{1,377.0}
\expandafter\def\csname eePCTfin1971\endcsname{82}
\expandafter\def\csname eePCTswe1971\endcsname{63}
\expandafter\def\csname eePCTwc1971\endcsname{68}
\expandafter\def\csname eePCTussr1971\endcsname{141}
\expandafter\def\csname eePCTlt1971\endcsname{132}
\expandafter\def\csname eePCTlv1971\endcsname{---}
\expandafter\def\csname eeGDP1972\endcsname{7,896}
\expandafter\def\csname eeGDPb1972\endcsname{12,587}
\expandafter\def\csname eePOP1972\endcsname{1,392.5}
\expandafter\def\csname eePCTfin1972\endcsname{76}
\expandafter\def\csname eePCTswe1972\endcsname{61}
\expandafter\def\csname eePCTwc1972\endcsname{65}
\expandafter\def\csname eePCTussr1972\endcsname{140}
\expandafter\def\csname eePCTlt1972\endcsname{126}
\expandafter\def\csname eePCTlv1972\endcsname{---}
\expandafter\def\csname eeGDP1973\endcsname{8,080}
\expandafter\def\csname eeGDPb1973\endcsname{12,879}
\expandafter\def\csname eePOP1973\endcsname{1,406.0}
\expandafter\def\csname eePCTfin1973\endcsname{73}
\expandafter\def\csname eePCTswe1973\endcsname{60}
\expandafter\def\csname eePCTwc1973\endcsname{63}
\expandafter\def\csname eePCTussr1973\endcsname{133}
\expandafter\def\csname eePCTlt1973\endcsname{125}
\expandafter\def\csname eePCTlv1973\endcsname{---}
\expandafter\def\csname eeGDP1974\endcsname{8,217}
\expandafter\def\csname eeGDPb1974\endcsname{13,098}
\expandafter\def\csname eePOP1974\endcsname{1,418.2}
\expandafter\def\csname eePCTfin1974\endcsname{72}
\expandafter\def\csname eePCTswe1974\endcsname{59}
\expandafter\def\csname eePCTwc1974\endcsname{64}
\expandafter\def\csname eePCTussr1974\endcsname{133}
\expandafter\def\csname eePCTlt1974\endcsname{125}
\expandafter\def\csname eePCTlv1974\endcsname{---}
\expandafter\def\csname eeGDP1975\endcsname{8,419}
\expandafter\def\csname eeGDPb1975\endcsname{13,419}
\expandafter\def\csname eePOP1975\endcsname{1,429.4}
\expandafter\def\csname eePCTfin1975\endcsname{74}
\expandafter\def\csname eePCTswe1975\endcsname{59}
\expandafter\def\csname eePCTwc1975\endcsname{66}
\expandafter\def\csname eePCTussr1975\endcsname{137}
\expandafter\def\csname eePCTlt1975\endcsname{124}
\expandafter\def\csname eePCTlv1975\endcsname{---}
\expandafter\def\csname eeGDP1976\endcsname{8,623}
\expandafter\def\csname eeGDPb1976\endcsname{13,745}
\expandafter\def\csname eePOP1976\endcsname{1,439.6}
\expandafter\def\csname eePCTfin1976\endcsname{76}
\expandafter\def\csname eePCTswe1976\endcsname{60}
\expandafter\def\csname eePCTwc1976\endcsname{65}
\expandafter\def\csname eePCTussr1976\endcsname{136}
\expandafter\def\csname eePCTlt1976\endcsname{124}
\expandafter\def\csname eePCTlv1976\endcsname{---}
\expandafter\def\csname eeGDP1977\endcsname{8,601}
\expandafter\def\csname eeGDPb1977\endcsname{13,709}
\expandafter\def\csname eePOP1977\endcsname{1,450.2}
\expandafter\def\csname eePCTfin1977\endcsname{76}
\expandafter\def\csname eePCTswe1977\endcsname{61}
\expandafter\def\csname eePCTwc1977\endcsname{63}
\expandafter\def\csname eePCTussr1977\endcsname{133}
\expandafter\def\csname eePCTlt1977\endcsname{122}
\expandafter\def\csname eePCTlv1977\endcsname{---}
\expandafter\def\csname eeGDP1978\endcsname{8,474}
\expandafter\def\csname eeGDPb1978\endcsname{13,508}
\expandafter\def\csname eePOP1978\endcsname{1,460.2}
\expandafter\def\csname eePCTfin1978\endcsname{73}
\expandafter\def\csname eePCTswe1978\endcsname{60}
\expandafter\def\csname eePCTwc1978\endcsname{61}
\expandafter\def\csname eePCTussr1978\endcsname{129}
\expandafter\def\csname eePCTlt1978\endcsname{118}
\expandafter\def\csname eePCTlv1978\endcsname{---}
\expandafter\def\csname eeGDP1979\endcsname{8,523}
\expandafter\def\csname eeGDPb1979\endcsname{13,585}
\expandafter\def\csname eePOP1979\endcsname{1,468.3}
\expandafter\def\csname eePCTfin1979\endcsname{69}
\expandafter\def\csname eePCTswe1979\endcsname{58}
\expandafter\def\csname eePCTwc1979\endcsname{59}
\expandafter\def\csname eePCTussr1979\endcsname{132}
\expandafter\def\csname eePCTlt1979\endcsname{121}
\expandafter\def\csname eePCTlv1979\endcsname{---}
\expandafter\def\csname eeGDP1980\endcsname{8,785}
\expandafter\def\csname eeGDPb1980\endcsname{14,003}
\expandafter\def\csname eePOP1980\endcsname{1,477.2}
\expandafter\def\csname eePCTfin1980\endcsname{68}
\expandafter\def\csname eePCTswe1980\endcsname{59}
\expandafter\def\csname eePCTwc1980\endcsname{61}
\expandafter\def\csname eePCTussr1980\endcsname{137}
\expandafter\def\csname eePCTlt1980\endcsname{127}
\expandafter\def\csname eePCTlv1980\endcsname{---}
\expandafter\def\csname eeGDP1981\endcsname{8,585}
\expandafter\def\csname eeGDPb1981\endcsname{13,684}
\expandafter\def\csname eePOP1981\endcsname{1,487.7}
\expandafter\def\csname eePCTfin1981\endcsname{65}
\expandafter\def\csname eePCTswe1981\endcsname{58}
\expandafter\def\csname eePCTwc1981\endcsname{59}
\expandafter\def\csname eePCTussr1981\endcsname{133}
\expandafter\def\csname eePCTlt1981\endcsname{120}
\expandafter\def\csname eePCTlv1981\endcsname{---}
\expandafter\def\csname eeGDP1982\endcsname{8,554}
\expandafter\def\csname eeGDPb1982\endcsname{13,635}
\expandafter\def\csname eePOP1982\endcsname{1,498.4}
\expandafter\def\csname eePCTfin1982\endcsname{63}
\expandafter\def\csname eePCTswe1982\endcsname{57}
\expandafter\def\csname eePCTwc1982\endcsname{59}
\expandafter\def\csname eePCTussr1982\endcsname{131}
\expandafter\def\csname eePCTlt1982\endcsname{119}
\expandafter\def\csname eePCTlv1982\endcsname{---}
\expandafter\def\csname eeGDP1983\endcsname{8,775}
\expandafter\def\csname eeGDPb1983\endcsname{13,988}
\expandafter\def\csname eePOP1983\endcsname{1,508.7}
\expandafter\def\csname eePCTfin1983\endcsname{64}
\expandafter\def\csname eePCTswe1983\endcsname{57}
\expandafter\def\csname eePCTwc1983\endcsname{59}
\expandafter\def\csname eePCTussr1983\endcsname{131}
\expandafter\def\csname eePCTlt1983\endcsname{120}
\expandafter\def\csname eePCTlv1983\endcsname{---}
\expandafter\def\csname eeGDP1984\endcsname{8,932}
\expandafter\def\csname eeGDPb1984\endcsname{14,238}
\expandafter\def\csname eePOP1984\endcsname{1,518.6}
\expandafter\def\csname eePCTfin1984\endcsname{63}
\expandafter\def\csname eePCTswe1984\endcsname{56}
\expandafter\def\csname eePCTwc1984\endcsname{59}
\expandafter\def\csname eePCTussr1984\endcsname{133}
\expandafter\def\csname eePCTlt1984\endcsname{116}
\expandafter\def\csname eePCTlv1984\endcsname{---}
\expandafter\def\csname eeGDP1985\endcsname{8,728}
\expandafter\def\csname eeGDPb1985\endcsname{13,913}
\expandafter\def\csname eePOP1985\endcsname{1,528.8}
\expandafter\def\csname eePCTfin1985\endcsname{60}
\expandafter\def\csname eePCTswe1985\endcsname{54}
\expandafter\def\csname eePCTwc1985\endcsname{56}
\expandafter\def\csname eePCTussr1985\endcsname{130}
\expandafter\def\csname eePCTlt1985\endcsname{113}
\expandafter\def\csname eePCTlv1985\endcsname{---}
\expandafter\def\csname eeGDP1986\endcsname{8,919}
\expandafter\def\csname eeGDPb1986\endcsname{14,216}
\expandafter\def\csname eePOP1986\endcsname{1,540.2}
\expandafter\def\csname eePCTfin1986\endcsname{60}
\expandafter\def\csname eePCTswe1986\endcsname{54}
\expandafter\def\csname eePCTwc1986\endcsname{56}
\expandafter\def\csname eePCTussr1986\endcsname{129}
\expandafter\def\csname eePCTlt1986\endcsname{114}
\expandafter\def\csname eePCTlv1986\endcsname{---}
\expandafter\def\csname eeGDP1987\endcsname{8,833}
\expandafter\def\csname eeGDPb1987\endcsname{14,079}
\expandafter\def\csname eePOP1987\endcsname{1,552.2}
\expandafter\def\csname eePCTfin1987\endcsname{57}
\expandafter\def\csname eePCTswe1987\endcsname{52}
\expandafter\def\csname eePCTwc1987\endcsname{54}
\expandafter\def\csname eePCTussr1987\endcsname{127}
\expandafter\def\csname eePCTlt1987\endcsname{114}
\expandafter\def\csname eePCTlv1987\endcsname{---}
\expandafter\def\csname eeGDP1988\endcsname{8,879}
\expandafter\def\csname eeGDPb1988\endcsname{14,153}
\expandafter\def\csname eePOP1988\endcsname{1,561.9}
\expandafter\def\csname eePCTfin1988\endcsname{55}
\expandafter\def\csname eePCTswe1988\endcsname{52}
\expandafter\def\csname eePCTwc1988\endcsname{53}
\expandafter\def\csname eePCTussr1988\endcsname{126}
\expandafter\def\csname eePCTlt1988\endcsname{113}
\expandafter\def\csname eePCTlv1988\endcsname{---}
\expandafter\def\csname eeGDP1989\endcsname{8,958}
\expandafter\def\csname eeGDPb1989\endcsname{14,280}
\expandafter\def\csname eePOP1989\endcsname{1,568.1}
\expandafter\def\csname eePCTfin1989\endcsname{53}
\expandafter\def\csname eePCTswe1989\endcsname{51}
\expandafter\def\csname eePCTwc1989\endcsname{52}
\expandafter\def\csname eePCTussr1989\endcsname{126}
\expandafter\def\csname eePCTlt1989\endcsname{114}
\expandafter\def\csname eePCTlv1989\endcsname{---}
\expandafter\def\csname eeGDP1990\endcsname{8,425}
\expandafter\def\csname eeGDPb1990\endcsname{13,429}
\expandafter\def\csname eePOP1990\endcsname{1,569.2}
\expandafter\def\csname eePCTfin1990\endcsname{50}
\expandafter\def\csname eePCTswe1990\endcsname{48}
\expandafter\def\csname eePCTwc1990\endcsname{48}
\expandafter\def\csname eePCTussr1990\endcsname{122}
\expandafter\def\csname eePCTlt1990\endcsname{113}
\expandafter\def\csname eePCTlv1990\endcsname{---}
\expandafter\def\csname eeGDP1991\endcsname{7,532}
\expandafter\def\csname eeGDPb1991\endcsname{12,006}
\expandafter\def\csname eePOP1991\endcsname{1,561.3}
\expandafter\def\csname eePCTfin1991\endcsname{48}
\expandafter\def\csname eePCTswe1991\endcsname{43}
\expandafter\def\csname eePCTwc1991\endcsname{42}
\expandafter\def\csname eePCTussr1991\endcsname{116}
\expandafter\def\csname eePCTlt1991\endcsname{104}
\expandafter\def\csname eePCTlv1991\endcsname{---}
\expandafter\def\csname eeGDP1992\endcsname{6,069}
\expandafter\def\csname eeGDPb1992\endcsname{9,675}
\expandafter\def\csname eePOP1992\endcsname{1,533.1}
\expandafter\def\csname eePCTfin1992\endcsname{40}
\expandafter\def\csname eePCTswe1992\endcsname{36}
\expandafter\def\csname eePCTwc1992\endcsname{34}
\expandafter\def\csname eePCTussr1992\endcsname{108}
\expandafter\def\csname eePCTlt1992\endcsname{101}
\expandafter\def\csname eePCTlv1992\endcsname{---}
\expandafter\def\csname eeGDP1993\endcsname{5,870}
\expandafter\def\csname eeGDPb1993\endcsname{9,357}
\expandafter\def\csname eePOP1993\endcsname{1,494.1}
\expandafter\def\csname eePCTfin1993\endcsname{39}
\expandafter\def\csname eePCTswe1993\endcsname{35}
\expandafter\def\csname eePCTwc1993\endcsname{33}
\expandafter\def\csname eePCTussr1993\endcsname{114}
\expandafter\def\csname eePCTlt1993\endcsname{108}
\expandafter\def\csname eePCTlv1993\endcsname{---}
\expandafter\def\csname eeGDP1994\endcsname{5,899}
\expandafter\def\csname eeGDPb1994\endcsname{9,402}
\expandafter\def\csname eePOP1994\endcsname{1,462.5}
\expandafter\def\csname eePCTfin1994\endcsname{38}
\expandafter\def\csname eePCTswe1994\endcsname{34}
\expandafter\def\csname eePCTwc1994\endcsname{32}
\expandafter\def\csname eePCTussr1994\endcsname{131}
\expandafter\def\csname eePCTlt1994\endcsname{118}
\expandafter\def\csname eePCTlv1994\endcsname{---}
\expandafter\def\csname eeGDP1995\endcsname{6,277}
\expandafter\def\csname eeGDPb1995\endcsname{10,005}
\expandafter\def\csname eePOP1995\endcsname{1,436.6}
\expandafter\def\csname eePCTfin1995\endcsname{39}
\expandafter\def\csname eePCTswe1995\endcsname{35}
\expandafter\def\csname eePCTwc1995\endcsname{33}
\expandafter\def\csname eePCTussr1995\endcsname{145}
\expandafter\def\csname eePCTlt1995\endcsname{121}
\expandafter\def\csname eePCTlv1995\endcsname{125}
\expandafter\def\csname eeGDP1996\endcsname{6,685}
\expandafter\def\csname eeGDPb1996\endcsname{10,655}
\expandafter\def\csname eePOP1996\endcsname{1,415.6}
\expandafter\def\csname eePCTfin1996\endcsname{40}
\expandafter\def\csname eePCTswe1996\endcsname{37}
\expandafter\def\csname eePCTwc1996\endcsname{35}
\expandafter\def\csname eePCTussr1996\endcsname{157}
\expandafter\def\csname eePCTlt1996\endcsname{121}
\expandafter\def\csname eePCTlv1996\endcsname{128}
\expandafter\def\csname eeGDP1997\endcsname{7,644}
\expandafter\def\csname eeGDPb1997\endcsname{12,184}
\expandafter\def\csname eePOP1997\endcsname{1,399.5}
\expandafter\def\csname eePCTfin1997\endcsname{43}
\expandafter\def\csname eePCTswe1997\endcsname{41}
\expandafter\def\csname eePCTwc1997\endcsname{39}
\expandafter\def\csname eePCTussr1997\endcsname{173}
\expandafter\def\csname eePCTlt1997\endcsname{127}
\expandafter\def\csname eePCTlv1997\endcsname{134}
\expandafter\def\csname eeGDP1998\endcsname{8,053}
\expandafter\def\csname eeGDPb1998\endcsname{12,836}
\expandafter\def\csname eePOP1998\endcsname{1,386.2}
\expandafter\def\csname eePCTfin1998\endcsname{43}
\expandafter\def\csname eePCTswe1998\endcsname{41}
\expandafter\def\csname eePCTwc1998\endcsname{40}
\expandafter\def\csname eePCTussr1998\endcsname{184}
\expandafter\def\csname eePCTlt1998\endcsname{124}
\expandafter\def\csname eePCTlv1998\endcsname{131}
\expandafter\def\csname eeGDP1999\endcsname{7,995}
\expandafter\def\csname eeGDPb1999\endcsname{12,744}
\expandafter\def\csname eePOP1999\endcsname{1,390.2}
\expandafter\def\csname eePCTfin1999\endcsname{41}
\expandafter\def\csname eePCTswe1999\endcsname{39}
\expandafter\def\csname eePCTwc1999\endcsname{39}
\expandafter\def\csname eePCTussr1999\endcsname{172}
\expandafter\def\csname eePCTlt1999\endcsname{123}
\expandafter\def\csname eePCTlv1999\endcsname{126}
\expandafter\def\csname eeGDP2000\endcsname{8,759}
\expandafter\def\csname eeGDPb2000\endcsname{13,961}
\expandafter\def\csname eePOP2000\endcsname{1,397.0}
\expandafter\def\csname eePCTfin2000\endcsname{43}
\expandafter\def\csname eePCTswe2000\endcsname{41}
\expandafter\def\csname eePCTwc2000\endcsname{41}
\expandafter\def\csname eePCTussr2000\endcsname{170}
\expandafter\def\csname eePCTlt2000\endcsname{129}
\expandafter\def\csname eePCTlv2000\endcsname{129}
\expandafter\def\csname eeGDP2001\endcsname{9,344}
\expandafter\def\csname eeGDPb2001\endcsname{14,894}
\expandafter\def\csname eePOP2001\endcsname{1,388.1}
\expandafter\def\csname eePCTfin2001\endcsname{44}
\expandafter\def\csname eePCTswe2001\endcsname{43}
\expandafter\def\csname eePCTwc2001\endcsname{43}
\expandafter\def\csname eePCTussr2001\endcsname{168}
\expandafter\def\csname eePCTlt2001\endcsname{128}
\expandafter\def\csname eePCTlv2001\endcsname{128}
\expandafter\def\csname eeGDP2002\endcsname{10,040}
\expandafter\def\csname eeGDPb2002\endcsname{16,004}
\expandafter\def\csname eePOP2002\endcsname{1,379.3}
\expandafter\def\csname eePCTfin2002\endcsname{47}
\expandafter\def\csname eePCTswe2002\endcsname{45}
\expandafter\def\csname eePCTwc2002\endcsname{45}
\expandafter\def\csname eePCTussr2002\endcsname{169}
\expandafter\def\csname eePCTlt2002\endcsname{128}
\expandafter\def\csname eePCTlv2002\endcsname{127}
\expandafter\def\csname eeGDP2003\endcsname{10,871}
\expandafter\def\csname eeGDPb2003\endcsname{17,328}
\expandafter\def\csname eePOP2003\endcsname{1,370.7}
\expandafter\def\csname eePCTfin2003\endcsname{50}
\expandafter\def\csname eePCTswe2003\endcsname{48}
\expandafter\def\csname eePCTwc2003\endcsname{49}
\expandafter\def\csname eePCTussr2003\endcsname{167}
\expandafter\def\csname eePCTlt2003\endcsname{125}
\expandafter\def\csname eePCTlv2003\endcsname{125}
\expandafter\def\csname eeGDP2004\endcsname{11,680}
\expandafter\def\csname eeGDPb2004\endcsname{18,618}
\expandafter\def\csname eePOP2004\endcsname{1,362.6}
\expandafter\def\csname eePCTfin2004\endcsname{52}
\expandafter\def\csname eePCTswe2004\endcsname{49}
\expandafter\def\csname eePCTwc2004\endcsname{51}
\expandafter\def\csname eePCTussr2004\endcsname{163}
\expandafter\def\csname eePCTlt2004\endcsname{124}
\expandafter\def\csname eePCTlv2004\endcsname{123}
\expandafter\def\csname eeGDP2005\endcsname{12,867}
\expandafter\def\csname eeGDPb2005\endcsname{20,509}
\expandafter\def\csname eePOP2005\endcsname{1,354.8}
\expandafter\def\csname eePCTfin2005\endcsname{56}
\expandafter\def\csname eePCTswe2005\endcsname{52}
\expandafter\def\csname eePCTwc2005\endcsname{55}
\expandafter\def\csname eePCTussr2005\endcsname{165}
\expandafter\def\csname eePCTlt2005\endcsname{125}
\expandafter\def\csname eePCTlv2005\endcsname{121}
\expandafter\def\csname eeGDP2006\endcsname{14,207}
\expandafter\def\csname eeGDPb2006\endcsname{22,645}
\expandafter\def\csname eePOP2006\endcsname{1,346.8}
\expandafter\def\csname eePCTfin2006\endcsname{59}
\expandafter\def\csname eePCTswe2006\endcsname{55}
\expandafter\def\csname eePCTwc2006\endcsname{59}
\expandafter\def\csname eePCTussr2006\endcsname{165}
\expandafter\def\csname eePCTlt2006\endcsname{126}
\expandafter\def\csname eePCTlv2006\endcsname{118}
\expandafter\def\csname eeGDP2007\endcsname{15,353}
\expandafter\def\csname eeGDPb2007\endcsname{24,473}
\expandafter\def\csname eePOP2007\endcsname{1,340.7}
\expandafter\def\csname eePCTfin2007\endcsname{61}
\expandafter\def\csname eePCTswe2007\endcsname{58}
\expandafter\def\csname eePCTwc2007\endcsname{61}
\expandafter\def\csname eePCTussr2007\endcsname{161}
\expandafter\def\csname eePCTlt2007\endcsname{122}
\expandafter\def\csname eePCTlv2007\endcsname{115}
\expandafter\def\csname eeGDP2008\endcsname{14,604}
\expandafter\def\csname eeGDPb2008\endcsname{23,279}
\expandafter\def\csname eePOP2008\endcsname{1,337.1}
\expandafter\def\csname eePCTfin2008\endcsname{58}
\expandafter\def\csname eePCTswe2008\endcsname{55}
\expandafter\def\csname eePCTwc2008\endcsname{58}
\expandafter\def\csname eePCTussr2008\endcsname{143}
\expandafter\def\csname eePCTlt2008\endcsname{111}
\expandafter\def\csname eePCTlv2008\endcsname{112}
\expandafter\def\csname eeGDP2009\endcsname{12,492}
\expandafter\def\csname eeGDPb2009\endcsname{19,912}
\expandafter\def\csname eePOP2009\endcsname{1,334.5}
\expandafter\def\csname eePCTfin2009\endcsname{54}
\expandafter\def\csname eePCTswe2009\endcsname{50}
\expandafter\def\csname eePCTwc2009\endcsname{52}
\expandafter\def\csname eePCTussr2009\endcsname{129}
\expandafter\def\csname eePCTlt2009\endcsname{111}
\expandafter\def\csname eePCTlv2009\endcsname{110}
\expandafter\def\csname eeGDP2010\endcsname{12,826}
\expandafter\def\csname eeGDPb2010\endcsname{20,445}
\expandafter\def\csname eePOP2010\endcsname{1,331.5}
\expandafter\def\csname eePCTfin2010\endcsname{54}
\expandafter\def\csname eePCTswe2010\endcsname{48}
\expandafter\def\csname eePCTwc2010\endcsname{52}
\expandafter\def\csname eePCTussr2010\endcsname{125}
\expandafter\def\csname eePCTlt2010\endcsname{110}
\expandafter\def\csname eePCTlv2010\endcsname{116}
\expandafter\def\csname eeGDP2011\endcsname{13,800}
\expandafter\def\csname eeGDPb2011\endcsname{21,997}
\expandafter\def\csname eePOP2011\endcsname{1,327.4}
\expandafter\def\csname eePCTfin2011\endcsname{57}
\expandafter\def\csname eePCTswe2011\endcsname{52}
\expandafter\def\csname eePCTwc2011\endcsname{54}
\expandafter\def\csname eePCTussr2011\endcsname{126}
\expandafter\def\csname eePCTlt2011\endcsname{109}
\expandafter\def\csname eePCTlv2011\endcsname{119}
\expandafter\def\csname eeGDP2012\endcsname{14,294}
\expandafter\def\csname eeGDPb2012\endcsname{22,784}
\expandafter\def\csname eePOP2012\endcsname{1,322.7}
\expandafter\def\csname eePCTfin2012\endcsname{60}
\expandafter\def\csname eePCTswe2012\endcsname{55}
\expandafter\def\csname eePCTwc2012\endcsname{56}
\expandafter\def\csname eePCTussr2012\endcsname{127}
\expandafter\def\csname eePCTlt2012\endcsname{107}
\expandafter\def\csname eePCTlv2012\endcsname{114}
\expandafter\def\csname eeGDP2013\endcsname{14,557}
\expandafter\def\csname eeGDPb2013\endcsname{23,204}
\expandafter\def\csname eePOP2013\endcsname{1,318.0}
\expandafter\def\csname eePCTfin2013\endcsname{62}
\expandafter\def\csname eePCTswe2013\endcsname{56}
\expandafter\def\csname eePCTwc2013\endcsname{57}
\expandafter\def\csname eePCTussr2013\endcsname{127}
\expandafter\def\csname eePCTlt2013\endcsname{104}
\expandafter\def\csname eePCTlv2013\endcsname{113}
\expandafter\def\csname eeGDP2014\endcsname{15,046}
\expandafter\def\csname eeGDPb2014\endcsname{23,983}
\expandafter\def\csname eePOP2014\endcsname{1,314.7}
\expandafter\def\csname eePCTfin2014\endcsname{65}
\expandafter\def\csname eePCTswe2014\endcsname{57}
\expandafter\def\csname eePCTwc2014\endcsname{58}
\expandafter\def\csname eePCTussr2014\endcsname{131}
\expandafter\def\csname eePCTlt2014\endcsname{103}
\expandafter\def\csname eePCTlv2014\endcsname{113}
\expandafter\def\csname eeGDP2015\endcsname{15,354}
\expandafter\def\csname eeGDPb2015\endcsname{24,474}
\expandafter\def\csname eePOP2015\endcsname{1,314.0}
\expandafter\def\csname eePCTfin2015\endcsname{66}
\expandafter\def\csname eePCTswe2015\endcsname{56}
\expandafter\def\csname eePCTwc2015\endcsname{58}
\expandafter\def\csname eePCTussr2015\endcsname{137}
\expandafter\def\csname eePCTlt2015\endcsname{102}
\expandafter\def\csname eePCTlv2015\endcsname{110}
\expandafter\def\csname eeGDP2016\endcsname{15,807}
\expandafter\def\csname eeGDPb2016\endcsname{25,196}
\expandafter\def\csname eePOP2016\endcsname{1,314.4}
\expandafter\def\csname eePCTfin2016\endcsname{66}
\expandafter\def\csname eePCTswe2016\endcsname{57}
\expandafter\def\csname eePCTwc2016\endcsname{60}
\expandafter\def\csname eePCTussr2016\endcsname{140}
\expandafter\def\csname eePCTlt2016\endcsname{101}
\expandafter\def\csname eePCTlv2016\endcsname{110}
\expandafter\def\csname eeGDP2017\endcsname{16,726}
\expandafter\def\csname eeGDPb2017\endcsname{26,662}
\expandafter\def\csname eePOP2017\endcsname{1,316.1}
\expandafter\def\csname eePCTfin2017\endcsname{68}
\expandafter\def\csname eePCTswe2017\endcsname{60}
\expandafter\def\csname eePCTwc2017\endcsname{62}
\expandafter\def\csname eePCTussr2017\endcsname{145}
\expandafter\def\csname eePCTlt2017\endcsname{101}
\expandafter\def\csname eePCTlv2017\endcsname{112}
\expandafter\def\csname eeGDP2018\endcsname{17,313}
\expandafter\def\csname eeGDPb2018\endcsname{27,596}
\expandafter\def\csname eePOP2018\endcsname{1,320.6}
\expandafter\def\csname eePCTfin2018\endcsname{70}
\expandafter\def\csname eePCTswe2018\endcsname{61}
\expandafter\def\csname eePCTwc2018\endcsname{63}
\expandafter\def\csname eePCTussr2018\endcsname{146}
\expandafter\def\csname eePCTlt2018\endcsname{100}
\expandafter\def\csname eePCTlv2018\endcsname{110}
\expandafter\def\csname eeGDP2019\endcsname{17,883}
\expandafter\def\csname eeGDPb2019\endcsname{28,506}
\expandafter\def\csname eePOP2019\endcsname{1,325.1}
\expandafter\def\csname eePCTfin2019\endcsname{71}
\expandafter\def\csname eePCTswe2019\endcsname{63}
\expandafter\def\csname eePCTwc2019\endcsname{65}
\expandafter\def\csname eePCTussr2019\endcsname{147}
\expandafter\def\csname eePCTlt2019\endcsname{98}
\expandafter\def\csname eePCTlv2019\endcsname{110}
\expandafter\def\csname eeGDP2020\endcsname{17,730}
\expandafter\def\csname eeGDPb2020\endcsname{28,262}
\expandafter\def\csname eePOP2020\endcsname{1,329.5}
\expandafter\def\csname eePCTfin2020\endcsname{73}
\expandafter\def\csname eePCTswe2020\endcsname{64}
\expandafter\def\csname eePCTwc2020\endcsname{68}
\expandafter\def\csname eePCTussr2020\endcsname{150}
\expandafter\def\csname eePCTlt2020\endcsname{98}
\expandafter\def\csname eePCTlv2020\endcsname{111}
\expandafter\def\csname eeGDP2021\endcsname{19,134}
\expandafter\def\csname eeGDPb2021\endcsname{30,499}
\expandafter\def\csname eePOP2021\endcsname{1,330.9}
\expandafter\def\csname eePCTfin2021\endcsname{76}
\expandafter\def\csname eePCTswe2021\endcsname{66}
\expandafter\def\csname eePCTwc2021\endcsname{71}
\expandafter\def\csname eePCTussr2021\endcsname{153}
\expandafter\def\csname eePCTlt2021\endcsname{100}
\expandafter\def\csname eePCTlv2021\endcsname{114}
\expandafter\def\csname eeGDP2022\endcsname{18,863}
\expandafter\def\csname eeGDPb2022\endcsname{30,067}
\expandafter\def\csname eePOP2022\endcsname{1,348.8}
\expandafter\def\csname eePCTfin2022\endcsname{74}
\expandafter\def\csname eePCTswe2022\endcsname{64}
\expandafter\def\csname eePCTwc2022\endcsname{68}
\expandafter\def\csname eePCTussr2022\endcsname{155}
\expandafter\def\csname eePCTlt2022\endcsname{97}
\expandafter\def\csname eePCTlv2022\endcsname{110}
\expandafter\def\csname eeGR1913x1938\endcsname{1.1}
\expandafter\def\csname eeCUM1913x1938\endcsname{30}
\expandafter\def\csname GRFinlandx1913x1938\endcsname{2.1}
\expandafter\def\csname GRSwedenx1913x1938\endcsname{2.2}
\expandafter\def\csname GRDenmarkx1913x1938\endcsname{1.6}
\expandafter\def\csname GRNordicsx1913x1938\endcsname{1.9}
\expandafter\def\csname GRWesternx1913x1938\endcsname{1.0}
\expandafter\def\csname GRPeripheryx1913x1938\endcsname{1.2}
\expandafter\def\csname GRUSSRx1913x1938\endcsname{1.7}
\expandafter\def\csname GRRussiax1913x1938\endcsname{2.7}
\expandafter\def\csname GRECEx1913x1938\endcsname{0.8}
\expandafter\def\csname GRPolandx1913x1938\endcsname{0.9}
\expandafter\def\csname GRCzechoslovakiax1913x1938\endcsname{1.2}
\expandafter\def\csname GRHungaryx1913x1938\endcsname{0.9}
\expandafter\def\csname GRAustriax1913x1938\endcsname{0.1}
\expandafter\def\csname GRIrelandx1913x1938\endcsname{0.4}
\expandafter\def\csname GRGreecex1913x1938\endcsname{2.5}
\expandafter\def\csname GRLithuaniax1913x1938\endcsname{0.9}
\expandafter\def\csname GRLatviax1913x1938\endcsname{1.1}
\expandafter\def\csname eeGR1919x1929\endcsname{2.8}
\expandafter\def\csname eeCUM1919x1929\endcsname{32}
\expandafter\def\csname GRFinlandx1919x1929\endcsname{5.1}
\expandafter\def\csname GRSwedenx1919x1929\endcsname{3.7}
\expandafter\def\csname GRDenmarkx1919x1929\endcsname{2.8}
\expandafter\def\csname GRNordicsx1919x1929\endcsname{3.3}
\expandafter\def\csname GRWesternx1919x1929\endcsname{3.5}
\expandafter\def\csname GRPeripheryx1919x1929\endcsname{3.1}
\expandafter\def\csname GRUSSRx1919x1929\endcsname{9.2}
\expandafter\def\csname GRRussiax1919x1929\endcsname{7.0}
\expandafter\def\csname GRECEx1919x1929\endcsname{4.8}
\expandafter\def\csname GRPolandx1919x1929\endcsname{4.6}
\expandafter\def\csname GRCzechoslovakiax1919x1929\endcsname{6.7}
\expandafter\def\csname GRHungaryx1919x1929\endcsname{5.3}
\expandafter\def\csname GRAustriax1919x1929\endcsname{5.1}
\expandafter\def\csname GRIrelandx1919x1929\endcsname{0.8}
\expandafter\def\csname GRGreecex1919x1929\endcsname{6.3}
\expandafter\def\csname GRLithuaniax1919x1929\endcsname{4.5}
\expandafter\def\csname GRLatviax1919x1929\endcsname{4.7}
\expandafter\def\csname eeGR1929x1938\endcsname{2.7}
\expandafter\def\csname eeCUM1929x1938\endcsname{27}
\expandafter\def\csname GRFinlandx1929x1938\endcsname{3.1}
\expandafter\def\csname GRSwedenx1929x1938\endcsname{2.1}
\expandafter\def\csname GRDenmarkx1929x1938\endcsname{1.4}
\expandafter\def\csname GRNordicsx1929x1938\endcsname{2.1}
\expandafter\def\csname GRWesternx1929x1938\endcsname{1.1}
\expandafter\def\csname GRPeripheryx1929x1938\endcsname{0.8}
\expandafter\def\csname GRUSSRx1929x1938\endcsname{5.0}
\expandafter\def\csname GRRussiax1929x1938\endcsname{6.5}
\expandafter\def\csname GRECEx1929x1938\endcsname{0.1}
\expandafter\def\csname GRPolandx1929x1938\endcsname{0.3}
\expandafter\def\csname GRCzechoslovakiax1929x1938\endcsname{$-$0.9}
\expandafter\def\csname GRHungaryx1929x1938\endcsname{0.8}
\expandafter\def\csname GRAustriax1929x1938\endcsname{$-$0.4}
\expandafter\def\csname GRIrelandx1929x1938\endcsname{0.9}
\expandafter\def\csname GRGreecex1929x1938\endcsname{1.5}
\expandafter\def\csname GRLithuaniax1929x1938\endcsname{3.0}
\expandafter\def\csname GRLatviax1929x1938\endcsname{3.1}
\expandafter\def\csname eeGR1938x1950\endcsname{2.6}
\expandafter\def\csname eeCUM1938x1950\endcsname{37}
\expandafter\def\csname GRFinlandx1938x1950\endcsname{1.4}
\expandafter\def\csname GRSwedenx1938x1950\endcsname{2.7}
\expandafter\def\csname GRDenmarkx1938x1950\endcsname{1.6}
\expandafter\def\csname GRNordicsx1938x1950\endcsname{2.1}
\expandafter\def\csname GRWesternx1938x1950\endcsname{0.6}
\expandafter\def\csname GRPeripheryx1938x1950\endcsname{$-$0.1}
\expandafter\def\csname GRUSSRx1938x1950\endcsname{2.4}
\expandafter\def\csname GRRussiax1938x1950\endcsname{3.1}
\expandafter\def\csname GRECEx1938x1950\endcsname{0.8}
\expandafter\def\csname GRPolandx1938x1950\endcsname{1.0}
\expandafter\def\csname GRCzechoslovakiax1938x1950\endcsname{1.9}
\expandafter\def\csname GRHungaryx1938x1950\endcsname{$-$0.6}
\expandafter\def\csname GRAustriax1938x1950\endcsname{0.3}
\expandafter\def\csname GRIrelandx1938x1950\endcsname{1.0}
\expandafter\def\csname GRGreecex1938x1950\endcsname{$-$2.8}
\expandafter\def\csname GRLithuaniax1938x1950\endcsname{0.9}
\expandafter\def\csname GRLatviax1938x1950\endcsname{1.5}
\expandafter\def\csname eeGR1950x1968\endcsname{3.9}
\expandafter\def\csname eeCUM1950x1968\endcsname{100}
\expandafter\def\csname GRFinlandx1950x1968\endcsname{3.6}
\expandafter\def\csname GRSwedenx1950x1968\endcsname{3.0}
\expandafter\def\csname GRDenmarkx1950x1968\endcsname{3.0}
\expandafter\def\csname GRNordicsx1950x1968\endcsname{3.1}
\expandafter\def\csname GRWesternx1950x1968\endcsname{3.6}
\expandafter\def\csname GRPeripheryx1950x1968\endcsname{5.2}
\expandafter\def\csname GRUSSRx1950x1968\endcsname{3.4}
\expandafter\def\csname GRRussiax1950x1968\endcsname{5.3}
\expandafter\def\csname GRECEx1950x1968\endcsname{3.9}
\expandafter\def\csname GRPolandx1950x1968\endcsname{3.2}
\expandafter\def\csname GRCzechoslovakiax1950x1968\endcsname{3.2}
\expandafter\def\csname GRHungaryx1950x1968\endcsname{3.9}
\expandafter\def\csname GRAustriax1950x1968\endcsname{4.8}
\expandafter\def\csname GRIrelandx1950x1968\endcsname{2.9}
\expandafter\def\csname GRGreecex1950x1968\endcsname{5.8}
\expandafter\def\csname GRLithuaniax1950x1968\endcsname{4.9}
\expandafter\def\csname GRLatviax1950x1968\endcsname{4.0}
\expandafter\def\csname eeGR1968x1989\endcsname{0.9}
\expandafter\def\csname eeCUM1968x1989\endcsname{22}
\expandafter\def\csname GRFinlandx1968x1989\endcsname{3.6}
\expandafter\def\csname GRSwedenx1968x1989\endcsname{2.0}
\expandafter\def\csname GRDenmarkx1968x1989\endcsname{2.1}
\expandafter\def\csname GRNordicsx1968x1989\endcsname{2.5}
\expandafter\def\csname GRWesternx1968x1989\endcsname{2.4}
\expandafter\def\csname GRPeripheryx1968x1989\endcsname{3.0}
\expandafter\def\csname GRUSSRx1968x1989\endcsname{1.5}
\expandafter\def\csname GRRussiax1968x1989\endcsname{2.0}
\expandafter\def\csname GRECEx1968x1989\endcsname{1.7}
\expandafter\def\csname GRPolandx1968x1989\endcsname{1.3}
\expandafter\def\csname GRCzechoslovakiax1968x1989\endcsname{1.6}
\expandafter\def\csname GRHungaryx1968x1989\endcsname{1.6}
\expandafter\def\csname GRAustriax1968x1989\endcsname{3.1}
\expandafter\def\csname GRIrelandx1968x1989\endcsname{3.1}
\expandafter\def\csname GRGreecex1968x1989\endcsname{3.2}
\expandafter\def\csname GRLithuaniax1968x1989\endcsname{1.8}
\expandafter\def\csname GRLatviax1968x1989\endcsname{0.7}
\expandafter\def\csname eeGR1938x1989\endcsname{2.4}
\expandafter\def\csname eeCUM1938x1989\endcsname{234}
\expandafter\def\csname GRFinlandx1938x1989\endcsname{3.1}
\expandafter\def\csname GRSwedenx1938x1989\endcsname{2.5}
\expandafter\def\csname GRDenmarkx1938x1989\endcsname{2.3}
\expandafter\def\csname GRNordicsx1938x1989\endcsname{2.6}
\expandafter\def\csname GRWesternx1938x1989\endcsname{2.4}
\expandafter\def\csname GRPeripheryx1938x1989\endcsname{3.0}
\expandafter\def\csname GRUSSRx1938x1989\endcsname{2.4}
\expandafter\def\csname GRRussiax1938x1989\endcsname{3.4}
\expandafter\def\csname GRECEx1938x1989\endcsname{2.3}
\expandafter\def\csname GRPolandx1938x1989\endcsname{1.9}
\expandafter\def\csname GRCzechoslovakiax1938x1989\endcsname{2.3}
\expandafter\def\csname GRHungaryx1938x1989\endcsname{1.9}
\expandafter\def\csname GRAustriax1938x1989\endcsname{3.0}
\expandafter\def\csname GRIrelandx1938x1989\endcsname{2.5}
\expandafter\def\csname GRGreecex1938x1989\endcsname{2.6}
\expandafter\def\csname GRLithuaniax1938x1989\endcsname{2.7}
\expandafter\def\csname GRLatviax1938x1989\endcsname{2.1}
\expandafter\def\csname eeGR1989x1995\endcsname{$-$5.8}
\expandafter\def\csname eeCUM1989x1995\endcsname{-30}
\expandafter\def\csname GRFinlandx1989x1995\endcsname{$-$0.8}
\expandafter\def\csname GRSwedenx1989x1995\endcsname{0.3}
\expandafter\def\csname GRDenmarkx1989x1995\endcsname{2.3}
\expandafter\def\csname GRNordicsx1989x1995\endcsname{1.6}
\expandafter\def\csname GRWesternx1989x1995\endcsname{1.3}
\expandafter\def\csname GRPeripheryx1989x1995\endcsname{2.1}
\expandafter\def\csname GRUSSRx1989x1995\endcsname{$-$8.0}
\expandafter\def\csname GRRussiax1989x1995\endcsname{$-$6.4}
\expandafter\def\csname GRECEx1989x1995\endcsname{$-$1.6}
\expandafter\def\csname GRPolandx1989x1995\endcsname{0.6}
\expandafter\def\csname GRCzechoslovakiax1989x1995\endcsname{$-$1.0}
\expandafter\def\csname GRHungaryx1989x1995\endcsname{$-$1.4}
\expandafter\def\csname GRAustriax1989x1995\endcsname{2.1}
\expandafter\def\csname GRIrelandx1989x1995\endcsname{5.8}
\expandafter\def\csname GRGreecex1989x1995\endcsname{1.0}
\expandafter\def\csname GRLithuaniax1989x1995\endcsname{$-$6.6}
\expandafter\def\csname GRLatviax1989x1995\endcsname{$-$8.9}
\expandafter\def\csname eeGR1995x2022\endcsname{4.2}
\expandafter\def\csname eeCUM1995x2022\endcsname{201}
\expandafter\def\csname GRFinlandx1995x2022\endcsname{1.7}
\expandafter\def\csname GRSwedenx1995x2022\endcsname{1.9}
\expandafter\def\csname GRDenmarkx1995x2022\endcsname{1.6}
\expandafter\def\csname GRNordicsx1995x2022\endcsname{2.1}
\expandafter\def\csname GRWesternx1995x2022\endcsname{1.5}
\expandafter\def\csname GRPeripheryx1995x2022\endcsname{1.4}
\expandafter\def\csname GRUSSRx1995x2022\endcsname{3.9}
\expandafter\def\csname GRRussiax1995x2022\endcsname{4.1}
\expandafter\def\csname GRECEx1995x2022\endcsname{4.4}
\expandafter\def\csname GRPolandx1995x2022\endcsname{4.7}
\expandafter\def\csname GRCzechoslovakiax1995x2022\endcsname{3.2}
\expandafter\def\csname GRHungaryx1995x2022\endcsname{4.0}
\expandafter\def\csname GRAustriax1995x2022\endcsname{1.5}
\expandafter\def\csname GRIrelandx1995x2022\endcsname{3.4}
\expandafter\def\csname GRGreecex1995x2022\endcsname{1.4}
\expandafter\def\csname GRLithuaniax1995x2022\endcsname{5.0}
\expandafter\def\csname GRLatviax1995x2022\endcsname{4.6}
\expandafter\def\csname eeGR1913x1919\endcsname{$-$4.1}
\expandafter\def\csname eeCUM1913x1919\endcsname{-22}
\expandafter\def\csname GRFinlandx1913x1919\endcsname{$-$3.9}
\expandafter\def\csname GRSwedenx1913x1919\endcsname{$-$0.3}
\expandafter\def\csname GRDenmarkx1913x1919\endcsname{$-$0.2}
\expandafter\def\csname GRNordicsx1913x1919\endcsname{$-$0.6}
\expandafter\def\csname GRWesternx1913x1919\endcsname{$-$3.0}
\expandafter\def\csname GRPeripheryx1913x1919\endcsname{$-$1.3}
\expandafter\def\csname GRUSSRx1913x1919\endcsname{$-$14.0}
\expandafter\def\csname GRRussiax1913x1919\endcsname{$-$8.9}
\expandafter\def\csname GRECEx1913x1919\endcsname{$-$4.3}
\expandafter\def\csname GRPolandx1913x1919\endcsname{$-$4.2}
\expandafter\def\csname GRCzechoslovakiax1913x1919\endcsname{$-$4.5}
\expandafter\def\csname GRHungaryx1913x1919\endcsname{$-$5.6}
\expandafter\def\csname GRAustriax1913x1919\endcsname{$-$6.9}
\expandafter\def\csname GRIrelandx1913x1919\endcsname{$-$0.9}
\expandafter\def\csname GRGreecex1913x1919\endcsname{$-$2.2}
\expandafter\def\csname GRLithuaniax1913x1919\endcsname{$-$7.6}
\expandafter\def\csname GRLatviax1913x1919\endcsname{$-$7.5}
\expandafter\def\csname eeGR1919x1938\endcsname{2.8}
\expandafter\def\csname eeCUM1919x1938\endcsname{68}
\expandafter\def\csname GRFinlandx1919x1938\endcsname{4.1}
\expandafter\def\csname GRSwedenx1919x1938\endcsname{2.9}
\expandafter\def\csname GRDenmarkx1919x1938\endcsname{2.1}
\expandafter\def\csname GRNordicsx1919x1938\endcsname{2.7}
\expandafter\def\csname GRWesternx1919x1938\endcsname{2.3}
\expandafter\def\csname GRPeripheryx1919x1938\endcsname{2.0}
\expandafter\def\csname GRUSSRx1919x1938\endcsname{7.2}
\expandafter\def\csname GRRussiax1919x1938\endcsname{6.7}
\expandafter\def\csname GRECEx1919x1938\endcsname{2.5}
\expandafter\def\csname GRPolandx1919x1938\endcsname{2.6}
\expandafter\def\csname GRCzechoslovakiax1919x1938\endcsname{3.0}
\expandafter\def\csname GRHungaryx1919x1938\endcsname{3.1}
\expandafter\def\csname GRAustriax1919x1938\endcsname{2.4}
\expandafter\def\csname GRIrelandx1919x1938\endcsname{0.9}
\expandafter\def\csname GRGreecex1919x1938\endcsname{4.0}
\expandafter\def\csname GRLithuaniax1919x1938\endcsname{3.8}
\expandafter\def\csname GRLatviax1919x1938\endcsname{4.0}
\expandafter\def\csname eeGR1940x1944\endcsname{$-$8.9}
\expandafter\def\csname eeCUM1940x1944\endcsname{-31}
\expandafter\def\csname GRFinlandx1940x1944\endcsname{3.4}
\expandafter\def\csname GRSwedenx1940x1944\endcsname{0.8}
\expandafter\def\csname GRDenmarkx1940x1944\endcsname{2.0}
\expandafter\def\csname GRNordicsx1940x1944\endcsname{1.1}
\expandafter\def\csname GRWesternx1940x1944\endcsname{$-$0.2}
\expandafter\def\csname GRPeripheryx1940x1944\endcsname{$-$7.8}
\expandafter\def\csname GRUSSRx1940x1944\endcsname{$-$1.6}
\expandafter\def\csname GRRussiax1940x1944\endcsname{$-$1.4}
\expandafter\def\csname GRECEx1940x1944\endcsname{$-$7.0}
\expandafter\def\csname GRPolandx1940x1944\endcsname{$-$11.6}
\expandafter\def\csname GRCzechoslovakiax1940x1944\endcsname{$-$5.1}
\expandafter\def\csname GRHungaryx1940x1944\endcsname{$-$4.3}
\expandafter\def\csname GRAustriax1940x1944\endcsname{1.2}
\expandafter\def\csname GRIrelandx1940x1944\endcsname{0.0}
\expandafter\def\csname GRGreecex1940x1944\endcsname{$-$15.8}
\expandafter\def\csname GRLithuaniax1940x1944\endcsname{$-$11.6}
\expandafter\def\csname GRLatviax1940x1944\endcsname{$-$7.1}
\expandafter\def\csname eeGR1944x1950\endcsname{11.5}
\expandafter\def\csname eeCUM1944x1950\endcsname{92}
\expandafter\def\csname GRFinlandx1944x1950\endcsname{2.4}
\expandafter\def\csname GRSwedenx1944x1950\endcsname{5.1}
\expandafter\def\csname GRDenmarkx1944x1950\endcsname{3.8}
\expandafter\def\csname GRNordicsx1944x1950\endcsname{4.6}
\expandafter\def\csname GRWesternx1944x1950\endcsname{0.8}
\expandafter\def\csname GRPeripheryx1944x1950\endcsname{5.6}
\expandafter\def\csname GRUSSRx1944x1950\endcsname{5.9}
\expandafter\def\csname GRRussiax1944x1950\endcsname{5.6}
\expandafter\def\csname GRECEx1944x1950\endcsname{8.0}
\expandafter\def\csname GRPolandx1944x1950\endcsname{14.0}
\expandafter\def\csname GRCzechoslovakiax1944x1950\endcsname{10.7}
\expandafter\def\csname GRHungaryx1944x1950\endcsname{2.0}
\expandafter\def\csname GRAustriax1944x1950\endcsname{$-$1.9}
\expandafter\def\csname GRIrelandx1944x1950\endcsname{2.1}
\expandafter\def\csname GRGreecex1944x1950\endcsname{9.4}
\expandafter\def\csname GRLithuaniax1944x1950\endcsname{13.7}
\expandafter\def\csname GRLatviax1944x1950\endcsname{8.9}
\expandafter\def\csname eeGR1950x1973\endcsname{3.5}
\expandafter\def\csname eeCUM1950x1973\endcsname{120}
\expandafter\def\csname GRFinlandx1950x1973\endcsname{4.3}
\expandafter\def\csname GRSwedenx1950x1973\endcsname{3.1}
\expandafter\def\csname GRDenmarkx1950x1973\endcsname{3.1}
\expandafter\def\csname GRNordicsx1950x1973\endcsname{3.3}
\expandafter\def\csname GRWesternx1950x1973\endcsname{3.7}
\expandafter\def\csname GRPeripheryx1950x1973\endcsname{5.1}
\expandafter\def\csname GRUSSRx1950x1973\endcsname{3.3}
\expandafter\def\csname GRRussiax1950x1973\endcsname{5.1}
\expandafter\def\csname GRECEx1950x1973\endcsname{3.9}
\expandafter\def\csname GRPolandx1950x1973\endcsname{3.5}
\expandafter\def\csname GRCzechoslovakiax1950x1973\endcsname{3.1}
\expandafter\def\csname GRHungaryx1950x1973\endcsname{3.6}
\expandafter\def\csname GRAustriax1950x1973\endcsname{4.9}
\expandafter\def\csname GRIrelandx1950x1973\endcsname{3.0}
\expandafter\def\csname GRGreecex1950x1973\endcsname{6.2}
\expandafter\def\csname GRLithuaniax1950x1973\endcsname{4.6}
\expandafter\def\csname GRLatviax1950x1973\endcsname{3.7}
\expandafter\def\csname eeGR1950x1978\endcsname{3.0}
\expandafter\def\csname eeCUM1950x1978\endcsname{131}
\expandafter\def\csname GRFinlandx1950x1978\endcsname{3.6}
\expandafter\def\csname GRSwedenx1950x1978\endcsname{2.7}
\expandafter\def\csname GRDenmarkx1950x1978\endcsname{2.7}
\expandafter\def\csname GRNordicsx1950x1978\endcsname{3.0}
\expandafter\def\csname GRWesternx1950x1978\endcsname{3.3}
\expandafter\def\csname GRPeripheryx1950x1978\endcsname{4.6}
\expandafter\def\csname GRUSSRx1950x1978\endcsname{3.0}
\expandafter\def\csname GRRussiax1950x1978\endcsname{4.7}
\expandafter\def\csname GRECEx1950x1978\endcsname{3.7}
\expandafter\def\csname GRPolandx1950x1978\endcsname{3.3}
\expandafter\def\csname GRCzechoslovakiax1950x1978\endcsname{2.9}
\expandafter\def\csname GRHungaryx1950x1978\endcsname{3.4}
\expandafter\def\csname GRAustriax1950x1978\endcsname{4.5}
\expandafter\def\csname GRIrelandx1950x1978\endcsname{3.2}
\expandafter\def\csname GRGreecex1950x1978\endcsname{5.6}
\expandafter\def\csname GRLithuaniax1950x1978\endcsname{4.2}
\expandafter\def\csname GRLatviax1950x1978\endcsname{3.1}
\expandafter\def\csname eeGR1973x1989\endcsname{0.6}
\expandafter\def\csname eeCUM1973x1989\endcsname{11}
\expandafter\def\csname GRFinlandx1973x1989\endcsname{2.7}
\expandafter\def\csname GRSwedenx1973x1989\endcsname{1.6}
\expandafter\def\csname GRDenmarkx1973x1989\endcsname{1.7}
\expandafter\def\csname GRNordicsx1973x1989\endcsname{2.1}
\expandafter\def\csname GRWesternx1973x1989\endcsname{1.9}
\expandafter\def\csname GRPeripheryx1973x1989\endcsname{2.6}
\expandafter\def\csname GRUSSRx1973x1989\endcsname{1.0}
\expandafter\def\csname GRRussiax1973x1989\endcsname{1.2}
\expandafter\def\csname GRECEx1973x1989\endcsname{1.0}
\expandafter\def\csname GRPolandx1973x1989\endcsname{0.4}
\expandafter\def\csname GRCzechoslovakiax1973x1989\endcsname{1.4}
\expandafter\def\csname GRHungaryx1973x1989\endcsname{1.3}
\expandafter\def\csname GRAustriax1973x1989\endcsname{2.4}
\expandafter\def\csname GRIrelandx1973x1989\endcsname{2.9}
\expandafter\def\csname GRGreecex1973x1989\endcsname{1.8}
\expandafter\def\csname GRLithuaniax1973x1989\endcsname{1.2}
\expandafter\def\csname GRLatviax1973x1989\endcsname{0.2}
\expandafter\def\csname eeGR1938x1968\endcsname{3.4}
\expandafter\def\csname eeCUM1938x1968\endcsname{174}
\expandafter\def\csname GRFinlandx1938x1968\endcsname{2.7}
\expandafter\def\csname GRSwedenx1938x1968\endcsname{2.9}
\expandafter\def\csname GRDenmarkx1938x1968\endcsname{2.4}
\expandafter\def\csname GRNordicsx1938x1968\endcsname{2.7}
\expandafter\def\csname GRWesternx1938x1968\endcsname{2.4}
\expandafter\def\csname GRPeripheryx1938x1968\endcsname{3.0}
\expandafter\def\csname GRUSSRx1938x1968\endcsname{3.0}
\expandafter\def\csname GRRussiax1938x1968\endcsname{4.4}
\expandafter\def\csname GRECEx1938x1968\endcsname{2.6}
\expandafter\def\csname GRPolandx1938x1968\endcsname{2.3}
\expandafter\def\csname GRCzechoslovakiax1938x1968\endcsname{2.7}
\expandafter\def\csname GRHungaryx1938x1968\endcsname{2.1}
\expandafter\def\csname GRAustriax1938x1968\endcsname{3.0}
\expandafter\def\csname GRIrelandx1938x1968\endcsname{2.1}
\expandafter\def\csname GRGreecex1938x1968\endcsname{2.3}
\expandafter\def\csname GRLithuaniax1938x1968\endcsname{3.3}
\expandafter\def\csname GRLatviax1938x1968\endcsname{3.0}
\expandafter\def\csname eeGR1987x1989\endcsname{0.7}
\expandafter\def\csname eeCUM1987x1989\endcsname{1}
\expandafter\def\csname GRFinlandx1987x1989\endcsname{5.0}
\expandafter\def\csname GRSwedenx1987x1989\endcsname{1.7}
\expandafter\def\csname GRDenmarkx1987x1989\endcsname{0.7}
\expandafter\def\csname GRNordicsx1987x1989\endcsname{1.7}
\expandafter\def\csname GRWesternx1987x1989\endcsname{3.0}
\expandafter\def\csname GRPeripheryx1987x1989\endcsname{4.4}
\expandafter\def\csname GRUSSRx1987x1989\endcsname{1.1}
\expandafter\def\csname GRRussiax1987x1989\endcsname{0.7}
\expandafter\def\csname GRECEx1987x1989\endcsname{$-$0.3}
\expandafter\def\csname GRPolandx1987x1989\endcsname{0.0}
\expandafter\def\csname GRCzechoslovakiax1987x1989\endcsname{1.4}
\expandafter\def\csname GRHungaryx1987x1989\endcsname{0.7}
\expandafter\def\csname GRAustriax1987x1989\endcsname{3.4}
\expandafter\def\csname GRIrelandx1987x1989\endcsname{5.9}
\expandafter\def\csname GRGreecex1987x1989\endcsname{3.9}
\expandafter\def\csname GRLithuaniax1987x1989\endcsname{0.8}
\expandafter\def\csname GRLatviax1987x1989\endcsname{$-$1.9}
\expandafter\def\csname eeGR1968x1973\endcsname{1.9}
\expandafter\def\csname eeCUM1968x1973\endcsname{10}
\expandafter\def\csname GRFinlandx1968x1973\endcsname{6.5}
\expandafter\def\csname GRSwedenx1968x1973\endcsname{3.1}
\expandafter\def\csname GRDenmarkx1968x1973\endcsname{3.3}
\expandafter\def\csname GRNordicsx1968x1973\endcsname{3.9}
\expandafter\def\csname GRWesternx1968x1973\endcsname{3.9}
\expandafter\def\csname GRPeripheryx1968x1973\endcsname{4.6}
\expandafter\def\csname GRUSSRx1968x1973\endcsname{3.1}
\expandafter\def\csname GRRussiax1968x1973\endcsname{4.6}
\expandafter\def\csname GRECEx1968x1973\endcsname{3.9}
\expandafter\def\csname GRPolandx1968x1973\endcsname{4.3}
\expandafter\def\csname GRCzechoslovakiax1968x1973\endcsname{2.5}
\expandafter\def\csname GRHungaryx1968x1973\endcsname{2.5}
\expandafter\def\csname GRAustriax1968x1973\endcsname{5.4}
\expandafter\def\csname GRIrelandx1968x1973\endcsname{3.5}
\expandafter\def\csname GRGreecex1968x1973\endcsname{7.8}
\expandafter\def\csname GRLithuaniax1968x1973\endcsname{3.7}
\expandafter\def\csname GRLatviax1968x1973\endcsname{2.5}
\expandafter\def\csname eeGR1980x1989\endcsname{0.2}
\expandafter\def\csname eeCUM1980x1989\endcsname{2}
\expandafter\def\csname GRFinlandx1980x1989\endcsname{3.0}
\expandafter\def\csname GRSwedenx1980x1989\endcsname{1.8}
\expandafter\def\csname GRDenmarkx1980x1989\endcsname{2.0}
\expandafter\def\csname GRNordicsx1980x1989\endcsname{2.1}
\expandafter\def\csname GRWesternx1980x1989\endcsname{2.0}
\expandafter\def\csname GRPeripheryx1980x1989\endcsname{2.4}
\expandafter\def\csname GRUSSRx1980x1989\endcsname{1.1}
\expandafter\def\csname GRRussiax1980x1989\endcsname{0.5}
\expandafter\def\csname GRECEx1980x1989\endcsname{0.2}
\expandafter\def\csname GRPolandx1980x1989\endcsname{$-$0.1}
\expandafter\def\csname GRCzechoslovakiax1980x1989\endcsname{1.0}
\expandafter\def\csname GRHungaryx1980x1989\endcsname{1.0}
\expandafter\def\csname GRAustriax1980x1989\endcsname{1.9}
\expandafter\def\csname GRIrelandx1980x1989\endcsname{2.7}
\expandafter\def\csname GRGreecex1980x1989\endcsname{1.3}
\expandafter\def\csname GRLithuaniax1980x1989\endcsname{1.4}
\expandafter\def\csname GRLatviax1980x1989\endcsname{$-$0.1}
\expandafter\def\csname eePEAKCrops\endcsname{1968}
\expandafter\def\csname eeREGAINCrops\endcsname{1947}
\expandafter\def\csname eeSGRCropsx1919x1929\endcsname{2.6}
\expandafter\def\csname eeSGRCropsx1929x1938\endcsname{2.8}
\expandafter\def\csname eeSGRCropsx1938x1950\endcsname{1.1}
\expandafter\def\csname eeSGRCropsx1950x1968\endcsname{2.2}
\expandafter\def\csname eeSGRCropsx1968x1989\endcsname{$-$1.1}
\expandafter\def\csname eeSGRCropsx1989x1991\endcsname{$-$8.7}
\expandafter\def\csname eePEAKLivestock\endcsname{1987}
\expandafter\def\csname eeREGAINLivestock\endcsname{1960}
\expandafter\def\csname eeSGRLivestockx1919x1929\endcsname{7.7}
\expandafter\def\csname eeSGRLivestockx1929x1938\endcsname{1.5}
\expandafter\def\csname eeSGRLivestockx1938x1950\endcsname{$-$3.9}
\expandafter\def\csname eeSGRLivestockx1950x1968\endcsname{3.7}
\expandafter\def\csname eeSGRLivestockx1968x1989\endcsname{2.0}
\expandafter\def\csname eeSGRLivestockx1989x1991\endcsname{$-$8.6}
\expandafter\def\csname eePEAKForestryfishing\endcsname{1976}
\expandafter\def\csname eeREGAINForestryfishing\endcsname{1947}
\expandafter\def\csname eeSGRForestryfishingx1919x1929\endcsname{4.7}
\expandafter\def\csname eeSGRForestryfishingx1929x1938\endcsname{0.6}
\expandafter\def\csname eeSGRForestryfishingx1938x1950\endcsname{0.9}
\expandafter\def\csname eeSGRForestryfishingx1950x1968\endcsname{7.6}
\expandafter\def\csname eeSGRForestryfishingx1968x1989\endcsname{1.4}
\expandafter\def\csname eeSGRForestryfishingx1989x1991\endcsname{$-$12.4}
\expandafter\def\csname eePEAKPrimarysectortotal\endcsname{1986}
\expandafter\def\csname eeREGAINPrimarysectortotal\endcsname{1959}
\expandafter\def\csname eeSGRPrimarysectortotalx1919x1929\endcsname{5.2}
\expandafter\def\csname eeSGRPrimarysectortotalx1929x1938\endcsname{1.9}
\expandafter\def\csname eeSGRPrimarysectortotalx1938x1950\endcsname{$-$1.3}
\expandafter\def\csname eeSGRPrimarysectortotalx1950x1968\endcsname{3.1}
\expandafter\def\csname eeSGRPrimarysectortotalx1968x1989\endcsname{0.6}
\expandafter\def\csname eeSGRPrimarysectortotalx1989x1991\endcsname{$-$8.9}
\expandafter\def\csname eePEAKFood\endcsname{1976}
\expandafter\def\csname eeREGAINFood\endcsname{1948}
\expandafter\def\csname eeSGRFoodx1919x1929\endcsname{---}
\expandafter\def\csname eeSGRFoodx1929x1938\endcsname{---}
\expandafter\def\csname eeSGRFoodx1938x1950\endcsname{2.2}
\expandafter\def\csname eeSGRFoodx1950x1968\endcsname{8.8}
\expandafter\def\csname eeSGRFoodx1968x1989\endcsname{2.1}
\expandafter\def\csname eeSGRFoodx1989x1991\endcsname{$-$9.7}
\expandafter\def\csname eePEAKTextiles\endcsname{1987}
\expandafter\def\csname eeREGAINTextiles\endcsname{1949}
\expandafter\def\csname eeSGRTextilesx1919x1929\endcsname{---}
\expandafter\def\csname eeSGRTextilesx1929x1938\endcsname{---}
\expandafter\def\csname eeSGRTextilesx1938x1950\endcsname{4.6}
\expandafter\def\csname eeSGRTextilesx1950x1968\endcsname{9.0}
\expandafter\def\csname eeSGRTextilesx1968x1989\endcsname{0.9}
\expandafter\def\csname eeSGRTextilesx1989x1991\endcsname{$-$5.5}
\expandafter\def\csname eePEAKOtherlightindustries\endcsname{1988}
\expandafter\def\csname eeREGAINOtherlightindustries\endcsname{1948}
\expandafter\def\csname eeSGROtherlightindustriesx1919x1929\endcsname{---}
\expandafter\def\csname eeSGROtherlightindustriesx1929x1938\endcsname{---}
\expandafter\def\csname eeSGROtherlightindustriesx1938x1950\endcsname{2.8}
\expandafter\def\csname eeSGROtherlightindustriesx1950x1968\endcsname{7.1}
\expandafter\def\csname eeSGROtherlightindustriesx1968x1989\endcsname{1.8}
\expandafter\def\csname eeSGROtherlightindustriesx1989x1991\endcsname{$-$12.4}
\expandafter\def\csname eePEAKHeavyindustries\endcsname{1980}
\expandafter\def\csname eeREGAINHeavyindustries\endcsname{1948}
\expandafter\def\csname eeSGRHeavyindustriesx1919x1929\endcsname{---}
\expandafter\def\csname eeSGRHeavyindustriesx1929x1938\endcsname{---}
\expandafter\def\csname eeSGRHeavyindustriesx1938x1950\endcsname{5.1}
\expandafter\def\csname eeSGRHeavyindustriesx1950x1968\endcsname{12.0}
\expandafter\def\csname eeSGRHeavyindustriesx1968x1989\endcsname{1.2}
\expandafter\def\csname eeSGRHeavyindustriesx1989x1991\endcsname{$-$8.6}
\expandafter\def\csname eePEAKIndustrytotal\endcsname{1988}
\expandafter\def\csname eeREGAINIndustrytotal\endcsname{1948}
\expandafter\def\csname eeSGRIndustrytotalx1919x1929\endcsname{11.5}
\expandafter\def\csname eeSGRIndustrytotalx1929x1938\endcsname{5.7}
\expandafter\def\csname eeSGRIndustrytotalx1938x1950\endcsname{5.0}
\expandafter\def\csname eeSGRIndustrytotalx1950x1968\endcsname{9.5}
\expandafter\def\csname eeSGRIndustrytotalx1968x1989\endcsname{1.4}
\expandafter\def\csname eeSGRIndustrytotalx1989x1991\endcsname{$-$9.1}
\expandafter\def\csname eePEAKConstruction\endcsname{1977}
\expandafter\def\csname eeREGAINConstruction\endcsname{1947}
\expandafter\def\csname eeSGRConstructionx1919x1929\endcsname{11.7}
\expandafter\def\csname eeSGRConstructionx1929x1938\endcsname{7.6}
\expandafter\def\csname eeSGRConstructionx1938x1950\endcsname{5.4}
\expandafter\def\csname eeSGRConstructionx1950x1968\endcsname{9.3}
\expandafter\def\csname eeSGRConstructionx1968x1989\endcsname{0.9}
\expandafter\def\csname eeSGRConstructionx1989x1991\endcsname{$-$9.9}
\expandafter\def\csname eePEAKTransportcommunications\endcsname{1989}
\expandafter\def\csname eeREGAINTransportcommunications\endcsname{1945}
\expandafter\def\csname eeSGRTransportcommunicationsx1919x1929\endcsname{10.7}
\expandafter\def\csname eeSGRTransportcommunicationsx1929x1938\endcsname{2.0}
\expandafter\def\csname eeSGRTransportcommunicationsx1938x1950\endcsname{8.0}
\expandafter\def\csname eeSGRTransportcommunicationsx1950x1968\endcsname{8.9}
\expandafter\def\csname eeSGRTransportcommunicationsx1968x1989\endcsname{4.9}
\expandafter\def\csname eeSGRTransportcommunicationsx1989x1991\endcsname{$-$4.2}
\expandafter\def\csname eePEAKTradeaccommodationfoodservice\endcsname{1990}
\expandafter\def\csname eeREGAINTradeaccommodationfoodservice\endcsname{1949}
\expandafter\def\csname eeSGRTradeaccommodationfoodservicex1919x1929\endcsname{---}
\expandafter\def\csname eeSGRTradeaccommodationfoodservicex1929x1938\endcsname{---}
\expandafter\def\csname eeSGRTradeaccommodationfoodservicex1938x1950\endcsname{1.2}
\expandafter\def\csname eeSGRTradeaccommodationfoodservicex1950x1968\endcsname{9.6}
\expandafter\def\csname eeSGRTradeaccommodationfoodservicex1968x1989\endcsname{4.7}
\expandafter\def\csname eeSGRTradeaccommodationfoodservicex1989x1991\endcsname{$-$22.0}
\expandafter\def\csname eePEAKHousing\endcsname{1986}
\expandafter\def\csname eeREGAINHousing\endcsname{1949}
\expandafter\def\csname eeSGRHousingx1919x1929\endcsname{1.8}
\expandafter\def\csname eeSGRHousingx1929x1938\endcsname{6.8}
\expandafter\def\csname eeSGRHousingx1938x1950\endcsname{1.1}
\expandafter\def\csname eeSGRHousingx1950x1968\endcsname{2.8}
\expandafter\def\csname eeSGRHousingx1968x1989\endcsname{2.0}
\expandafter\def\csname eeSGRHousingx1989x1991\endcsname{$-$4.6}
\expandafter\def\csname eePEAKPublicservicestotal\endcsname{1984}
\expandafter\def\csname eeREGAINPublicservicestotal\endcsname{1945}
\expandafter\def\csname eeSGRPublicservicestotalx1919x1929\endcsname{$-$5.3}
\expandafter\def\csname eeSGRPublicservicestotalx1929x1938\endcsname{1.3}
\expandafter\def\csname eeSGRPublicservicestotalx1938x1950\endcsname{4.0}
\expandafter\def\csname eeSGRPublicservicestotalx1950x1968\endcsname{1.2}
\expandafter\def\csname eeSGRPublicservicestotalx1968x1989\endcsname{1.8}
\expandafter\def\csname eeSGRPublicservicestotalx1989x1991\endcsname{$-$2.5}
\expandafter\def\csname eePEAKTotalGVAGDP\endcsname{1989}
\expandafter\def\csname eeREGAINTotalGVAGDP\endcsname{1947}
\expandafter\def\csname eeSGRTotalGVAGDPx1919x1929\endcsname{3.3}
\expandafter\def\csname eeSGRTotalGVAGDPx1929x1938\endcsname{2.9}
\expandafter\def\csname eeSGRTotalGVAGDPx1938x1950\endcsname{1.9}
\expandafter\def\csname eeSGRTotalGVAGDPx1950x1968\endcsname{5.4}
\expandafter\def\csname eeSGRTotalGVAGDPx1968x1989\endcsname{1.7}
\expandafter\def\csname eeSGRTotalGVAGDPx1989x1991\endcsname{$-$8.7}
\newcommand{\speak}[1]{\csname eePEAK#1\endcsname}
\newcommand{\sregain}[1]{\csname eeREGAIN#1\endcsname}
\newcommand{\sgr}[3]{\csname eeSGR#1x#2x#3\endcsname}
\expandafter\def\csname eeTGR1919x1929\endcsname{3.3}
\expandafter\def\csname eeTGR1929x1938\endcsname{2.9}
\expandafter\def\csname eeTGR1938x1950\endcsname{1.9}
\expandafter\def\csname eeTGR1950x1968\endcsname{5.4}
\expandafter\def\csname eeTGR1968x1989\endcsname{1.7}
\expandafter\def\csname eeTGR1989x1991\endcsname{$-$8.7}
\expandafter\def\csname eeTGR1950x1978\endcsname{4.3}
\expandafter\def\csname eeTGR1950x1973\endcsname{4.9}
\expandafter\def\csname eeTGR1973x1989\endcsname{1.3}
\expandafter\def\csname eeTGR1938x1989\endcsname{3.1}
\expandafter\def\csname eeTGR1919x1938\endcsname{3.1}
\newcommand{\tgr}[2]{\csname eeTGR#1x#2\endcsname}
\expandafter\def\csname eeREGAINnineteen\endcsname{2001}
\expandafter\def\csname eeTROUGHyear\endcsname{1993}
\expandafter\def\csname eeTROUGHval\endcsname{5,870}
\expandafter\def\csname eeTROUGHpct\endcsname{34}


% ---- Section formatting (EHR style: "I | UPPERCASE TITLE") ----
\usepackage{titlesec}
\renewcommand{\thesection}{\Roman{section}}   % cross-references print "Section IV"
\titleformat{\section}{\normalfont\large\bfseries}{\Roman{section}~|~}{0em}{\MakeUppercase}
\titleformat{\subsection}{\normalfont\normalsize\bfseries\itshape}{}{0em}{}

% ---- Line numbers for review (comment out for submission) ----
% \usepackage{lineno}
% \linenumbers

% =============================================================
\begin{document}
% =============================================================

% ---- Title page ----
\begin{titlepage}
    \centering
    \vspace*{2cm}
    {\large\bfseries New GDP series for interwar and Soviet Estonia:\\
    a Finnish miracle that never was\par}
    \vspace{1.5cm}
    {\normalsize
    Adomas Klimantas\\
    ISM University of Management and Economics\\
    \texttt{a.klimantas@ism.lt}\\[0.8em]
    Martin Klesment\\
    University of Tartu\\
    \texttt{martin.klesment@ut.ee}
    \par}
    \vspace{1.5cm}
    {\small\today\par}
    \vspace{1.5cm}
    \begin{abstract}
        \noindent
        This paper constructs the first internationally comparable annual series of Estonian GDP per capita for 1919--95 and splices it to existing estimates to give an unbroken series for 1913 and 1919--2022 in 1990 and 2011 international dollars, ready for the Maddison Project Database. The reconstruction proceeds in three steps: a new 1938 benchmark based on a bilateral price comparison with Sweden; a recalculated and extended production-side series for 1919--40; and a new series for the Estonian SSR for 1937--95, built from physical output and employment series drawn from Soviet yearbooks and the Estonian archives and anchored to Bolotin's 1990 benchmark. Independently constructed series agree within 4 per cent where they overlap in 1937--40, and within 1 per cent at the 1991 splice. Estonia grew by \gr{1919}{1929} per cent a year in the 1920s and \gr{1929}{1938} per cent in the 1930s, by \gr{1950}{1968} per cent a year in 1950--68 and by \gr{1968}{1989} per cent in 1968--89; the transition cut income per head by a third. Estonia stood at \pctfin{1919} per cent of Finnish income per head in 1919 but at only \pctfin{1938} per cent in 1938: almost half of the gap that separated the two countries in 1989 had opened before Soviet occupation. Finland is a revealing comparator for Estonia, but not a measure of what Estonia lost.
    \end{abstract}
    \vspace{0.8cm}
    {\small\textbf{JEL codes:} N14, O47, E01, P27\par}
    {\small\textbf{Keywords:} GDP, historical national accounts, Estonia, Finland, Soviet economy,
    Maddison Project Database\par}
\end{titlepage}

\newpage

% =============================================================
% INTRODUCTION (unnumbered, EHR style)
% =============================================================
\section*{Introduction}

For Estonians, Finland has long served as an `alter ego': the neighbour across the Gulf that shared their language, their Lutheran faith and their nineteenth-century awakening, and that showed, through the television signals picked up in Tallinn from the 1960s, what Estonia might have become had it not spent half a century inside the Soviet Union.\footnote{\autocite[83, 86]{lepppantti2013}.} The comparison is more than nostalgia. It underpins the belief, widespread in the Baltic states, that living standards before the Second World War were comparable to those of the Finns, so that the whole of the gap between the two countries at the restoration of independence in 1991 can be charged to Soviet rule.\footnote{\autocite[154]{norkus_restorations}; \autocite[206]{dellenbrant1992}; \autocite[81--2]{bocharnikova2021}.} In its strongest form, found in the Estonian government's `White Book' on the losses inflicted by occupation, the argument becomes a counterfactual accounting exercise: had Estonia developed as Finland did, its GDP per capita in the early 2000s would have been four to five times what it actually was.\footnote{\autocite[150]{kukk2005}; \autocite[260]{klesment2009}.} Others have doubted that the two economies were ever close, and have pointed out that Finland's own post-war growth was exceptional even by Western European standards, so that the divergence tells us as much about Finland as about Estonia.\footnote{\autocite[18]{norkus2016}; \autocite[201]{koll1994gap}.}

What has prevented these claims from being tested is the absence of an annual, internationally comparable GDP series for Estonia before 1995. The existing estimates are partial and rest on conversions that are no longer defensible. Valge's pioneering production-side series covers 1923--38 and was converted into international dollars by assuming, on the authority of Colin Clark's rough figures for 1925--34, that Estonian income per head stood at 90 per cent of Finland's.\footnote{\autocite[2725--6]{valge2003all}; \autocite[60]{clark1938}.} Klesment, Puur and Valge produced a series for 1950--90 from four sectors of material production, scaled to a 1990 level obtained by deflating a purchasing-power-parity estimate for 2005, of undocumented derivation, published by the UN Economic Commission for Europe (UNECE).\footnote{\autocite[39--44]{klesment2010}.} The Maddison Project Database (MPD) fills the gap between 1973 and 1990 by backcasting from modern benchmarks with Soviet net material product (NMP) growth rates, a procedure that Maddison himself regarded as inadmissible for the Soviet Union as a whole, and that places Estonian income in 1989 almost a third above the figure reported here.\footnote{\autocite[307--8]{Maddison1998b}; \autocite{mpd2020}, as quoted in \autocite[155]{norkus_restorations}. See Table~\ref{tab:previous}.} The only estimates that are methodologically consistent with the Maddison tradition are the indirect approximations of Norkus and Markevi\v{c}i\={u}t\.{e} for four benchmark years, 1913, 1922, 1929 and 1938.\footnote{\autocite{norkusmark2021}.} Between those years, and for the whole of the Soviet period, there is nothing to which a Finnish, Swedish or Soviet series can be compared.

This paper fills that gap. It makes three contributions. First, it constructs a new benchmark estimate of Estonian GDP per capita for 1938, replacing Valge's Finnish bridge with a direct bilateral comparison of consumer prices with Sweden of the kind already used for Lithuania in 1937 and Latvia in 1935, and it uses the new benchmark to recalculate and extend the production-side series to every year from 1919 to 1940.\footnote{\autocite{klimantas2019}; \autocite{norkuslatvia1935}; \autocite{klimantas2023lithuania}; \autocite{Klimantas2024}.} Second, and principally, it provides the first annual series for the Estonian Soviet Socialist Republic (ESSR) covering 1937--95, built from 101 physical product series, agricultural output, transport volumes and employment records drawn from Soviet statistical yearbooks and the Estonian National Archives, and anchored to Bolotin's purchasing-power-parity benchmark for 1990. Third, it splices the two new series to each other, to the 1913 benchmark of Norkus and Markevi\v{c}i\={u}t\.{e}, and to adjusted post-1990 MPD data, to give a continuous series for 1913 and 1919--2022 in both 1990 and 2011 international dollars, together with a consistent population series. Because the interwar and Soviet-period series are built from different sources and anchored to benchmarks fifty-two years apart, their agreement where they overlap, within 4 per cent in each of the years 1937--40 and within 1 per cent at the 1991 splice, is the strongest available check on the whole reconstruction.

The new series alters the received picture in three respects. First, Estonia's interwar growth was real but unremarkable: income per head rose by \gr{1919}{1929} per cent a year in the 1920s and \gr{1929}{1938} per cent a year in the 1930s, against \grx{Finland}{1919}{1929} and \grx{Finland}{1929}{1938} per cent in Finland. Estonia emerged from the war of independence in 1919 at \pctfin{1919} per cent of Finnish income per head, close to the near-parity of 1913, but stood at only \pctfin{1938} per cent in 1938. Measured in logarithms, almost half of the gap between the two countries in 1989 had opened before a Soviet soldier crossed the border. Second, the Soviet period divides sharply in two. Between 1950 and 1968 Estonian income per head grew by \gr{1950}{1968} per cent a year, the fastest sustained growth in the country's history before 1995 and faster than in Finland or the Western European core (hereafter `Western core'), though not than the Western periphery (Table~\ref{tab:growth}); between 1968 and 1989 it grew by \gr{1968}{1989} per cent a year, and the ground gained was lost again. Estonia stood at \pctwc{1968} per cent of the Western core in 1968 and at \pctwc{1989} per cent in 1989, almost exactly where it had been in 1913. Within the Soviet Union, by contrast, Estonia's position barely changed, from \pctussr{1940} per cent of the Union average in 1940 to \pctussr{1989} per cent in 1989, one of the richest republics throughout. Third, the transition shock was severe: income per head fell by \eeTROUGHpct{} per cent between 1989 and the trough of \eeTROUGHyear{}, and did not regain its 1989 level until 2001.

The paper also speaks to a wider literature. The economic history of the Soviet Union has been rewritten over the past two decades from archival sources, but that work has concentrated on Russia; the non-Russian republics, which held half the Union's population, remain almost entirely without modern national accounts.\footnote{\autocite{zhuravskaya2024}; \autocite[277--302]{vonyoinmorys2021}; \autocite{voskoboynikov2021}; \autocite{kalinovsky2022}.} The Estonian series is built with the same methods as the recent series for Lithuania, so that the two can be compared directly and the procedure repeated for the remaining republics,\footnote{\autocite{klimantas2026soviet}.} and it adds a small, open economy with a strong pre-war statistical tradition to the European peripheral economies for which annual interwar accounts now exist.\footnote{\autocite[877--80]{kenny_irelandgdp}; \autocite[1--24]{broadberryinmorys}.}

Section~\ref{sec:background} sets the scene and reviews the existing estimates; Section~\ref{sec:method} describes sources and method; Section~\ref{sec:results} presents the series; Section~\ref{sec:limitations} validates it and assesses its biases; Section~\ref{sec:implications} reinterprets Estonian growth under capitalism, planning and transition; Section~\ref{sec:finland} returns to Finland; Section~\ref{sec:conclusion} concludes.

% =============================================================
% I. HISTORICAL BACKGROUND AND PREVIOUS ESTIMATES
% =============================================================
\section{Historical background and previous estimates}
\label{sec:background}

The territory of present-day Estonia passed from Sweden to Russia in 1710--21 and remained within the Russian Empire until 1918. Two inheritances from the Swedish and Baltic-German centuries distinguished it from the rest of the empire. Lutheran parish schooling made the Estonian and Latvian governorates the most literate in the Tsar's dominions, and the early emancipation of the serfs in 1816--19, followed from the 1860s by the sale of farmsteads to Estonian peasants, created a class of proprietors with an orientation towards market agriculture well before 1914.\footnote{\autocite[4--5]{baltichistorykasekamp2017}; \autocite[149--66]{eellend2010}.} By the eve of the First World War Tallinn and Narva had become significant manufacturing centres serving the imperial market, so that Estonia was industrialised out of proportion to its size.\footnote{\autocite[102]{ehbaltics_tarvel}.} The best available estimate places Estonian income per head in 1913 at \ee{1913}~\intlnineteen{}, within a few per cent of Finland's and about half the level of the Western core (Table~\ref{tab:relative}).\footnote{\autocite[603]{norkusmark2021}.}

The war, the German occupation of 1918 and the war of independence of 1918--20 reversed these gains. Factories were evacuated eastwards, the urban population fell sharply and the Russian market, which had absorbed the bulk of pre-war Baltic manufactures, was lost for good.\footnote{\autocite[107]{ehbaltics_tarvel}; \autocite[193--210]{koll1994gap}.} The new republic was therefore built on the land. The land reform of 1919, among the most radical in Europe, broke up the estates of the Baltic-German nobility and created more than 50,000 new smallholdings; co-operative dairies and export associations reoriented agriculture towards butter, bacon and eggs for the British and German markets, in conscious imitation of Denmark.\footnote{\autocite[64--97]{jorgensen2006}; \autocite{norkus_denmarkisation}.} Agriculture still employed some 63 per cent of the labour force in 1934 (Table~\ref{tab:structure}).\footnote{\autocite[181]{morysinterwarinmorys2021}.} The Depression brought a banking crisis in 1931--2 and a devaluation in 1933 that pegged the kroon to sterling, and the authoritarian government installed in 1934 presided over a state-led drive into import-substituting industry now described as `economic nationalism'.\footnote{\autocite[74]{bernankejames01}; \autocite[11--29]{koll1998}; \autocite[118--35]{klesment2005}; \autocite[157--78]{feest2025}.} Output was still expanding when the Molotov--Ribbentrop pact assigned Estonia to the Soviet sphere in August 1939 and Soviet forces occupied the country in June 1940.

Soviet rule abolished the interwar economic order within a year. Banks, industry and urban property were nationalised in 1940--1; the German occupation of 1941--4 and the renewed fighting of 1944 followed; and in March 1949 the deportation of some 20,000 people broke the resistance of the countryside to collectivisation, which was essentially complete within months.\footnote{\autocite[265]{taagepera_coll01}; \autocite[12]{deportations_ee}; \autocite[100]{feest2014}.} Industry was rebuilt as a specialised component of the all-Union economy: oil shale and electricity in the north-east, cotton textiles in Narva and Tallinn, and later electrical engineering, with the labour supplied by immigration from Russia on a scale that reduced Estonians to little more than three-fifths of the population by 1989.\footnote{\autocite[158--217]{case_estonia}; \autocite[47]{ahderajasalu1992}.} Estonia was consistently among the two or three richest republics of the Union on every Soviet measure of income and consumption.\footnote{\autocite[169]{schroeder1973}; \autocite[35--6]{bs_livingstandards1988}.} Independence was restored in August 1991; a currency board followed in June 1992, prices were freed and enterprises privatised, and the collapse of eastern markets produced the deepest contraction in the country's recorded history before growth resumed in 1995.\footnote{\autocite[1--8]{wb1993estonia}; \autocite[443--4]{transition_progress}.}

The quantitative record of this history is thin. Estonia published no national income accounts before 1940, and the Soviet statistical system measured net material product, which excludes most services and values output at administered prices. Three long series have been attempted, besides the benchmark approximations of Norkus and Markevi\v{c}i\={u}t\.{e}, and their results are compared with the new series in Table~\ref{tab:previous}.

% tab:previous
\begin{table}[htbp]
\centering\small
\caption[The new series compared with previous estimates of Estonian GDP per capita]{\textsc{The new series compared with previous estimates of Estonian GDP per capita}}
\label{tab:previous}
\begin{threeparttable}
\begin{tabular}{lrrrrrrr}
\toprule
& This paper & \multicolumn{2}{c}{Valge (2003)} & \multicolumn{2}{c}{Klesment et al. (2010)} & \multicolumn{2}{c}{MPD (2020)} \\
\cmidrule(lr){3-4}\cmidrule(lr){5-6}\cmidrule(lr){7-8}
Year & (\intlnineteen{}) & Level & Diff.\,\% & Level & Diff.\,\% & Level & Diff.\,\% \\
\midrule
1923 & 1,913 & 1,811 & $-$5.3 & --- & --- & --- & --- \\
1924 & 1,895 & 2,337 & +23.4 & --- & --- & --- & --- \\
1929 & 2,113 & 2,182 & +3.3 & --- & --- & --- & --- \\
1932 & 2,105 & 2,110 & +0.2 & --- & --- & --- & --- \\
1938 & 2,685 & 2,745 & +2.2 & --- & --- & --- & --- \\
1950 & 3,667 & --- & --- & 1,676 & $-$54.3 & --- & --- \\
1960 & 5,319 & --- & --- & 3,949 & $-$25.8 & --- & --- \\
1970 & 7,688 & --- & --- & 5,377 & $-$30.1 & --- & --- \\
1973 & 8,080 & --- & --- & 5,621 & $-$30.4 & 8,657 & +7.1 \\
1980 & 8,785 & --- & --- & 5,834 & $-$33.6 & --- & --- \\
1989 & 8,958 & --- & --- & 6,692 & $-$25.3 & 11,780 & +31.5 \\
1990 & 8,425 & --- & --- & 6,444 & $-$23.5 & --- & --- \\
1995 & 6,277 & --- & --- & 5,084 & $-$19.0 & --- & --- \\
\bottomrule
\end{tabular}
\begin{tablenotes}[para,flushleft]
\footnotesize
\textit{Sources:} this paper: see online appendix. Valge: \cite[2726]{valge2003all}, Table 50, Geary--Khamis dollars obtained by scaling Finnish income with Colin Clark's 1925--34 ratio. Klesment et al.: \cite[44]{klesment2010}, Table 2.2, UNECE 2005 purchasing-power parities deflated to 1990 US dollars with the consumer price index, four sectors only. MPD: \cite{mpd2020}, as quoted in \cite[155]{norkus_restorations}, converted from \intltwenty{} at 1.594. \textit{Notes:} Diff. = previous estimate relative to this paper, per cent. `---': not available.
\end{tablenotes}
\end{threeparttable}
\end{table}


The foundation for the interwar period is Valge's production-side estimate of GDP at current and constant prices for 1923--38.\footnote{\autocite{valge2003all}.} Its coverage of agriculture, industry, construction and housing is retained here, but two features limit its use for international comparison. Valge deflated current-price GDP with the cost-of-living index, so that the rapid fall in food prices after 1923 produces a rise of 29 per cent in real income per head in a single year, 1924, when crop, livestock, industrial and transport output each moved by less than 5 per cent.\footnote{\autocite[2718, 2721]{valge2003all}; Klesment's re-deflation with wholesale price indices reduces the jump to 21 per cent (\autocite[29--31]{klesment2008a}; \autocite[37--8]{klesment2010}).} And he expressed the result in international dollars indirectly, by taking Colin Clark's estimate that Estonian output per head averaged 90 per cent of Finland's in 1925--34 and applying that ratio to Maddison's Finnish figures.\footnote{\autocite[2725]{valge2003all}; \autocite[60]{clark1938}.} The resulting level is not far from the one reported here for the 1930s (Table~\ref{tab:previous}), but only because a domestic-currency GDP that is too low is offset by a conversion ratio that is too high. Norkus and Markevi\v{c}i\={u}t\.{e} have since produced indirect, income-side approximations for 1913, 1922, 1929 and 1938 that are consistent with the MPD; their 1913 figure is used here as the pre-war anchor, and their 1938 figure, about 5 per cent below the new benchmark, is an independent check on it.\footnote{\autocite{norkusmark2021}; \autocite[154]{norkus_restorations}.}

For the Soviet period the only long series is that of Klesment, Puur and Valge, who extrapolated a 1990 level derived from UNECE purchasing-power parities backwards with an official rouble GDP index for 1980--9 and with a composite index of physical output in agriculture, industry, construction and transport for 1950--79.\footnote{\autocite[39--44]{klesment2010}.} The four sectors cover only about two-thirds of GDP, since trade and the non-material sphere are omitted, and the UNECE-based level understates income: even in 1995 the series stands 19 per cent below the estimate obtained here from the MPD's 2011 benchmark. Its growth rate of income per head for 1950--78, 4.4 per cent a year, is well above that of the present series, for reasons discussed in Section~\ref{sec:results}. The MPD itself reports Estonian GDP per capita for 1973 and annually from 1980. Its 1990 level descends from Maddison's application of Bolotin's republic-to-Union ratios, taken for 1991 rather than 1990, to a Soviet total itself derived from a bilateral Austria--USSR comparison; its earlier years are extrapolated backwards with NMP growth rates adjusted by the Union-wide ratio of GDP to NMP, a short cut that the literature on material-product accounting had warned against because the ratio varies across countries and over time, and that overstates growth in the final Soviet decades. The MPD figure for 1989 is 32 per cent above the one reported here.\footnote{\autocite[181--2]{maddisonworldeconomy}; \autocite[307--8, 318]{Maddison1998b}; \autocite[593]{ceska1989}; \autocite{Campbell1985}; \autocite[279--80]{vonyoinmorys2021}; Table~\ref{tab:previous}.}

% =============================================================
% II. SOURCES AND METHOD
% =============================================================
\section{Sources and method}
\label{sec:method}

\subsection{The production approach and the reliability of Soviet statistics}

For Estonia only the output side of national income is feasible in either period: consistent annual data on expenditure or factor incomes do not exist for the interwar republic, whose statistical office never compiled national accounts, and for the Soviet period the monetary aggregates that an expenditure or income estimate would require are precisely the figures that cannot be trusted. The series presented here are therefore constructed, in both periods, by extrapolating a benchmark estimate of GDP with constant-price indices of sectoral value-added, each sector weighted by its share of GDP in the benchmark year:
\[
    \mathit{GDP}_{t} = \mathit{GDP}_{\mathrm{bench}} \times \sum_{i=1}^{n} w_{i}\, I_{it},
\]
where $w_i$ is sector $i$'s share of benchmark GDP and $I_{it}$ its volume index in year $t$. The procedure is standard where expenditure and income data are sparse, and is the one used for Ireland, the European periphery and, in the companion studies, Lithuania and Latvia.\footnote{\autocite[880--3]{kenny_irelandgdp}; \autocite[22]{broadberryinmorys}; \autocite{klimantas2023lithuania}; \autocite{Klimantas2024}.}

For the Soviet period this approach also has an intellectual pedigree that bears directly on the reliability of the result. Soviet practice measured net material product at prevailing prices, excluded most services, double-counted intermediate transactions and, above all, understated inflation, so that official constant-price growth was too high: in Maddison's verdict, the official series `overstated growth, and understated inflation'.\footnote{\autocite[307]{Maddison1998b}.} The Western response, begun by Bergson in the 1950s and institutionalised by the CIA from the 1960s, was to rebuild Soviet GNP from the bottom up, repricing output at factor cost, imputing the missing services and, crucially, replacing the official value indices with indices built from physical output series weighted by value-added.\footnote{\autocite{bergson1950}; \autocite{bergson1961}; \autocite[48]{cia1982a}; \autocite[328]{Maddison1995a}.} The CIA estimates were attacked from both sides. Navrozov and others argued that by accepting Soviet physical data they still overstated growth; Khanin, working from Soviet materials in the 1980s, obtained much lower growth rates for the Stalin period but, as Harrison showed, results for 1950--87 that coincided with the CIA's at 3.8 per cent a year.\footnote{\autocite{assessingCIA1988}; \autocite[144, 159]{harrison1993}.} The CIA's own indices were reworked at 1982 prices for the Joint Economic Committee in 1990, the form in which they entered the later literature and the MPD.\footnote{\autocite{cia1990}.} The assessments by Becker and by Schroeder, the review by Maddison and the recent survey by Vonyó and Markevich all reach the same conclusion: the CIA series contain errors in both directions, most of the criticism `was not well founded', and the physical-output method remains the best available basis for measuring Soviet growth.\footnote{\autocite{becker1994}; \autocite{schroeder1995}; \autocite[309]{Maddison1998b}; \autocite[279--80]{vonyoinmorys2021}.} Ponomarenko's archival reconstruction of Russian national accounts, the fullest attempt so far to recast Soviet data in the System of National Accounts (SNA), reaches the same judgement from inside the Soviet statistical system. The physical output data `are apparently of high quality and can be used', whereas the official value indices were inflated by the pricing of new products; Khanin's alternative, in Ponomarenko's view, rests on a `rather primitive methodological base' compared with the CIA's.\footnote{\autocite[16--18, 73--4]{ponomarenko2002}.} The literature on the Soviet economy has grown rapidly since the opening of the archives, as Zhuravskaya, Guriev and Markevich document, but its quantitative core for the post-war decades still rests on physical output series of this kind.\footnote{\autocite[47--114]{zhuravskaya2024}; \autocite{voskoboynikov2021}.}

It matters, therefore, to be precise about which Soviet data this paper uses and which it avoids. The distortions documented in the literature were concentrated in value aggregates: in the national income and gross output indices at `comparable prices', in the reporting of hidden price increases as real growth, and in the double counting of the material product system.\footnote{\autocite[307]{Maddison1998b}; \autocite{vonyo2017}.} Khanin's classification of the official data, as Harrison summarises it, is that `the ease of distortion remained greater for value series than for physical volumes', and that distortion was concentrated in heterogeneous products, above all in machinery and construction, while physical series for homogeneous goods such as grain, milk, electricity, cement or railway freight are comparatively reliable.\footnote{\autocite[144, 152]{harrison1993}.} The sectoral indices below rest almost entirely on physical quantities of this kind, on employment and on inputs; the one place where value data at comparable prices enter, trade turnover, which also feeds the productivity adjustment for finance, is flagged and its bias bounded in Section~\ref{sec:limitations}. The claim, then, is not that Soviet statistics were honest, but that the particular statistics used here were recorded for operational planning and had little scope for manipulation.

\subsection{The benchmarks}

\textit{The 1938 benchmark.} Valge's estimate of current-price value-added for 1938 in agriculture, forestry and fishing (250.1 million kroon), industry (132.1 million), construction (14.0 million) and housing (29.0 million) is accepted without revision. His treatment of the remaining services is not. For all other services Valge assumed uniform value-added per worker and scaled a single figure by the service labour force, obtaining 122.9 million kroon in all, less than the contribution of industry and implausibly low for an economy at Estonia's level of urbanisation.\footnote{\autocite[2718]{valge2003all}.} This paper rebuilds service value-added from its income components, as in the Lithuanian and Latvian benchmark studies: compensation of employees is estimated from the labour force and average earnings of each service branch, and the remaining components, operating surplus, depreciation and net production taxes, are added in the proportions observed in the Latvian benchmark of 1935, the closest comparator with a full set of accounts.\footnote{\autocite{klimantas2019}; \autocite{norkuslatvia1935}.} Commerce yields 91.5 million kroon, transport and communications 32.2 million, public services 99.3 million, of which 21.4 million is defence, and household services 5.5 million: 228.5 million in all. Gross value-added rises to 654 million kroon and GDP at market prices, with 17.2 million of taxes less subsidies, to 671 million, or about 592 kroon per head against Valge's 502. The sectoral composition is reported in Table~\ref{tab:benchmarks}.

% tab:benchmarks
\begin{table}[htbp]
\centering\small
\caption[Benchmark GDP by sector: Estonia 1938 and 1990, current prices]{\textsc{Benchmark GDP by sector: Estonia 1938 and 1990, current prices}}
\label{tab:benchmarks}
\begin{threeparttable}
\begin{tabular}{lrrrr}
\toprule
& \multicolumn{2}{c}{1938} & \multicolumn{2}{c}{1990} \\
\cmidrule(lr){2-3}\cmidrule(lr){4-5}
Sector & m.\,\eek{} & \% & m.\,roubles & \% \\
\midrule
Agriculture, forestry \& fishing & 250.1 & 38.3 & 1,250.0 & 16.9 \\
Industry & 132.1 & 20.2 & 2,203.4 & 29.8 \\
Construction & 14.0 & 2.1 & 818.2 & 11.1 \\
Transport \& communications & 32.2 & 4.9 & 535.7 & 7.3 \\
Trade, finance, housing \& other commercial services & 120.5 & 18.4 & 1,322.3 & 17.9 \\
Public administration, defence, education \& health & 99.3 & 15.2 & 1,064.4 & 14.4 \\
Other services & 5.5 & 0.8 & 190.1 & 2.6 \\
\midrule
Gross value-added & 653.7 & 100.0 & 7,384.1 & 100.0 \\
Taxes less subsidies on products & 17.2 & & --- & \\
GDP at market prices & 670.9 & & 7,384.1 & \\
GDP per capita, national currency & 592.4 & & 4,705.7 & \\
\textit{GDP per capita, \intlnineteen{}} & \textit{2,685} & & \textit{8,425} & \\
\bottomrule
\end{tabular}
\begin{tablenotes}[para,flushleft]
\footnotesize
\textit{Sources:} 1938: Valge's estimates for agriculture, industry, construction and housing, with services re-estimated as described in Section~\ref{sec:method} (see online appendix). 1990: \cite[184--6]{bolotin1992}, converted to roubles with Bolotin's sectoral purchasing-power parities; services other than trade and transport split using the 1987 Soviet input--output table (\cite{wbinputoutput1987}). \textit{Notes:} the 1990 figure includes taxes less subsidies within sectoral values (`---': included in sectoral values). Population: mid-year.
\end{tablenotes}
\end{threeparttable}
\end{table}


To express this figure in international dollars, a bilateral comparison of consumer prices with Sweden was carried out for 1938, using prices of some fifty consumer goods and services in both countries, weighted by the structure of household expenditure in each. The Fisher price ratio is 1.79 Swedish kronor per Estonian kroon, against an official exchange rate of 1.06: the kroon bought in Estonia roughly 70 per cent more than its exchange value would have bought in Sweden. Applied to Swedish GDP per head of 1,931 kronor, the ratio implies that Estonian GDP per head was 54.8 per cent of Sweden's, and since Sweden stood at 4,901~\intlnineteen{} in 1938, Estonia's 1938 benchmark is \ee{1938}~\intlnineteen{}, or \eeb{1938}~\intltwenty{}.\footnote{\autocite{mpd2023}, with \intltwenty{} converted at 1.594 (\autocite[12]{bolt_zanden_newgdp2021}). The price comparison follows the procedure of \autocite{klimantas2019} and \autocite{norkuslatvia1935}; the underlying prices and weights are given in the online appendix.} The result is 2 per cent below Valge's 2,745, but the proximity is accidental: his domestic-currency estimate was too low by 15 per cent and his Finnish bridge too generous by a similar margin. Expressed as a share of Finland, the new benchmark gives \pctfin{1938} per cent for 1938 against Clark's 90 per cent; Norkus and Markevi\v{c}i\={u}t\.{e}'s independent approximation gives 71 per cent.\footnote{\autocite[154]{norkus_restorations}.}

\textit{The 1990 benchmark.} The Soviet-period series is anchored to Bolotin's estimate of ESSR GDP in 1990, the first systematic attempt to express the accounts of all fifteen republics in international dollars.\footnote{\autocite[181--92]{bolotin1992}.} Bolotin started from the republican GDP totals for 1990 compiled by Goskomstat in 1990--1, 1990 being the first year for which republican accounts on the SNA definition were compiled, and converted them into dollars with purchasing-power parities estimated, as he put it, `in practically full conformity' with the International Comparison Project (ICP). The parities, by sector of origin, were 0.40 dollars per rouble in agriculture, 1.80 in industry, 1.10 in construction, 2.80 in transport, 1.45 in trade and 3.30 in other services, or 1.82 for GDP as a whole.\footnote{\autocite[181, 184]{bolotin1992}.} For Estonia the result is \ee{bench1990}~\intlnineteen{} per head, the highest of the Baltic republics. Table~\ref{tab:benchmarks} shows the sectoral composition in roubles, with services other than trade and transport allocated by the 1987 Soviet input--output table;\footnote{\autocite{wbinputoutput1987}.} these rouble-price shares, domestic-price weights like the kroon shares of 1938, are the weights with which the volume indices are combined, the dollar conversion applying to the total alone.

Three features of this benchmark need to be stated plainly, because the reconstruction stands or falls with it. First, its level is conservative by the standards of the time. Bolotin's figure for the Soviet Union as a whole, 6,275 dollars per head, is below Maddison's own estimate for 1987, and the 1990 comparison on which it draws, a bilateral Austria--USSR price comparison conducted in the ICP framework, placed the USSR at 32 per cent of the United States where the CIA had placed it at 42 per cent.\footnote{\autocite[318--19]{Maddison1998b}.} Second, the parities are Union-wide, so that the conversion assumes that relative prices in Estonia were those of the Soviet Union as a whole. For the material sectors this is a reasonable approximation in a planned economy with administered prices; for services, where Bolotin's parity of 3.30 dollars per rouble reflects the very low rouble wages of the non-market sector, it is less certain. A separate question is the allocation of value-added among services with the 1987 input--output table: the resulting shares of trade and services in Estonian GDP at rouble prices (Table~\ref{tab:benchmarks}) are higher than those implied by the World Bank's 1990 sectoral data for Estonia, 16 per cent for agriculture, 47 per cent for industry and construction and 38 per cent for services.\footnote{\autocite{wbstatisticalhandbook1994}; \autocite[1--8]{wb1993estonia}.} The sensitivity of the series to this allocation is examined in Section~\ref{sec:limitations}. Third, there is no alternative. Ponomarenko's retrospective accounts, the most thorough archival reconstruction of Soviet republican national accounts, cover Russia alone, make no dollar conversion and, in their author's words, show that the official exchange rate `was greatly overvalued', so that any republican benchmark must rest on purchasing-power parities of Bolotin's kind.\footnote{\autocite[6, 15, 82]{ponomarenko2002}.} The CIA's 1982 benchmark cannot be decomposed by republic, and the 1987 input--output table understates services too severely to serve as a benchmark.\footnote{\autocite[66--89]{wbinputoutput1987}.} An independent Estonian benchmark with republic-specific prices is the natural next step; it would alter only the level to which the volume indices are anchored, not the indices themselves.

The choice of 1990 rather than an earlier year means that the Soviet-period series is extrapolated backwards for more than fifty years from end-of-period weights. The practical reason is that 1990 is the only year for which republican accounts on an SNA definition and purchasing-power parities exist together. Two further considerations favour it: the reporting of 1989--90, at the height of glasnost, was more transparent than that of any earlier Soviet year, the manipulation of earlier returns having been publicly acknowledged by Baltic statisticians in 1988;\footnote{\autocite{adlys1988}.} and a late benchmark captures the economy at its most diversified, so that branches that barely existed in 1940, electrical engineering, chemicals or oil-shale electricity, receive weights appropriate to their eventual size. The index-number consequences are treated in Section~\ref{sec:limitations}; the short answer is that they can be tested directly, because the Soviet-period series meets an independent series with 1938 weights in 1937--40.

\subsection{Volume indices and their sources}

Table~\ref{tab:sources} summarises the indicator and sources for each sector in each period; the online appendix documents every series year by year. Three general points precede the detail. First, the interwar indices are built from the publications of the Estonian State Statistical Central Bureau, the agricultural surveys of 1925, 1929 and 1939 and the compilations of Valge and of Klesment and Valge, and are weighted by the revised 1938 benchmark.\footnote{\autocite{valge2003all}; \autocite{klesment2007}.} Second, the Soviet-period indices are built from the statistical yearbooks of the ESSR, \textit{Eesti NSV rahvamajandus}, which were published from 1957 and give annual physical output from the mid-1950s. These are supplemented, for the years 1945--55 that the yearbooks cover incompletely or not at all, by the returns of the republican statistical administration held in the Estonian National Archives (fond R-10) and, for the weights within industrial branches, by an archival compilation of 1986 output values for several hundred products (fond R-1969).\footnote{See Table~\ref{tab:sources}. The archival sources are listed in the online appendix, section A.2.} Archival material is thus used systematically, not merely to fill gaps. Third, the series cover the civilian economy as reported to the republican statistical office. Military production, organised through all-Union enterprises and never reported at the republican level, is not captured: Estonia hosted fewer such plants than Latvia, but the largest, the Dvigatel works in Tallinn and the uranium plant at Sillam\"{a}e, lay outside republican statistics, an unquantifiable source of understatement;\footnote{\autocite[150]{kukk2005}; \autocite[66]{ahderajasalu1992}.} the size of the Soviet forces stationed in the republic is likewise unknown and is treated as a source of bias in Section~\ref{sec:limitations}.

% tab:sources
\begin{table}[htbp]
\centering\footnotesize
\caption[Volume indicators and sources by sector]{\textsc{Volume indicators and sources by sector}}
\label{tab:sources}
\begin{threeparttable}
\begin{tabular}{>{\RaggedRight}p{2.6cm}>{\RaggedRight}p{5.3cm}>{\RaggedRight}p{5.3cm}c}
\toprule
Sector & Interwar series, 1919--40 (1938 weights) & Soviet-period series, 1937--95 (1990 weights) & App. \\
\midrule
Crops & Physical output of 15 crops weighted by 1938 value-added (Valge; 1939 agricultural survey); yields from \cite{klesment2007} & Physical output; weights from World Bank (1996) and Eurostat agricultural accounts; 1945--65 minor crops interpolated between quinquennial observations & A.1 \\
Livestock & Physical output of meat, milk, eggs and wool; 1938 value-added weights & Physical output; during 1941--44 animal numbers multiplied by pre-war yields & A.1 \\
Forestry and fishing & Timber removals and fish catch; timber for 1919--21 and fish for 1919--27 extrapolated from the crops-and-livestock index & Timber removals and fish catch; 1941--44 interpolated & A.1 \\
Industry & Employment adjusted for real wages (no usable annual output data; Valge's current-price series cannot be deflated by branch) & 101 physical product series in 11 branches; branch weights from the 1990 benchmark, product weights from 1986 archival output data (ERA) & A.2 \\
Construction & Domestic cement use (production plus net imports), as in Valge; 1919--21 from population & Weighted index of material inputs (1987 Soviet input--output weights), following the CIA procedure & A.3 \\
Transport and communications & Railway passenger- and tonne-kilometres, merchant fleet tonnage, postal items, telegrams and telephone calls & Freight tonne-kilometres and passenger-kilometres for all modes; postal, telephone and broadcasting indicators; 1941--45 interpolated between 1940 and the first post-war observation, with 1943--45 movements from Lithuanian ratios (Section~\ref{sec:method}) & A.4 \\
Trade, finance and insurance & Marketed value-added of agriculture and industry plus deflated imports (`commerce' index); finance assumed to move with commerce & Soviet retail and wholesale turnover at comparable prices (archival and published); finance and insurance from employment with a productivity adjustment & A.5 \\
Housing & Valge's housing value-added for 1929--38, extrapolated with the urban population & Housing stock from construction and population data & A.5 \\
Public administration, defence, health and education & Civil services assumed to move with GDP; defence from military personnel (15,900 in 1921--39) & Employment by branch (Soviet yearbooks, archival records); defence from the 1940 and 1990 totals & A.6 \\
Other services & Total population & Employment & A.6 \\
\bottomrule
\end{tabular}
\begin{tablenotes}[para,flushleft]
\footnotesize
\textit{Sources:} Interwar: \cite{valge2003all}; \cite{klesment2007}; interwar statistical yearbooks and agricultural surveys. Soviet period: \textit{Eesti NSV rahvamajandus} statistical yearbooks (1957--89), \textit{Eesti statistika aastaraamat} (1990--98), Estonian National Archives (ERA, fond R-10 and R-1969), German occupation statistical bulletins (\textit{Statistische Berichte f\"ur das Ostland}, 1942--44), and World Bank (1993, 1996). Full source lists by series are given in the online appendix, sections A.1--A.6.
\end{tablenotes}
\end{threeparttable}
\end{table}


\textit{Agriculture, forestry and fishing.} In both periods the index combines physical output of crops (fifteen products, including fodder) and of livestock products (meat, milk, eggs and wool) with timber removals and the fish catch, each sub-index weighted by its benchmark-year value-added. Interwar output comes from Klesment and Valge's compilation and the agricultural surveys, Soviet-period output from the yearbooks and, for the early post-war years, the archives, where some minor crops are reported only quinquennially and have been interpolated.\footnote{\autocite{klesment2007}; online appendix, section A.1.} The weights for the Soviet period derive from the World Bank's 1990 agricultural accounts for Estonia and, for the split within grains, from the Eurostat economic accounts for agriculture.

\textit{Industry.} No annual volume index of interwar industrial output exists, and Valge's current-price series cannot be deflated by branch for want of producer price indices. The interwar industrial index is therefore built from industrial employment adjusted for the movement of real wages, a recognised second-best when the share of industry in GDP is modest.\footnote{\autocite[22]{broadberryinmorys}; \autocite[427--8]{thomasfeinstein2003}.} For the Soviet period, by contrast, physical output series could be assembled for 101 products in eleven branches: electricity, petroleum products, other fuels including oil shale, metals, chemicals, machinery and metalworking, wood and paper, building materials, light industry, food and a residual. Products are aggregated to branches with the 1986 archival weights and branches to industry with the weights of the 1990 benchmark, following the CIA's procedure for Soviet GNP.\footnote{\autocite[48]{cia1982a}; \autocite[312]{maddison2010}.}

\textit{Construction.} The interwar index follows Valge in tracking domestic consumption of cement, production plus net imports; it is volatile in the 1920s, when timber construction, which the index does not capture, was still important.\footnote{\autocite[Table 35]{valge2003all}.} The Soviet-period index is a weighted index of the material inputs into construction, with weights from the 1987 input--output table, which is the CIA's procedure.\footnote{\autocite[66--89]{wbinputoutput1987}.}

\textit{Transport and communications.} The interwar index combines railway passenger- and tonne-kilometres, the tonnage of the merchant fleet, and postal items, telegrams and telephone calls. The Soviet-period transport index combines freight tonne-kilometres and passenger-kilometres across all modes; communications combine postal and telegraph traffic, telephone calls and radio and television sets in use.

\textit{Trade, finance and housing.} For the interwar years a single `commerce' index tracks the marketed value-added of agriculture and industry together with deflated imports, and finance and insurance are assumed to move with it. Under planning, industrial products circulated by allocation rather than through markets, so the Soviet-period trade index is built instead from the retail and wholesale turnover statistics at comparable prices compiled by the republican statistical office, most of them archival; these are the only value series used in the reconstruction, and the hidden inflation they may contain is bounded in Section~\ref{sec:limitations}. Finance and insurance follow employment, adjusted for productivity with the ratio of trade turnover to trade employment, as the CIA did; housing follows the housing stock, derived from construction and population.\footnote{\autocite[319]{Maddison1998b}; \autocite[48]{cia1982a}.}

\textit{Public and other services.} Interwar public services other than defence are assumed to move with the rest of the economy, since no consistent series of civil employment exists; defence follows military personnel, which is known for 1919--20 and 1939 and held constant in the intervening years. Soviet-period public administration, health and education follow employment by branch, as in the CIA's accounts, with defence derived from the 1940 and 1990 totals. Other, mainly domestic, services follow population before 1940 and employment thereafter.

\textit{The war years.} Estonia was under German civil administration from 1941 to 1944, and the \textit{Reichskommissariat Ostland} published statistical bulletins that cover agriculture, the extractive industries and employment but little else. Where these fail, forestry, fishing, transport and several service branches are interpolated between 1940 and the first post-war observation, and the movement between 1943 and 1945, for which no Estonian figures could be found, is taken from the ratios observed in Lithuania, where documentation is better. The estimates for 1941--5 are therefore `rough' in Feinstein's sense, with a margin of 15--25 per cent, and the 1944 trough is an order of magnitude rather than a measurement.\footnote{\autocite[21]{feinstein1972}.} Because the Soviet-period series is extrapolated backwards through the war, an error in the wartime linkage would shift its level in 1937--40, so that the overlap tests the linkage together with the rest of the extrapolation; the spliced series, which follows the interwar series to 1940, is affected only through the 1940 average.

\subsection{Population}

Estonia's borders changed once, when the area east of Narva and most of Petseri county, together about 6 per cent of the interwar population, were transferred to the Russian republic in 1944--5. The interwar figures refer to the borders of the republic and the Soviet-period and modern figures to present-day borders, which is the convention of the MPD; the transferred districts were rural and poorer than the national average, so that the break in income per head is likely to be small, and no adjustment has been made.\footnote{\autocite[30]{bolt_zanden_newgdp2021}.} The population series rests on the censuses of 1922, 1934, 1959, 1970, 1979 and 1989 and the year-end estimates of Statistics Estonia and the `Baltic 100' project, with mid-year figures as the average of consecutive year-ends; wartime losses, deportations and the flight of some 70,000 people to the West in 1944 are taken from the demographic literature, and the Soviet garrison of 1940--1 is added to the resident population (Table~\ref{tab:series}).

\subsection{Aggregation and splicing}

Each of the two new series is first compiled in constant benchmark-year prices and then expressed in \intlnineteen{} through its benchmark. The interwar series covers 1919--40 and the Soviet-period series 1937--95, so that the two overlap for four years. They are spliced at 1940 by taking the average of the two estimates for that year, and the Soviet-period series is in turn spliced at 1991 to a modern series obtained by backcasting the MPD's 2011 benchmark with the World Bank's constant-price growth rates of GDP per capita, which are used in preference to the MPD's own 1990--2010 growth rates for reasons explained in Section~\ref{sec:limitations}.\footnote{\autocite{mpd2023}; \autocite{wdi2024}.} The 1913 figure is the benchmark of Norkus and Markevi\v{c}i\={u}t\.{e}, 2,058~\intlnineteen{}, taken as published; there are no estimates for 1914--18, and the change between 1913 and 1919 therefore combines their income-side approximation with the new production-side estimate and should be read with that caveat. Following Bolt and van Zanden, the series is compiled in \intlnineteen{} and converted to \intltwenty{} at 1.594, and both versions are reported.\footnote{\autocite[12, 30]{bolt_zanden_newgdp2021}.}

% =============================================================
% III. THE NEW SERIES
% =============================================================
\section{The new series}
\label{sec:results}

The annual series of GDP per capita and population for 1913 and 1919--2022 is reported in full in Table~\ref{tab:series} of the appendix, and the sectoral volume indices from which it is built in Tables~\ref{tab:appsec1} and~\ref{tab:appsec2}. Table~\ref{tab:growth} summarises growth by sub-period alongside the comparators that recur in the discussion below, and Table~\ref{tab:sectoral} gives the growth of each sector.

% tab:growth
\begin{table}[htbp]
\centering\small\setlength{\tabcolsep}{3.5pt}
\caption[Average annual growth of GDP per capita by sub-period (per cent)]{\textsc{Average annual growth of GDP per capita by sub-period (per cent)}}
\label{tab:growth}
\begin{threeparttable}
\begin{tabular}{lrrrrrrrr}
\toprule
Period & Estonia & Lithuania & Finland & Nordics & Western core & Western periphery & ECE & USSR \\
\midrule
\textit{1913--38} & 1.1 & 0.9 & 2.1 & 1.9 & 1.0 & 1.2 & 0.8 & 1.7 \\
\midrule
\textit{1919--29} & 2.8 & 4.5 & 5.1 & 3.3 & 3.5 & 3.1 & 4.8 & 9.2 \\
\textit{1929--38} & 2.7 & 3.0 & 3.1 & 2.1 & 1.1 & 0.8 & 0.1 & 5.0 \\
\midrule
\textit{1938--50} & 2.6 & 0.9 & 1.4 & 2.1 & 0.6 & $-$0.1 & 0.8 & 2.4 \\
\textit{1950--68} & 3.9 & 4.9 & 3.6 & 3.1 & 3.6 & 5.2 & 3.9 & 3.4 \\
\textit{1968--89} & 0.9 & 1.8 & 3.6 & 2.5 & 2.4 & 3.0 & 1.7 & 1.5 \\
\textit{1938--89} & 2.4 & 2.7 & 3.1 & 2.6 & 2.4 & 3.0 & 2.3 & 2.4 \\
\midrule
\textit{1989--95} & $-$5.8 & $-$6.6 & $-$0.8 & 1.6 & 1.3 & 2.1 & $-$1.6 & $-$8.0 \\
\textit{1995--2022} & 4.2 & 5.0 & 1.7 & 2.1 & 1.5 & 1.4 & 4.4 & 3.9 \\
\bottomrule
\end{tabular}
\begin{tablenotes}[para,flushleft]
\footnotesize
\textit{Sources:} Estonia: see online appendix. Lithuania: \cite{klimantas2026soviet}. Other series: \cite{mpd2023}, supplemented as documented in the online appendix. \textit{Notes:} compound annual rates between end-point years, \intlnineteen{}. Nordics: Denmark, Finland, Iceland, Norway and Sweden; Western periphery: Finland, Greece, Iceland, Ireland, Italy, Portugal and Spain; Western core: Austria, Belgium, Denmark, France, Germany (West Germany 1950--97), the Netherlands, Norway, Sweden, Switzerland and the UK; ECE (East-Central Europe): Bulgaria, Czechoslovakia, Hungary, Poland, Romania and Yugoslavia; USSR: Soviet Union within 1990 borders (former Soviet Union after 1991).
\end{tablenotes}
\end{threeparttable}
\end{table}


% tab:sectoral
\begin{table}[htbp]
\centering\scriptsize\setlength{\tabcolsep}{3pt}
\caption[Average annual growth of sectoral value-added volume indices, Estonia (per cent)]{\textsc{Average annual growth of sectoral value-added volume indices, Estonia (per cent)}}
\label{tab:sectoral}
\begin{threeparttable}
\begin{tabular}{lrrrrrrcc}
\toprule
Sector & 1919--29 & 1929--38 & 1938--50 & 1950--68 & 1968--89 & 1989--91 & 1938 level & Peak \\
& & & & & & & regained in & year \\
\midrule
Crops & 2.6 & 2.8 & 1.1 & 2.2 & $-$1.1 & $-$8.7 & 1947 & 1968 \\
Livestock & 7.7 & 1.5 & $-$3.9 & 3.7 & 2.0 & $-$8.6 & 1960 & 1987 \\
Forestry and fishing & 4.7 & 0.6 & 0.9 & 7.6 & 1.4 & $-$12.4 & 1947 & 1976 \\
\quad Primary sector & 5.2 & 1.9 & $-$1.3 & 3.1 & 0.6 & $-$8.9 & 1959 & 1986 \\
\midrule
Food & --- & --- & 2.2 & 8.8 & 2.1 & $-$9.7 & 1948 & 1976 \\
Textiles & --- & --- & 4.6 & 9.0 & 0.9 & $-$5.5 & 1949 & 1987 \\
Other light industries & --- & --- & 2.8 & 7.1 & 1.8 & $-$12.4 & 1948 & 1988 \\
Heavy industries & --- & --- & 5.1 & 12.0 & 1.2 & $-$8.6 & 1948 & 1980 \\
\quad Industry & 11.5 & 5.7 & 5.0 & 9.5 & 1.4 & $-$9.1 & 1948 & 1988 \\
\midrule
Construction & 11.7 & 7.6 & 5.4 & 9.3 & 0.9 & $-$9.9 & 1947 & 1977 \\
Transport and communications & 10.7 & 2.0 & 8.0 & 8.9 & 4.9 & $-$4.2 & 1945 & 1989 \\
Trade, accommodation and catering & --- & --- & 1.2 & 9.6 & 4.7 & $-$22.0 & 1949 & 1990 \\
Housing & 1.8 & 6.8 & 1.1 & 2.8 & 2.0 & $-$4.6 & 1949 & 1986 \\
Public services & $-$5.3 & 1.3 & 4.0 & 1.2 & 1.8 & $-$2.5 & 1945 & 1984 \\
\midrule
\textbf{GDP} & 3.3 & 2.9 & 1.9 & 5.4 & 1.7 & $-$8.7 & 1947 & 1989 \\
\bottomrule
\end{tabular}
\begin{tablenotes}[para,flushleft]
\footnotesize
\textit{Sources:} see online appendix. \textit{Notes:} compound annual growth between end-point years of constant-price value-added indices (1938 = 100). Industrial branch indices are available from 1938 only. `1938 level regained': first post-war year (1945 or later) in which the index is at or above its 1938 value; `Peak year': year of the highest index value up to 1990. Public services comprise administration, defence, health and education. Growth in 1938--50 is affected by the wartime trough of 1944, for which Estonian sectoral movements are partly extrapolated from Lithuanian ratios (Section~\ref{sec:limitations}). The 1989--91 column ends in 1991, the splicing year beyond which the final series follows adjusted Maddison Project data.
\end{tablenotes}
\end{threeparttable}
\end{table}


On the eve of the First World War Estonian income per head stood at \ee{1913}~\intlnineteen{}. The war and the fighting of 1918--20 cut it, on the two benchmarks available, which rest on different methods, by more than a fifth, to \ee{1919}~\intlnineteen{} in 1919, a level that was not regained until 1929. Recovery in the 1920s was steady rather than rapid, \gr{1919}{1929} per cent a year, and was carried by the sectors that had suffered most in the war: industry, construction and transport each grew by more than 10 per cent a year from their very low 1919 base, while agriculture, which had been less damaged, grew by \sgr{Primarysectortotal}{1919}{1929} per cent a year (Table~\ref{tab:sectoral}). The Depression was shallow. Income per head fell from \ee{1930}~\intlnineteen{} in 1930 to \ee{1932}~\intlnineteen{} in 1932, a decline of 8 per cent, and had recovered by 1934. The second half of the decade was the fastest phase of interwar growth: with construction growing by \sgr{Construction}{1929}{1938} per cent a year over 1929--38, industry by \sgr{Industrytotal}{1929}{1938} per cent and housing by \sgr{Housing}{1929}{1938} per cent, income per head reached \ee{1938}~\intlnineteen{} in 1938 and an interwar peak of \ee{1940}~\intlnineteen{} in 1940, the year of the Soviet occupation. Over the two decades of independence as a whole, income per head grew by \gr{1919}{1938} per cent a year between 1919 and 1938, but by only \gr{1913}{1938} per cent a year measured from 1913: the recovery from the war absorbed most of the gain.

The Second World War reduced income per head to \ee{1944}~\intlnineteen{} in 1944, a third below the level of 1940, though as Section~\ref{sec:method} explained the 1944 figure is partly extrapolated and should not be pressed. Reconstruction was rapid: GDP regained its 1938 level in \sregain{TotalGVAGDP}, industry in \sregain{Industrytotal} and income per head reached \ee{1950}~\intlnineteen{} in 1950, more than a third above its 1938 level. The result is striking beside Lithuania, where income per head grew by less than 1 per cent a year over 1938--50, and Finland, where it grew by \grx{Finland}{1938}{1950} per cent (Table~\ref{tab:growth}), and two things account for most of it. The first is the denominator: the resident population was 8 per cent smaller in 1950 than in 1938, so that total GDP rose by a quarter where income per head rose by a third. The second is the weight of investment in oil shale, electricity and textiles in the fourth five-year plan, which took industrial output to 1.8 times its 1938 level and transport to 2.5 times by 1950 (Table~\ref{tab:appsec1}). Because the Soviet-period series is anchored in 1990, the 1950 level does not depend on the wartime observations but on growth over 1950--90. Agriculture did not share in this recovery. Crop output was back at its 1938 level by \sregain{Crops}, but livestock output, which had been halved by the war and was then disrupted by collectivisation, did not regain it until \sregain{Livestock}, and the agricultural sector as a whole not until \sregain{Primarysectortotal}, twenty years after the pre-war peak.

The eighteen years from 1950 to 1968 are the most dynamic in the series before 1995. Income per head grew by \gr{1950}{1968} per cent a year, from \ee{1950} to \ee{1968}~\intlnineteen{}, with total GDP growing by \tgr{1950}{1968} per cent a year. Heavy industry grew by \sgr{Heavyindustries}{1950}{1968} per cent a year, textiles by \sgr{Textiles}{1950}{1968} per cent and food processing by \sgr{Food}{1950}{1968} per cent; construction, transport and trade each grew by around 9 per cent a year; and agriculture, now recovering, by \sgr{Primarysectortotal}{1950}{1968} per cent. Growth then slowed sharply and, after 1973, almost ceased. Income per head rose by \gr{1968}{1989} per cent a year between 1968 and 1989, and by only \gr{1973}{1989} per cent a year after 1973; between 1980 and 1989 it rose by \gr{1980}{1989} per cent a year. The sectoral indices show that the slowdown was general. Food processing peaked in \speak{Food}, construction in \speak{Construction}, heavy industry in \speak{Heavyindustries} and the housing stock in \speak{Housing}; only textiles, other light industry, livestock, transport and trade continued to grow into the late 1980s, and industrial output as a whole peaked in \speak{Industrytotal}. The series reaches its Soviet-era maximum of \ee{1989}~\intlnineteen{} in 1989.

The collapse that followed was deeper than anything in the interwar record. Income per head fell by \eeTROUGHpct{} per cent between 1989 and \eeTROUGHyear{}, to \eeTROUGHval~\intlnineteen{}, with the sharpest fall, of 19 per cent, in 1992, the year in which the kroon replaced the rouble and trade with the former Union contracted. Growth resumed in 1995, averaged \gr{1995}{2022} per cent a year to 2022, and carried income per head past its 1989 level in \eeREGAINnineteen{} and to \ee{2022}~\intlnineteen{} (\eeb{2022}~\intltwenty{}) in 2022, second among the Baltic states behind Lithuania.

Table~\ref{tab:previous} compares the new series with the three earlier estimates. For the 1930s the new series and Valge's differ by only a few per cent, which, as explained above, reflects offsetting errors in his domestic-currency total and his conversion ratio rather than agreement in method. For the 1920s the differences are larger and change sign, and the 29 per cent rise that Valge's deflation produces for 1924 becomes a fall of 1 per cent in the new series, in which crops, livestock, industry and transport move by between $-$2 and $+$5 per cent in that year and the small net fall reflects the volatile cement-based construction index (Table~\ref{tab:appsec1}). The series of Klesment, Puur and Valge lies between a quarter and a third below the new one throughout 1960--89, and more than half below in 1950, because of its low UNECE benchmark and the omission of trade and the non-material sphere; its implied growth of income per head over 1950--78, about 4.4 per cent a year, is well above the \gr{1950}{1978} per cent found here, because their index omits the non-material sphere, whose employment-based components grew far more slowly than material production, and weights the four material branches by their shares of net material product, which overstate the share of industry in GDP. The MPD figures, by contrast, lie above the new series, by 7 per cent in 1973 and by 32 per cent in 1989, because they carry modern benchmarks backwards with NMP growth rates that overstate real growth in the final Soviet decades. The new series thus lowers Estonian income in the late Soviet period relative to the MPD, raises it relative to the only previous long series, and replaces both for the years they cover.

% =============================================================
% IV. VALIDATION, LIMITATIONS AND SENSITIVITY
% =============================================================
\section{Validation, limitations and sensitivity}
\label{sec:limitations}

\subsection{The overlaps}

The reconstruction can be tested at two points. The interwar series, built from Estonian sources and weighted by the 1938 benchmark, and the Soviet-period series, built from Soviet and archival sources and weighted by the 1990 benchmark, both cover the years 1937--40; and the Soviet-period series and the modern series, backcast from the MPD's 2011 benchmark, both cover 1990--5. Table~\ref{tab:overlaps} reports the results.

% tab:overlaps
\begin{table}[htbp]
\centering\small
\caption[Overlapping estimates of Estonian GDP per capita at the two splicing points]{\textsc{Overlapping estimates of Estonian GDP per capita at the two splicing points}}
\label{tab:overlaps}
\begin{threeparttable}
\begin{tabular}{lrrcc}
\toprule
Year & Earlier-benchmark series & Later-benchmark series & Difference (\%) & Feinstein grade \\
& (\intlnineteen{}) & (\intlnineteen{}) & & \\
\midrule
\multicolumn{5}{l}{\textit{(a) Interwar series (1938 benchmark) and Soviet-period series (1990 benchmark)}} \\
1937 & 2,722 & 2,678 & $-$1.6 & Firm \\
1938 & 2,685 & 2,782 & +3.6 & Firm \\
1939 & 2,736 & 2,654 & $-$3.0 & Firm \\
1940$^{\dagger}$ & 2,776 & 2,767 & $-$0.3 & Firm \\
\midrule
\multicolumn{5}{l}{\textit{(b) Soviet-period series (1990 benchmark) and adjusted MPD series (2011 benchmark)}} \\
1990 & 8,425 & 8,172 & +3.0 & Firm \\
1991$^{\dagger}$ & 7,506 & 7,558 & $-$0.7 & Firm \\
1992 & 5,400 & 6,069 & $-$12.4 & Good \\
1993 & 4,916 & 5,870 & $-$19.4 & Rough \\
1994 & 4,973 & 5,899 & $-$18.6 & Rough \\
1995 & 4,989 & 6,277 & $-$25.8 & Conjectural \\
\bottomrule
\end{tabular}
\begin{tablenotes}[para,flushleft]
\footnotesize
\textit{Sources:} see online appendix; adjusted MPD series from \cite{mpd2023} with 1990--2010 growth replaced by World Bank constant-price growth backcast from the 2011 benchmark (\cite{wdi2024}). \textit{Notes:} Difference = (later $-$ earlier)/earlier in panel (a) and (earlier $-$ later)/earlier in panel (b), so that a positive number always means that the 1990-based series is the higher one. Feinstein's grades: firm ${<}5$ per cent, good 5--15, rough 15--25, conjectural ${>}25$ (\cite[21]{feinstein1972}). $^{\dagger}$~Splice year: the final series takes the arithmetic average of the two estimates.
\end{tablenotes}
\end{threeparttable}
\end{table}


In 1937--40 the two series differ by $-$1.6, $+$3.6, $-$3.0 and $-$0.3 per cent. All four differences fall within the 5 per cent margin that Feinstein reserved for `firm' estimates, and they do not share a sign.\footnote{\autocite[21]{feinstein1972}.} The test is a demanding one. The two series share no sources: the interwar indices come from the republic's own statistics and the Soviet-period indices from the yearbooks of the ESSR and the post-war archives. They share no weights, since one uses the sectoral structure of 1938 and the other that of 1990. And they are anchored to benchmarks constructed by different methods and separated by fifty-two years, one a bilateral price comparison with Sweden, the other a multilateral purchasing-power-parity study of the Soviet republics. That the Soviet-period series, extrapolated backwards over half a century of war, collectivisation and forced industrialisation, arrives within 4 per cent of an independently measured level is the strongest evidence available that the net error of the extrapolation is small. It cannot exclude errors that happen to offset, such as an overstatement of growth from hidden inflation in trade turnover against an understatement from the index-number assumptions. But the two index-number biases discussed below, late base-year weights and fixed value-added weights, would both understate growth and so would both place the backcast 1937--40 level above the interwar series; the overlap therefore tests their combined size. A net understatement of growth of as little as 0.1 percentage point a year over 1940--90 would place the Soviet-period series 5 per cent above the interwar series in every overlap year, which is not observed. Broadberry and Burhop proposed that time-series projections should be regarded as consistent with direct benchmarks when they agree within 10 per cent; a projection over thirty years in their own work landed within 4 per cent of the benchmark and was judged a success.\footnote{\autocite[323--6]{broad2007}; \autocite[933]{broad2008}.} The Estonian overlap meets that standard over a longer span, and does so in each of four consecutive years.

The second overlap is less tidy, and the reasons are instructive. In 1990 the Soviet-period series stands 3.0 per cent above the modern series and in 1991 0.7 per cent below it; thereafter the gap widens to 12 per cent in 1992 and 26 per cent by 1995. The close agreement in 1990--1 is itself a validation of the 1990 benchmark, and in particular of the conversion from roubles to dollars on which a reviewer might fasten. Bolotin's dollar figure for Estonia was obtained with purchasing-power parities, not exchange rates, and the modern series against which it is tested derives from an entirely independent ICP round, that of 2011, carried backwards with the World Bank's constant-price growth rates. Two routes that share no data arrive at the same level of Estonian income in 1990--1 within 3 per cent. The divergence after 1991 arises because the Soviet-period series continues to value the physical collapse of Soviet-era output at 1990 weights and, in trade, to rely on deflated turnover figures that understate inflation: it records what happened to the Soviet-era economy, not what replaced it. The modern series is therefore preferred from 1991, the year in which the two agree, and the MPD's own growth rates for 1990--2010, which are chained to its NMP-based level for the late Soviet years, are replaced by the World Bank's. The MPD's 1989 figure for Estonia lies 32 per cent above the new series (Table~\ref{tab:previous}); the implied fall in income during the transition in the MPD is 36 per cent, against \eeTROUGHpct{} per cent here, so that the two sources agree on the depth of the collapse even though they disagree on the level from which it started.\footnote{\autocite{mpd2023}; the comparison is developed in \autocite{klimantas2026soviet}.}

\subsection{A single late benchmark and the index-number problem}

The Soviet-period series extrapolates 1990 weights backwards over 1937--90, a period of profound structural change. An index with a late base year assigns to each sector the relative prices of the end of the period, and since the sectors that grew fastest were, as a rule, those whose relative prices fell, late weights give less importance to fast-growing sectors and so understate growth; conversely, early weights overstate it. Feinstein and Thomas put the point generally: an index with an early base `will always show a greater increase in quantities than one constructed with a base year late in the period'.\footnote{\autocite[516]{feinsteinthomas2002}.} This is the effect that Gerschenkron first documented for Soviet machinery, and it was large in the Stalin period: Bergson's GNP for 1928--37 grows by 175 per cent at 1928 prices and by 62 per cent at 1937 prices.\footnote{\autocite[453]{gerschenkron1957}; \autocite[517]{feinsteinthomas2002}.} Three considerations establish that it is far smaller for the series presented here.

First, the effect is concentrated in the early phases of industrialisation, when relative prices change fastest; Jonas and Sardy showed that `most of the action which leads to the Gerschenkron effect takes place early'.\footnote{\autocite[85]{jonassardy1970}.} For the post-war Soviet economy the measured effect is modest. When the CIA moved its base year from 1970 to 1982 prices, its estimate of Soviet GNP growth over 1951--84 fell by about half a percentage point a year, from 5.6 to 5.2 per cent for the 1950s and from 3.2 to 2.4 per cent for the 1970s; an earlier comparison of 1960, 1970 and 1976 prices produced differences of a few tenths of a point, which the CIA judged `not large'.\footnote{\autocite[8]{cia1990}; \autocite[23--4]{cia1982a}; \autocite[156]{harrison1993}.} Half a point a year over forty years compounds to roughly a fifth of the level; this indicates the order of magnitude to be expected for the post-war decades rather than a bound, since the Estonian extrapolation also spans the forced industrialisation of the 1940s and 1950s, in which the effect is largest. Second, the bound must therefore come from the data, and the series can be tested. If 1990 weights understated Estonian growth over 1940--90, the Soviet-period series would arrive in 1937--40 systematically above the interwar series, which is weighted by the prices of 1938. It does not: the four differences are small and of both signs. The structural change within industry, where the Gerschenkron effect originates, is moreover handled at a finer level than the aggregate weights suggest, since the 101 product series are aggregated within eleven branches with 1986 weights before the branches are combined. Third, the procedure is the one used by every reconstruction of Soviet output with which the series might be compared. The CIA used 1970 weights for 1950--80 and 1982 weights for 1950--90, and Maddison, reviewing the practice, found `nothing untoward' in it; Kuboniwa's and Ponomarenko's reconstructions of Russian GDP for 1961--90 use, among their variants, the fixed structure of Russia's 1990 GDP.\footnote{\autocite[312]{Maddison1998b}; \autocite[57, 62]{ponomarenko2002}.} A chain-linked alternative would require annual current-price value-added by sector, which does not exist for any Soviet republic, and a 1938-weighted alternative for the whole period would require the sectors of 1990 to be valued at prices of 1938, when several of them did not exist. The reader who wishes to see the question settled empirically should turn to the 1937--40 overlap.

\subsection{Fixed value-added weights}

A related assumption is that the ratio of value-added to gross output within each sector remained at its benchmark-year level. The volume indices track gross output, or inputs, or employment, and are combined with value-added weights; the method is sometimes called single deflation, although the term properly describes the deflation of nominal value-added by an output price index, a procedure that this paper does not use and that performs much worse.\footnote{\autocite[443--4]{thomasfeinstein2003}.} Thomas and Feinstein's evaluation of the method actually used here, output indicators with value-added weights, found that it understated true value-added growth by between 0.15 and 1 percentage point a year in their simulations, depending on the pace of technical change, and they described it as `the next best procedure' to double deflation where intermediate inputs cannot be observed.\footnote{\autocite[427--8, 444]{thomasfeinstein2003}.} Double deflation is not available for either period. The interwar statistical office recorded no intermediate consumption, and for the Soviet period, as Ponomarenko notes, the price statistics that double deflation would require `practically did not exist'; the CIA applied double deflation only to agriculture and used gross output as the index of value-added everywhere else, on the argument, due to Hill, that gross output is a better proxy than inputs unless output and input growth diverge by more than two points a year.\footnote{\autocite[18--19]{ponomarenko2002}; \autocite[45]{cia1982a}.}

The bias is bounded with the help of the Swedish historical national accounts, which permit single- and double-deflated series to be compared over comparable phases of industrialisation.\footnote{\autocite[19]{swedengdp_15602010}.} For the interwar period, when input costs outside agriculture moved little, the maximum error is about 2 per cent of the level of GDP. For the Soviet period the issue is more serious, because the sustained shift towards intermediate-intensive heavy industry makes it likely that value-added ratios changed, and the maximum error is put at 14 per cent of the level of GDP over the half-century, the derivation being given in the online appendix. The sign adopted in Table~\ref{tab:sensitivity} is that of Thomas and Feinstein's simulations, an understatement of growth and hence an overstatement of the backcast early levels; 14 per cent over fifty years corresponds to about a quarter of a percentage point a year, in the lower part of their range, which is appropriate because the Estonian indices track gross output product by product rather than by broad branch. Two observations limit the practical concern. The error accumulates over the extrapolation, so that it is largest in the earliest years of the Soviet-period series, precisely where the 1937--40 overlap shows the series to be within 4 per cent of an independent level. And the opposite sign cannot be excluded, since the Soviet shift towards intermediate-intensive heavy industry would have raised rather than lowered value-added ratios if the rising degree of fabrication outpaced the relative fall in intermediate-input prices; this is why the overlap, not the simulations, is decisive. Gerschenkron himself, confronting constant gross-to-net ratios over half a century in Jostock's German accounts, thought them `not so very implausible', since the rising degree of fabrication is likely to have been offset by the relative fall in the prices of value-added components.\footnote{\autocite[453]{gerschenkron1957}.}

\subsection{Other sources of bias}

Table~\ref{tab:sensitivity} lists the remaining sources of bias, the sectors they affect and the maximum impact of each on the level of GDP; the bounds for the occupation period are derived in the same way as in the companion Lithuanian study, with the Estonian sectoral weights.

% tab:sensitivity
\begin{table}[htbp]
\centering\footnotesize\setlength{\tabcolsep}{4pt}
\caption[Sources of potential bias and their maximum impact on the level of GDP]{\textsc{Sources of potential bias and their maximum impact on the level of GDP}}
\label{tab:sensitivity}
\begin{threeparttable}
\begin{tabular}{p{2.5cm}p{6.6cm}p{2.4cm}r}
\toprule
Period & Source of bias & Sectors affected & Max. impact \\
& & & (\% of GDP) \\
\midrule
Interwar, 1919--40 & Service-sector value-added in the 1938 benchmark & Services (34\% of GDP) & $\pm$3 \\
 & Fixed value-added weights (`single deflation') & All & $\pm$2 \\
 & Industrial output proxied by employment and real wages & Industry & $\pm$1 \\
 & Taxes less subsidies not tracked annually & All & $\pm$1 \\
 & No productivity adjustment in public services & Public services & $\pm$3 \\
 & Other services proxied by population & Other services & $\pm$1 \\
 & Military personnel known only for 1919--20 and 1939 & Defence & $\pm$1 \\
\midrule
Soviet period, 1940--90 & Service-sector allocation in the 1990 benchmark & Services & $\pm$4 \\
 & Taxes less subsidies absorbed in sectoral values & All & $-$0.5 \\
 & Fixed value-added weights (`single deflation') & All & $-$14 \\
 & Hidden inflation in comparable-price trade turnover & Trade & +2 \\
 & Size of Soviet forces stationed in Estonia unknown & Defence & $\pm$5 \\
 & No productivity adjustment in employment-based services & Services & $-$10 \\
 & Wartime movements 1943--5 partly borrowed from Lithuania & All (1941--5 only) & 15--25 \\
\bottomrule
\end{tabular}
\begin{tablenotes}[para,flushleft]
\footnotesize
\textit{Sources:} authors' assessment; see Section~\ref{sec:limitations} and the online appendix for the derivation of each bound. \textit{Notes:} signed entries indicate the direction in which the published series would be revised if the bias were corrected; $\pm$ entries are symmetric. The bounds are not additive.
\end{tablenotes}
\end{threeparttable}
\end{table}


Three deserve comment. The service sector of the 1938 benchmark, about one-third of GDP, rests on earnings and employment data and on Latvian structural ratios; an error of 10 per cent in the service estimate moves GDP by 3 per cent, and Norkus and Markevi\v{c}i\={u}t\.{e}'s independent approximation, 5 per cent below the new benchmark, suggests that any error is of that order. In the 1990 benchmark the allocation of value-added among services depends on Bolotin's Union-wide parities and the 1987 input--output table; a reallocation of a quarter of service value-added between branches moves GDP by about 4 per cent. The Soviet-period trade index uses turnover at comparable prices, the one value series in the reconstruction; Kontorovich's estimates of hidden inflation in Soviet retail trade imply an upward bias of at most 2 per cent of GDP.\footnote{\autocite[688--9]{kontorovich_2001}.} Defence after 1940 is unobserved, since the Soviet forces stationed in Estonia were an all-Union institution and their size is unknown between 1940 and 1990; the defence component of public services is interpolated between the 1940 and 1990 totals in step with civil employment, and the maximum bias is 5 per cent of GDP in either direction. The employment-based indices for health, education and administration make no allowance for productivity change, which understates service growth if productivity rose; the maximum effect is 10 per cent of the level. Finally, unregistered activity is estimated to have amounted to 10--20 per cent of Soviet GDP in 1989, but OECD economies displayed shadow economies of comparable size at the same date, so that the omission does not systematically impair international comparison.\footnote{\autocite[532--51]{kim2003}; \autocite[5]{baltictransitioncontraction}; \autocite[90--1]{shadow_europe}.}

Taken together, the quantifiable biases imply a margin of about 6 per cent around the interwar levels and, around the Soviet-period levels, of 14 per cent for the largest single bias, or about a fifth in the earliest years if the two same-signed biases of Table~\ref{tab:sensitivity}, fixed value-added weights and unadjusted service productivity, were both at their maxima; the overlap shows that the realised error is far smaller. The war years are subject to a wider margin. The sub-period growth rates on which the argument of the following sections rests are robust to all of them: no plausible combination of the biases listed in Table~\ref{tab:sensitivity} would raise growth in 1968--89 to the level of 1950--68, close the gap with Finland in 1938, or reverse Estonia's position relative to the Soviet average.

% =============================================================
% V. ESTONIAN GROWTH IN COMPARATIVE PERSPECTIVE
% =============================================================
\section{Estonian growth in comparative perspective}
\label{sec:implications}

The series spans three regimes: the capitalist republic of 1919--40, the planned economy of 1940--91 and the market economy restored after 1991. The questions that the literature on catch-up growth asks of each are the same. Gerschenkron's late industrialiser can borrow technology rather than invent it and so grow faster than the leader; the neoclassical model of Solow predicts the same from diminishing returns to capital.\footnote{\autocite[8]{gerschenkron1962}; \autocite[70--1]{solow1956}.} But Abramovitz and David showed that the potential is realised only where a follower possesses the `social capability' to absorb what it borrows, the literacy, institutions and commercial organisation that make imported technology productive, and where the leader's technology is congruent with the follower's endowments.\footnote{\autocite{abramovitz_david1996}.} North, and Acemoglu and Robinson, make the institutional condition explicit: convergence is conditional on institutions that reward productive investment.\footnote{\autocite{north1990}; \autocite{acemoglu_whynationsfail}.} Estonia is an unusually clean test of these propositions, because its social capability was, by the standards of the European periphery, high throughout, while its institutions were changed twice from outside. Tables~\ref{tab:relative}, \ref{tab:interwar}, \ref{tab:trade} and~\ref{tab:structure} provide the comparative evidence.

% tab:relative
\begin{table}[htbp]
\centering\small
\caption[Estonian GDP per capita relative to comparators, selected years]{\textsc{Estonian GDP per capita relative to comparators, selected years}}
\label{tab:relative}
\begin{threeparttable}
\begin{tabular}{lrrrrrrr}
\toprule
& Estonia & \multicolumn{6}{c}{Estonia as a percentage of} \\
\cmidrule(lr){3-8}
Year & (\intlnineteen{}) & Finland & Sweden & Western core & USSR & Lithuania & Latvia \\
\midrule
1913 & 2,058 & 97 & 72 & 51 & 146 & 126 & 87 \\
1919 & 1,600 & 96 & 57 & 48 & 279 & 157 & 108 \\
1929 & 2,113 & 78 & 52 & 45 & 152 & 134 & 90 \\
1938 & 2,685 & 75 & 55 & 52 & 125 & 130 & 87 \\
\midrule
1940 & 2,771 & 86 & 57 & 52 & 129 & 159 & --- \\
1950 & 3,667 & 86 & 54 & 66 & 129 & 160 & --- \\
1968 & 7,348 & 91 & 64 & 70 & 141 & 136 & --- \\
1973 & 8,080 & 73 & 60 & 63 & 133 & 125 & --- \\
1980 & 8,785 & 68 & 59 & 61 & 137 & 127 & --- \\
1989 & 8,958 & 53 & 51 & 52 & 126 & 114 & --- \\
\midrule
1995 & 6,277 & 39 & 35 & 33 & 145 & 121 & 125 \\
2022 & 18,863 & 74 & 64 & 68 & 155 & 97 & 110 \\
\bottomrule
\end{tabular}
\begin{tablenotes}[para,flushleft]
\footnotesize
\textit{Sources:} Estonia: see online appendix. Lithuania: \cite{klimantas2026soviet}; Latvia: \cite{Klimantas2024} (1940--89 not shown). Finland, Sweden, Western core and USSR: \cite{mpd2023} and online appendix. \textit{Notes:} USSR within 1990 borders; 1995 and 2022 refer to the former Soviet Union aggregate.
\end{tablenotes}
\end{threeparttable}
\end{table}


\subsection{Interwar Estonia: recovery, depression and economic nationalism}

The interwar record falls into two phases. The 1920s were a recovery: income per head grew by \gr{1919}{1929} per cent a year but regained its 1913 level only in 1929, and then by barely 3 per cent. The sectors that led, industry, construction and transport, were those that had been most damaged, and the limits of the recovery were set by the permanent loss of the imperial market. Pre-war Estonian industry had been built to serve Russia; the Bolshevik revolution closed that market, and the new republic's attempts to recover a Russian transit trade through Tallinn in the early 1920s came to little.\footnote{\autocite[193--210]{koll1994gap}.} Finland, which had never depended on the Russian market to the same degree and whose forest industries found Western buyers at once, grew by \grx{Finland}{1919}{1929} per cent a year over the same decade (Table~\ref{tab:growth}). In the 1920s, therefore, Estonia's social capability was not in doubt, but the technology it had absorbed before 1914 was no longer congruent with the markets available to it.

The 1930s look different. The Depression was shallow, 8 per cent from peak to trough, and short, and the recovery after 1933 was among the fastest in Europe: Estonian income per head grew by \gr{1929}{1938} per cent a year between 1929 and 1938, when the Western core managed \grx{Western}{1929}{1938} per cent and East-Central Europe \grx{ECE}{1929}{1938} per cent (Table~\ref{tab:growth}). The timing follows the exchange rate. Estonia clung to gold through the banking crisis of 1931--2, suffered the deflation that gold imposed, and recovered only after it devalued and pegged to sterling in June 1933, a pattern consistent with that which Eichengreen and Sachs identified for Europe as a whole and that the Nordic economies displayed.\footnote{\autocite[74]{bernankejames01}; \autocite{RefWorks:doc:5ea847d2e4b0aa025739fe41}; \autocite{grytten_nordicdepression}; \autocite[287--99]{kuokstis2016}.} The content of the recovery was the policy that Kõll and Valge have called economic nationalism: state investment in oil shale, textiles and food processing, protection of the home market and the organisation of agricultural exports through co-operatives in imitation of Denmark.\footnote{\autocite[11--29]{koll1998}; \autocite{norkus_denmarkisation}; \autocite[118--35]{klesment2005}.} The sectoral indices confirm that the 1930s acceleration was industrial and urban: construction grew by \sgr{Construction}{1929}{1938} per cent a year, industry by \sgr{Industrytotal}{1929}{1938} per cent and the housing stock by \sgr{Housing}{1929}{1938} per cent, while agriculture grew by \sgr{Primarysectortotal}{1929}{1938} per cent.

Table~\ref{tab:interwar} places the result in its regional setting. In 1938 Estonia's income per head, \ee{1938}~\intlnineteen{}, was above that of Lithuania, Poland and Hungary, slightly below Czechoslovakia's and well below Latvia's, Finland's and the Nordic average. Its growth over 1929--38 was exceeded in the region only by its two Baltic neighbours, by Finland and by the Soviet Union, whose own figures for the period are the subject of the debate reviewed in Section~\ref{sec:method}, and was far above that of the Western core and of East-Central Europe. Table~\ref{tab:structure} shows why the result should be read as a success and why it was not a greater one. Estonia's adult literacy, around 92 per cent, was second only to Latvia's in the European periphery and above Finland's; its university enrolment per head was among the highest in Europe; its land was in the hands of smallholders after the reform of 1919. These are the ingredients of social capability. Yet 63 per cent of the labour force still worked in agriculture in 1934, a share as high as Finland's, though below Latvia's, and agriculture produced some 38 per cent of GDP in 1937 (Table~\ref{tab:structure}). An economy of that structure, open to trade (Table~\ref{tab:trade}) but exporting butter, bacon, timber and flax, could grow at the pace of its agricultural productivity and of the markets Britain and Germany allowed it, and no faster; the industrial drive of 1934--9 was a response to that constraint, not an escape from it.

% tab:interwar
\begin{table}[htbp]
\centering\small
\caption[GDP per capita in Estonia and European comparators, 1913--38]{\textsc{GDP per capita in Estonia and European comparators, 1913--38}}
\label{tab:interwar}
\begin{threeparttable}
\begin{tabular}{lrrrrrrrr}
\toprule
& \multicolumn{5}{c}{\intlnineteen{}} & \multicolumn{3}{c}{Growth, \% p.a.} \\
\cmidrule(lr){2-6}\cmidrule(lr){7-9}
Country & 1913 & 1919 & 1929 & 1932 & 1938 & 1919--29 & 1929--38 & 1913--38 \\
\midrule
Denmark & 3,912 & 3,860 & 5,075 & 5,169 & 5,762 & 2.8 & 1.4 & 1.6 \\
Sweden & 2,874 & 2,825 & 4,063 & 3,968 & 4,901 & 3.7 & 2.1 & 2.2 \\
\textbf{Finland} & 2,111 & 1,658 & 2,717 & 2,550 & 3,589 & 5.1 & 3.1 & 2.1 \\
Austria & 3,465 & 2,259 & 3,699 & 2,940 & 3,559 & 5.1 & $-$0.4 & 0.1 \\
Latvia & 2,361 & 1,481 & 2,351 & 2,286 & 3,099 & 4.7 & 3.1 & 1.1 \\
Ireland & 2,736 & 2,595 & 2,824 & 2,885 & 3,052 & 0.8 & 0.9 & 0.4 \\
Czechoslovakia & 2,096 & 1,588 & 3,042 & 2,680 & 2,803 & 6.7 & $-$0.9 & 1.2 \\
\textbf{Estonia} & 2,058 & 1,600 & 2,113 & 2,105 & 2,685 & 2.8 & 2.7 & 1.1 \\
Greece & 1,455 & 1,274 & 2,342 & 2,289 & 2,677 & 6.3 & 1.5 & 2.5 \\
Hungary & 2,098 & 1,484 & 2,476 & 2,192 & 2,655 & 5.3 & 0.8 & 0.9 \\
Poland & 1,739 & 1,348 & 2,117 & 1,658 & 2,182 & 4.6 & 0.3 & 0.9 \\
USSR & 1,414 & 573 & 1,386 & 1,439 & 2,150 & 9.2 & 5.0 & 1.7 \\
Lithuania & 1,632 & 1,017 & 1,577 & 1,702 & 2,057 & 4.5 & 3.0 & 0.9 \\
\bottomrule
\end{tabular}
\begin{tablenotes}[para,flushleft]
\footnotesize
\textit{Sources:} Estonia: see online appendix. Latvia: 1913, \cite[603]{norkusmark2021}; 1919--38, \cite[316]{klimantas2025thesis}, table E.1, the series published for 1920--39 in \cite{Klimantas2024}, whose values for 1929 and 1932 (2,349 and 2,288) differ from those used here by less than 0.1 per cent. Lithuania: \cite{klimantas2023lithuania}. Other countries: \cite{mpd2023}. \textit{Notes:} countries ordered by 1938 level.
\end{tablenotes}
\end{threeparttable}
\end{table}


% tab:trade
\begin{table}[htbp]
\centering\small
\caption[Trade openness (exports plus imports as a percentage of GDP), period averages]{\textsc{Trade openness (exports plus imports as a percentage of GDP), period averages}}
\label{tab:trade}
\begin{threeparttable}
\begin{tabular}{lrrr}
\toprule
Country & 1924--29 & 1930--34 & 1935--38 \\
\midrule
\textbf{Estonia} & 33.4 & 25.6 & 31.5 \\
Latvia & 38.1 & 26.9 & 21.8 \\
Lithuania & 19.9 & 21.8 & 21.6 \\
\textbf{Finland} & 41.2 & 26.8 & 34.3 \\
Denmark & 51.8 & 35.2 & 36.3 \\
Norway & 56.1 & 35.6 & 41.1 \\
Sweden & 36.8 & 23.9 & 29.3 \\
Poland & 13.3 & 10.3 & 11.8 \\
Czechoslovakia & 36.8 & 23.6 & 26.5 \\
Hungary & 21.7 & 14.1 & 16.8 \\
Germany & 25.0 & 16.9 & 15.6 \\
United Kingdom & 36.9 & 23.3 & 25.6 \\
\bottomrule
\end{tabular}
\begin{tablenotes}[para,flushleft]
\footnotesize
\textit{Sources:} see the online appendix for the compilation of trade and GDP series. Baltic exports and imports: \cite{lidaexportsee}; \cite{lidaimportsee}; \cite{lidaexportslv}; \cite{lidaimportslv}; \cite{lidaexp}; \cite{lidaimp}; \cite[123, 161]{balticstates1938}. Other countries: \cite{worldtrade1949}. \textit{Notes:} unweighted averages of annual ratios; missing years omitted.
\end{tablenotes}
\end{threeparttable}
\end{table}


% tab:structure
\begin{table}[htbp]
\centering\footnotesize\setlength{\tabcolsep}{4pt}
\caption[Social capability and economic structure around 1930: Estonia, Finland and European comparators]{\textsc{Social capability and economic structure around 1930: Estonia, Finland and European comparators}}
\label{tab:structure}
\begin{threeparttable}
\begin{tabular}{lrrrr}
\toprule
Country & Literacy & Agriculture & University & Agriculture \\
& (\% of adults) & (\% of employment) & students (per 1,000) & (\% of GDP) \\
& c.\,1930 & c.\,1930 & 1930 & c.\,1937 \\
\midrule
Sweden & 99.6 & 36.0 & 1.6 & 13.0 \\
United Kingdom & 97.5 & 6.0 & 3.1 & 4.0 \\
Germany & 97.5 & 29.0 & --- & 16.0 \\
Denmark & 97.5 & 35.3 & 2.7 & 18.0 \\
Norway & 97.5 & 35.5 & 2.2 & 13.0 \\
Austria & 97.5 & 39.0 & 3.2 & 14.0 \\
Ireland & 97.5 & 48.5 & 1.5 & 24.8 \\
Czechoslovakia & 96.0 & 37.5 & --- & 24.0 \\
Latvia & 93.4 & 67.2 & 4.9 & 39.6 \\
\textbf{Estonia} & 92.0 & 63.0 & 3.5 & 37.6 \\
Hungary & 90.0 & 54.2 & --- & 31.0 \\
\textbf{Finland} & 84.1 & 64.6 & 2.0 & 32.0 \\
Italy & 78.4 & 48.2 & 1.1 & 29.0 \\
Lithuania & 77.0 & 75.2 & 1.8 & 43.8 \\
Poland & 73.0 & 65.9 & --- & --- \\
Bulgaria & 69.0 & 71.8 & --- & 52.1 \\
Spain & 68.0 & 51.9 & 1.7 & 22.0 \\
Greece & 58.1 & 56.6 & 1.1 & 54.4 \\
Romania & 57.1 & 77.2 & --- & --- \\
Yugoslavia & 55.0 & 79.7 & --- & 41.0 \\
USSR & 50.0 & 73.0 & --- & 29.0 \\
Portugal & 39.8 & 48.0 & 0.8 & 30.0 \\
\bottomrule
\end{tabular}
\begin{tablenotes}[para,flushleft]
\footnotesize
\textit{Sources:} literacy: \cite{cronin_literacy}; \cite[246--77]{kirk_literacy}; \cite[528]{bostonencikl_rast}. Agricultural employment: \cite[181]{morysinterwarinmorys2021}; \cite{mitchelstat}; \cite[395, 409, 418]{sov_labour1939}. University students and agricultural shares of GDP: see online appendix, section A.7. \textit{Notes:} countries ordered by literacy; a literacy value of 97.5 denotes near-universal literacy where the source gives no census figure. `---': not available.
\end{tablenotes}
\end{threeparttable}
\end{table}


\subsection{Under Soviet rule: collectivisation, industrialisation and stagnation}

The Soviet period opens with a decade of destruction. Income per head fell by about a third between 1940 and 1944 on the present estimates, and the population by a fifth between 1939 and 1945 through war deaths, deportation and flight.\footnote{\autocite[265]{taagepera_coll01}; \autocite[12]{deportations_ee}.} The recovery of industry and construction after 1945 was rapid, as Soviet reconstruction everywhere was: both regained their 1938 level by 1948, and industry grew by around 20 and construction by around 14 per cent a year between 1945 and 1953 (Tables~\ref{tab:sectoral} and~\ref{tab:appsec1}). Agriculture tells a different story, and one that the sectoral indices allow to be told with some precision. Livestock output, halved by the war, was still 38 per cent below its 1938 level in 1950, the year after collectivisation; crop output recovered sooner, but the agricultural sector as a whole stood at 60--65 per cent of its 1938 level in 1945--6, fluctuated between 77 and 93 per cent of it in 1947--57, and exceeded it only in \sregain{Primarysectortotal}. The index shows a shallower trough than Mertelsmann's archival estimate, which puts official harvests 10--15 per cent above reality and output after collectivisation at perhaps 40 per cent of the 1939 level, but the two agree that the sector did not recover for a decade; Mertelsmann also finds that the collective and state farms produced little more than a third of agricultural output in the early 1950s, the rest coming from the private plots of their members, which still supplied 28 per cent of output in 1960.\footnote{\autocite[156--8]{olaf_stalin}; \autocite[79--109]{feest2014}; \autocite[257]{klesment2009}; \autocite[160]{hansen1994}.} Nothing comparable to the Soviet famine of 1932--3 occurred, probably because the Baltic collectivisation of 1949 did not coincide with a procurement crisis and because, as Harrison has argued for Soviet agriculture generally, the system had by then learned to extract the surplus without destroying the herds.\footnote{\autocite[192--208]{harrison1996}.} The cost was a lost decade and a permanent shift of effort towards animal husbandry, in which Estonia came to lead the Union in output per head, while crop output peaked in \speak{Crops} and never grew again.\footnote{\autocite[234]{years_of_dependence}.}

Between 1950 and 1968 Estonia experienced what Gerschenkron would have recognised as a classic burst of state-led industrialisation. Total GDP grew by \tgr{1950}{1968} per cent a year and income per head by \gr{1950}{1968} per cent, faster than in Finland, the Nordic countries or the Western core (Table~\ref{tab:growth}). The engine was investment directed from Moscow into three complexes. The first was the oil-shale basin of the north-east, which fed two of the largest thermal power stations in the Union, completed in 1965 and 1973; electricity output rose from 2 billion kilowatt-hours in 1960 to 12 billion in 1970, of which Estonia itself consumed a third. The second was cotton textiles, which Moscow expanded in Estonia rather than in cotton-growing Uzbekistan because Estonia had the skilled labour and the mills. The third was the machine-building and electronics plants that supplied the Union market.\footnote{\autocite[239]{years_of_dependence}; \autocite[41--7]{ahderajasalu1992}; \autocite[248--58]{pihlamagi2013}.} The sovnarkhoz reform of 1957--65, which briefly placed industrial management in Tallinn, saw industrial output double in seven years, with 70 per cent of investment going to industry and the share of light industry in the total deliberately reduced.\footnote{\autocite[248--52]{pihlamagi2013}; \autocite[67--70]{purre1965}.} Heavy industry grew by \sgr{Heavyindustries}{1950}{1968} per cent a year, textiles by \sgr{Textiles}{1950}{1968} per cent and construction by \sgr{Construction}{1950}{1968} per cent. The labour came partly from a countryside that collectivisation had released and partly from immigration: Estonia's population grew by half between 1950 and 1989 while the Estonian share of it fell from more than 90 to little more than 60 per cent.\footnote{\autocite[158--217]{case_estonia}.}

How did this growth compare with that of the Union into which Estonia had been absorbed? Table~\ref{tab:relative} gives the answer, and it is not the one that either Soviet or nationalist accounts have assumed. Estonia entered the Union in 1940 with income per head \pctussr{1940} per cent of the Soviet average; it stood at \pctussr{1968} per cent in 1968 and at \pctussr{1989} per cent in 1989. The gap neither closed nor widened appreciably over half a century: Estonia grew with the Union, somewhat faster in the 1950s and 1960s and somewhat slower thereafter, and remained one of its two or three richest republics throughout. The result matches the Soviet-era estimates of relative republican income, which placed Estonia at 137--145 per cent of the Union average in 1960--70 by net material product and at 134 per cent in 1970 on Koropeckyj's reworking, and it matches the CIA's finding that Estonian national income per head grew by 6.3 per cent a year in 1961--9 against 5.5 per cent for the Union.\footnote{\autocite[169, 174]{schroeder1973}; \autocite[64--71]{bs_livingstandards1976}; \autocite[8--10]{cia1972}; \autocite[35--8]{bs_livingstandards1988}.} It also bears on the question whether Moscow bought the loyalty of its western borderlands with investment. Investment per head in Estonia was above the Union average in the 1960s, second only to Kazakhstan, but the republic's share of investment in its own national income was below the Union average. Van Selm and Dölle's reconstruction of inter-republican transfers finds Estonia's net balance close to zero from 1966 to 1978, becoming a net recipient of 4--6 per cent of national income only after 1979 and of 13 per cent in 1988, when the Union economy was already unravelling.\footnote{\autocite[22--9]{cia1972}; \autocite[53]{adirim1988}; \autocite[136--40]{vanselmdolle1993}.} For the years after 1966 for which transfers can be measured, Estonia's high income under Soviet rule was earned rather than transferred; and over the period as a whole it was inherited rather than created, since the republic was richer than the Union when it was annexed and remained so in roughly the same proportion.

The deceleration after 1968 was general and, by the late 1970s, complete. Income per head grew by \gr{1968}{1989} per cent a year over 1968--89 and by \gr{1973}{1989} per cent after 1973, and the sectoral peaks cluster in a single decade: food processing in \speak{Food}, construction in \speak{Construction}, electricity and oil-shale output in 1979--80, heavy industry in \speak{Heavyindustries}.\footnote{Electricity peaked in 1979 and fell by 9 per cent to 1985; oil-shale mining levelled off after the arrival of natural gas (\autocite[239]{years_of_dependence}; \autocite{taagepera2018}).} The sources point to three causes. The Union-wide slowdown in investment growth that the tenth and eleventh five-year plans deliberately imposed fell on a republic whose growth had been investment-led, and whose own labour reserves were exhausted: contemporaries noted that further industrial expansion in Estonia required either immigration or capital deepening, and the planners increasingly withheld both.\footnote{\autocite[128--63]{hanson2003}; \autocite[164--76]{kolde1978}; \autocite[441]{dienes1972}.} The oil-shale complex, built on a resource that was expensive to mine and polluting to burn, lost its rationale once Siberian gas reached the Baltic.\footnote{\autocite[188--9]{olaf_stalin}.} And the gains in reported output after 1975 were increasingly nominal: Finnish and Estonian economists who compared the two economies in 1992 concluded that the rapid growth of Estonian industry to 1975 reflected real investment, but that `the further growth of production can be explained mainly by the change in prices'.\footnote{\autocite[39--40]{ahderajasalu1992}.} The physical indices used here confirm that judgement: industrial output grew by \sgr{Industrytotal}{1968}{1989} per cent a year after 1968, and the branches that continued to grow, textiles and other light industry, did so slowly, while cotton cloth output was flat from 1970.\footnote{\autocite[41]{ahderajasalu1992}.} Two further irregularities in the indices deserve mention. Construction, which had grown by about 9 per cent a year since 1950, peaked in \speak{Construction} and stagnated for the rest of the Soviet period as the Union-wide investment squeeze took hold, and crop output fell by a quarter in each of the poor harvests of 1978 and 1981 that affected the Union as a whole; and the housing stock, which had grown by 2 per cent a year since the 1950s, stopped growing in 1986 and declined thereafter, as completions fell from 13,500 dwellings a year in 1985 to 5,200 in 1991.\footnote{\autocite[134--50]{hanson2003}; \autocite[60, 66]{ahderajasalu1992}; \autocite[216]{wb1993estonia}.}

Read through the lens of growth theory, the Soviet episode is coherent. The catch-up of 1950--68 was what the models predict when a capital-poor economy with high social capability is made to invest: Estonia's literate and urbanised labour force absorbed Soviet technology quickly, and the mass-production methods of mid-century were congruent with its endowments. The stagnation that followed is equally what the models predict when the institutional condition fails. Extensive growth exhausted its inputs, and the system could not replace them with the productivity growth that flexible production, foreign technology and competition were delivering across the Gulf: the share of services in Estonian GDP was 42 per cent in 1990, barely above its interwar level, against about 60 per cent in Finland.\footnote{Table~\ref{tab:benchmarks}; \autocite[11--12]{ahderajasalu1992}; \autocite{klimantas2026soviet}.} Estonia's income relative to the Western core, which had risen from \pctwc{1950} per cent in 1950 to \pctwc{1968} per cent in 1968, fell back to \pctwc{1989} per cent by 1989, almost exactly its level of 1913. Two decades of Soviet growth had been cancelled by two decades of Soviet stagnation.

\subsection{Transition and recovery}

The transition compressed into four years a fall in income per head of \eeTROUGHpct{} per cent, deeper than the Depression and comparable only to the two world wars. The sectoral indices to 1991 show the collapse beginning in Soviet-era industry and construction, which contracted at \sgr{Industrytotal}{1989}{1991} and \sgr{Construction}{1989}{1991} per cent a year even before the rouble was abandoned. After 1991 the reconstruction follows the modern series, which records the loss of eastern markets and cheap Soviet inputs, the currency board of 1992 and the stabilisation that followed.\footnote{\autocite[1--8]{wb1993estonia}; \autocite[164--6]{norkus_restorations}; \autocite[482--509]{shteinbuka1993}.} By comparison with the other Soviet republics the fall was contained, probably because Estonia had less heavy industry to lose and reoriented trade towards Finland and Sweden faster than any other successor state.\footnote{\autocite[22]{transition_greenes}; \autocite[443--4]{transition_progress}.} Recovery from 1995 at \gr{1995}{2022} per cent a year carried Estonia from \pctwc{1995} per cent of the Western core in 1995 to \pctwc{2022} per cent in 2022, close to the \pctwc{1968} per cent reached in 1968 at the height of Soviet growth. Whether the catch-up of the market era will continue or stall in the middle-income trap that Staehr has described is a question for another paper; what the long series shows is that Estonia converged on Western Europe under capitalism twice, in the 1930s and after 1995, and under planning once, briefly, before diverging again.\footnote{\autocite[274--80]{staehr2015}.}

% =============================================================
% VI. A FINNISH MIRACLE THAT NEVER WAS
% =============================================================
\section{A Finnish miracle that never was}
\label{sec:finland}

The comparison with Finland has shaped the way Estonians and their historians think about the twentieth century. The two peoples speak related languages, were shaped by Lutheran schooling and Swedish law, emancipated their peasantries early and built co-operative agrarian economies after 1918; in the Soviet period Finland, visible on television from Tallinn, became `proof of what Estonia would have been if the country had not ended up under the Soviet occupation'.\footnote{\autocite[83]{lepppantti2013}.} The scholarly version of the argument has two steps. The first is that the two countries stood at the same level before 1940: Dellenbrant wrote that their levels of development were `close to equal around 1940', Soviet-era comparisons of real wages put Estonian workers ahead of Finnish ones as late as 1966, and the White Book's estimate of the losses of occupation rests on the assumption that Estonia would have followed the Finnish path.\footnote{\autocite[206]{dellenbrant1992}; \autocite[251, 257]{klesment2009}; \autocite[141--74]{kukk2005}.} The second is that the gap observed in 1991, when Estonian living standards were put at 40 per cent of Finland's, measures the cost of Soviet rule.\footnote{\autocite[206]{dellenbrant1992}.} Sceptics have questioned both steps. Kõll argued in 1994 that the belief in Baltic--Nordic parity before the war was `misleading'; an Estonian economist of the 1930s had already written that `Denmark, Sweden, Finland, and even Latvia are ahead of us regarding their wealth'; and Norkus, working from the benchmark estimates, concluded that the gap `cannot only be considered as the legacy of the Soviet occupation'.\footnote{\autocite[201]{koll1994gap}; \autocite[131]{klesment2005}; \autocite[18]{norkus2016}; \autocite[154]{norkus_restorations}.} The new series allows the two steps to be examined year by year.

Table~\ref{tab:relative} gives Estonian income per head as a percentage of Finland's at the benchmark years, and Table~\ref{tab:growth} the growth rates between them. In 1913 the two countries were indeed close: Estonia stood at \pctfin{1913} per cent of the Finnish level. In 1919, after wars that had damaged both, it stood at \pctfin{1919} per cent. The parity did not survive the 1920s. Finland grew by \grx{Finland}{1919}{1929} per cent a year to 1929 and Estonia by \gr{1919}{1929} per cent, and by 1929 Estonia was at \pctfin{1929} per cent of the Finnish level; the 1930s, in which Estonia's growth of \gr{1929}{1938} per cent a year was close to Finland's \grx{Finland}{1929}{1938} per cent, left the ratio at \pctfin{1938} per cent in 1938. Clark's 90 per cent, which Valge adopted and which has anchored the literature since, is thus too high by a sixth; Bocharnikova's finding that the purchasing power of wages in Estonia was only 4 per cent below Finland's in 1938 refers to urban wages alone and says nothing about the nearly two-thirds of the Estonian labour force on the land.\footnote{\autocite[81--2]{bocharnikova2021}.} The difference in 1938 was not one of social capability. Table~\ref{tab:structure} shows Estonian literacy above Finnish, Estonian university enrolment per head well above Finnish, and the same share of the labour force in agriculture; Table~\ref{tab:trade} shows both economies open, though Finland more so. The difference is better explained by resources and markets. Finland's forests, sawmills and pulp mills gave it an export staple for which Western demand recovered quickly after 1920 and which carried a growing engineering sector with it; Estonia's butter and bacon faced Danish competition and British quotas, and its industry had lost the only large market it had ever served.\footnote{\autocite[184, 198--200]{finlandeh01}; \autocite[193--210]{koll1994gap}.}

The second step of the argument fares no better, though for a different reason. Soviet rule did open a gap, but the timing is not what the argument supposes. Finland lost the Second World War and paid reparations; Estonia was fought over twice and collectivised. In 1950 the ratio stood at \pctfin{1950} per cent, higher than in 1938, and in 1968, after the two decades of Soviet industrialisation examined above, at \pctfin{1968} per cent, the highest level since 1919. Measured against Finland, the first half of Soviet rule was a period of convergence. The divergence that the literature attributes to occupation is a phenomenon of its second half: between 1968 and 1989 Finland grew by \grx{Finland}{1968}{1989} per cent a year, half as fast again as the Western core, and Estonia by \gr{1968}{1989} per cent, so that the ratio fell to \pctfin{1973} per cent in 1973 and \pctfin{1989} per cent in 1989. Those who lived through it dated the divergence correctly: Dellenbrant, drawing on Kukk's calculations, placed its beginning in the late 1960s.\footnote{\autocite[205]{dellenbrant1992}.} What the new series adds is the magnitude. Expressed in logarithms, the Estonian shortfall against Finland was 0.29 in 1938 and 0.63 in 1989; almost half of the gap at the end of Soviet rule had opened before it began, and the whole of the increase after 1938 occurred after 1968.

Two conclusions follow. The first concerns the use of Finland as a yardstick. Finland is a revealing comparator for Estonia, because it isolates the effect of institutions on two economies with similar endowments of human capital. But it is not a measure of what Estonia lost. The two economies were already a quarter apart in 1938, and Finland's growth after 1968, \grx{Finland}{1968}{1989} per cent a year, was exceptional by the standards of the Western core (\grx{Western}{1968}{1989} per cent) and of the Nordic countries (\grx{Nordics}{1968}{1989} per cent), driven by a forest-products and engineering complex that Estonia never possessed and by access to a Soviet market on favourable bilateral terms that Estonia, inside the Union, could not exploit.\footnote{\autocite[184--200]{finlandeh01}; \autocite[23--4]{norkus2016}.} The Western periphery as a whole, the group of poorer Western economies to which Estonia belonged by income in 1938, grew by \grx{Periphery}{1968}{1989} per cent a year over the same years (Table~\ref{tab:growth}); its Iberian and Greek members, starting from levels below Estonia's of 1950, grew as fast as Finland, while Ireland grew more slowly. Measured against the core or the periphery rather than against Finland, Estonia's shortfall in 1968--89 is between one and a half and two points a year rather than more than two and a half: still large, but a different magnitude from the one the White Book implies, and one that cannot be read off the Finnish series alone. The second concerns the Soviet record itself. The years in which Estonia kept pace with Finland were the years of forced industrialisation, high investment and immigration; the years in which it fell behind were those in which the planned economy, having exhausted the sources of extensive growth, could neither raise productivity nor change its structure. The gap with Finland in 1991 is therefore likely to overstate the cost of occupation. The part of it that can be charged to the system is the stagnation of the last two decades, and the measure of that cost is Estonia's shortfall over those years against the Western economies of its own level rather than against Finland: a less comforting conclusion than the one the White Book drew, and a more exact one.

% =============================================================
% VII. CONCLUSION
% =============================================================
\section{Conclusion}
\label{sec:conclusion}

This paper has constructed the first annual, internationally comparable series of Estonian GDP per capita for 1919--95 and joined it to existing estimates to give a continuous series for 1913 and 1919--2022 in 1990 and 2011 international dollars. It rests on a new 1938 benchmark from a direct price comparison with Sweden, a recalculated production-side series for 1919--40, and a new series for the ESSR for 1937--95 built from physical output, transport and employment data and anchored to Bolotin's 1990 benchmark. Built from different sources, with different weights and from benchmarks half a century apart, the two new series agree within 4 per cent in every year of their 1937--40 overlap, and the Soviet-period series meets a modern series derived from the 2011 ICP round within 1 per cent at the 1991 splice. The series supersedes the three partial estimates that preceded it, lowering late-Soviet income by about a quarter relative to the Maddison Project's NMP-based figures and raising it by a similar margin relative to the only earlier long series, and it is ready for the Maddison Project Database.

Three findings stand out. First, Estonia's interwar growth was respectable but not exceptional: \gr{1919}{1929} per cent a year in the 1920s, as the country recovered from war and the loss of its imperial market, and \gr{1929}{1938} per cent in the 1930s, when devaluation and a state-led industrial drive produced a rapid recovery from the Depression. Second, Soviet rule brought two decades of rapid growth and two of stagnation: \gr{1950}{1968} per cent a year in 1950--68, faster than in Finland or the Western core, as investment directed from Moscow built an oil-shale, textile and engineering economy on an unusually literate labour force, and \gr{1968}{1989} per cent in 1968--89, as the sources of extensive growth ran out. Estonia's income relative to the Western core rose from two-thirds in 1950 to seven-tenths in 1968 and fell back to half by 1989, while within the Union it remained what it had been at annexation, a republic between a fifth and a half richer than the average. Third, the transition cut income per head by a third in four years, and recovery after 1995 carried Estonia back, by 2022, close to the share of Western European income it had briefly reached in 1968.

The comparison with Finland, which gives the paper its title, needs restating. The two countries were close in 1913 and 1919, but by 1938 Estonia stood at three-quarters of the Finnish level, and almost half of the gap of 1989 had opened under independence, not occupation; the further divergence was concentrated after 1968, Estonia having kept pace with Finland in the preceding two decades. Finland shows what institutions can do with a given endowment of human capital, but not what Estonia would have become, because the Finnish path depended on resources and markets that Estonia did not have. The costs of Soviet rule were real, and the series locates most of them in the stagnation of 1968--89; measured against the Western core or periphery rather than against a Finland whose growth in those decades was exceptional, they are smaller than the Finnish comparison implies, and they are not the costs of a Finnish miracle that Estonia was denied.

Several extensions are open: a dedicated study of the 1938 service sector of the kind now available for Lithuania and Latvia; a Soviet-period benchmark built from the Estonian archives with republican rather than Union-wide prices, which would test the level to which the volume indices are anchored; expenditure-side cross-checks for a few interwar years; and the application of the procedure to the other non-Russian republics, whose national accounts remain the largest gap in the quantitative economic history of the Soviet Union.


% =============================================================
% BIBLIOGRAPHY
% =============================================================
\newpage
{\sloppy\printbibliography}

% =============================================================
% STATISTICAL APPENDIX
% =============================================================
\newpage
% =============================================================
% STATISTICAL APPENDIX
% =============================================================
\section*{Statistical appendix}
\addcontentsline{toc}{section}{Statistical appendix}
\setcounter{table}{0}
\renewcommand{\thetable}{A\arabic{table}}

This appendix reports the annual series on which the paper rests. Table~\ref{tab:series} gives GDP per capita in 1990 and 2011 international dollars and the mid-year population for 1913 and every year from 1919 to 2022; Tables~\ref{tab:appsec1} and~\ref{tab:appsec2} give the constant-price volume indices of sectoral value-added for 1919--95 from which the GDP series of 1919--91 is built, with 1938 equal to 100. The indices for 1992--5 are those of the Soviet-period series and are superseded, in the final series, by the modern series from 1991 (Section~\ref{sec:limitations}). The online appendix, available as an Excel workbook, contains the same series together with the 1938 and 1990 benchmarks by sector, the overlapping estimates at the two splicing points, the prices and weights of the 1938 comparison with Sweden, the product-by-product composition of the industrial index, and a source list for every series by year, keyed to the sections A.1--A.6 referred to in Table~\ref{tab:sources}. The interwar series refer to the borders of the republic and the series from 1945 to present-day borders (Section~\ref{sec:method}).

% tab:series -- appendix table A1
\begingroup\setlength{\tabcolsep}{3pt}\footnotesize
\begin{longtable}{rrrr@{\hspace{1.2em}}rrrr}
\caption[Estonian GDP per capita and population, 1913 and 1919--2022]{\textsc{Estonian GDP per capita and population, 1913 and 1919--2022}}\label{tab:series}\\
\toprule
Year & \intlnineteen{} & \intltwenty{} & Population & Year & \intlnineteen{} & \intltwenty{} & Population \\
\midrule
\endfirsthead
\multicolumn{8}{l}{\textit{Table~\ref{tab:series} (continued)}}\\
\toprule
Year & \intlnineteen{} & \intltwenty{} & Population & Year & \intlnineteen{} & \intltwenty{} & Population \\
\midrule
\endhead
\midrule\multicolumn{8}{r}{\textit{continued}}\\
\endfoot
\bottomrule
\multicolumn{8}{p{0.96\linewidth}}{\footnotesize\textit{Sources:} see online appendix. GDP per capita in international Geary--Khamis dollars of 1990 and of 2011 (\intltwenty{} = \intlnineteen{} $\times$ 1.594, \cite[12]{bolt_zanden_newgdp2021}); population in thousands, mid-year, within present-day borders. No estimates for 1914--18.}\\
\endlastfoot
1913 & 2,058 & 3,280 & 1,164.0 & 1971 & 8,002 & 12,755 & 1,377.0 \\
\midrule
1919 & 1,600 & 2,550 & 1,064.0 & 1972 & 7,896 & 12,587 & 1,392.5 \\
1920 & 1,475 & 2,351 & 1,067.8 & 1973 & 8,080 & 12,879 & 1,406.0 \\
1921 & 1,673 & 2,666 & 1,087.1 & 1974 & 8,217 & 13,098 & 1,418.2 \\
1922 & 1,756 & 2,800 & 1,102.4 & 1975 & 8,419 & 13,419 & 1,429.4 \\
1923 & 1,913 & 3,050 & 1,110.8 & 1976 & 8,623 & 13,745 & 1,439.6 \\
1924 & 1,895 & 3,020 & 1,115.6 & 1977 & 8,601 & 13,709 & 1,450.2 \\
1925 & 2,038 & 3,249 & 1,117.0 & 1978 & 8,474 & 13,508 & 1,460.2 \\
1926 & 2,030 & 3,235 & 1,116.8 & 1979 & 8,523 & 13,585 & 1,468.3 \\
1927 & 2,043 & 3,256 & 1,115.6 & 1980 & 8,785 & 14,003 & 1,477.2 \\
1928 & 2,046 & 3,262 & 1,115.7 & 1981 & 8,585 & 13,684 & 1,487.7 \\
1929 & 2,113 & 3,368 & 1,115.7 & 1982 & 8,554 & 13,635 & 1,498.4 \\
1930 & 2,282 & 3,638 & 1,116.1 & 1983 & 8,775 & 13,988 & 1,508.7 \\
1931 & 2,176 & 3,469 & 1,118.4 & 1984 & 8,932 & 14,238 & 1,518.6 \\
1932 & 2,105 & 3,356 & 1,121.5 & 1985 & 8,728 & 13,913 & 1,528.8 \\
1933 & 2,186 & 3,485 & 1,124.3 & 1986 & 8,919 & 14,216 & 1,540.2 \\
1934 & 2,353 & 3,751 & 1,126.4 & 1987 & 8,833 & 14,079 & 1,552.2 \\
1935 & 2,355 & 3,754 & 1,128.9 & 1988 & 8,879 & 14,153 & 1,561.9 \\
1936 & 2,492 & 3,972 & 1,130.0 & 1989 & 8,958 & 14,280 & 1,568.1 \\
1937 & 2,722 & 4,338 & 1,130.7 & 1990 & 8,425 & 13,429 & 1,569.2 \\
1938 & 2,685 & 4,280 & 1,132.5 & 1991 & 7,532 & 12,006 & 1,561.3 \\
1939 & 2,736 & 4,361 & 1,140.4 & 1992 & 6,069 & 9,675 & 1,533.1 \\
1940 & 2,771 & 4,418 & 1,122.7 & 1993 & 5,870 & 9,357 & 1,494.1 \\
1941 & 2,390 & 3,809 & 1,049.1 & 1994 & 5,899 & 9,402 & 1,462.5 \\
1942 & 2,382 & 3,797 & 1,007.0 & 1995 & 6,277 & 10,005 & 1,436.6 \\
1943 & 2,328 & 3,710 & 981.9 & 1996 & 6,685 & 10,655 & 1,415.6 \\
1944 & 1,909 & 3,044 & 917.8 & 1997 & 7,644 & 12,184 & 1,399.5 \\
1945 & 2,656 & 4,234 & 903.1 & 1998 & 8,053 & 12,836 & 1,386.2 \\
1946 & 2,833 & 4,516 & 938.7 & 1999 & 7,995 & 12,744 & 1,390.2 \\
1947 & 3,123 & 4,978 & 990.2 & 2000 & 8,759 & 13,961 & 1,397.0 \\
1948 & 3,323 & 5,297 & 1,026.3 & 2001 & 9,344 & 14,894 & 1,388.1 \\
1949 & 3,474 & 5,538 & 1,026.1 & 2002 & 10,040 & 16,004 & 1,379.3 \\
1950 & 3,667 & 5,846 & 1,036.4 & 2003 & 10,871 & 17,328 & 1,370.7 \\
1951 & 3,709 & 5,912 & 1,061.6 & 2004 & 11,680 & 18,618 & 1,362.6 \\
1952 & 3,758 & 5,990 & 1,083.1 & 2005 & 12,867 & 20,509 & 1,354.8 \\
1953 & 3,747 & 5,973 & 1,106.5 & 2006 & 14,207 & 22,645 & 1,346.8 \\
1954 & 3,787 & 6,036 & 1,128.9 & 2007 & 15,353 & 24,473 & 1,340.7 \\
1955 & 3,951 & 6,298 & 1,144.2 & 2008 & 14,604 & 23,279 & 1,337.1 \\
1956 & 4,243 & 6,763 & 1,157.9 & 2009 & 12,492 & 19,912 & 1,334.5 \\
1957 & 4,460 & 7,109 & 1,171.9 & 2010 & 12,826 & 20,445 & 1,331.5 \\
1958 & 4,727 & 7,535 & 1,187.8 & 2011 & 13,800 & 21,997 & 1,327.4 \\
1959 & 5,010 & 7,986 & 1,201.6 & 2012 & 14,294 & 22,784 & 1,322.7 \\
1960 & 5,319 & 8,478 & 1,211.5 & 2013 & 14,557 & 23,204 & 1,318.0 \\
1961 & 5,513 & 8,788 & 1,225.1 & 2014 & 15,046 & 23,983 & 1,314.7 \\
1962 & 5,547 & 8,841 & 1,241.6 & 2015 & 15,354 & 24,474 & 1,314.0 \\
1963 & 5,840 & 9,309 & 1,258.9 & 2016 & 15,807 & 25,196 & 1,314.4 \\
1964 & 6,150 & 9,803 & 1,277.1 & 2017 & 16,726 & 26,662 & 1,316.1 \\
1965 & 6,466 & 10,308 & 1,294.6 & 2018 & 17,313 & 27,596 & 1,320.6 \\
1966 & 6,613 & 10,541 & 1,308.6 & 2019 & 17,883 & 28,506 & 1,325.1 \\
1967 & 7,013 & 11,178 & 1,318.9 & 2020 & 17,730 & 28,262 & 1,329.5 \\
1968 & 7,348 & 11,712 & 1,331.2 & 2021 & 19,134 & 30,499 & 1,330.9 \\
1969 & 7,442 & 11,863 & 1,345.2 & 2022 & 18,863 & 30,067 & 1,348.8 \\
1970 & 7,688 & 12,255 & 1,360.1 &  & & &  \\
\end{longtable}
\endgroup


% tab:appsec1
\begingroup\setlength{\tabcolsep}{2.5pt}\scriptsize
\begin{longtable}{rrrrrrrrrrrr}
\caption[Sectoral volume indices, Estonia 1919--95: primary and secondary sectors]{\textsc{Sectoral volume indices, Estonia 1919--95: primary and secondary sectors}}\label{tab:appsec1}\\
\toprule
Year & Crops & Livestock & Forestry, fishing & Primary & Food & Textiles & Other light & Heavy & Industry & Construction & Secondary \\
\midrule
\endfirsthead
\multicolumn{12}{l}{\textit{Table~\ref{tab:appsec1} (continued)}}\\
\toprule
Year & Crops & Livestock & Forestry, fishing & Primary & Food & Textiles & Other light & Heavy & Industry & Construction & Secondary \\
\midrule
\endhead
\midrule\multicolumn{12}{r}{\textit{continued}}\\
\endfoot
\bottomrule
\multicolumn{12}{p{0.96\linewidth}}{\scriptsize\textit{Sources:} see online appendix. \textit{Notes:} constant-price value-added volume indices, 1938 = 100; `---': not available. Commerce$^{*}$: interwar index of trade, finance and insurance combined.}\\
\endlastfoot
1919 & 60.6 & 41.5 & 60.0 & 50.5 & --- & --- & --- & --- & 20.4 & 17.1 & 20.1 \\
1920 & 62.7 & 58.6 & 72.5 & 61.2 & --- & --- & --- & --- & 26.3 & 17.2 & 25.4 \\
1921 & 77.6 & 66.9 & 86.0 & 72.6 & --- & --- & --- & --- & 35.3 & 17.5 & 33.6 \\
1922 & 74.6 & 72.0 & 87.7 & 74.1 & --- & --- & --- & --- & 42.7 & 17.8 & 40.3 \\
1923 & 72.2 & 74.9 & 93.4 & 75.1 & --- & --- & --- & --- & 55.1 & 61.7 & 55.7 \\
1924 & 71.0 & 78.4 & 108.1 & 77.4 & --- & --- & --- & --- & 54.8 & 18.4 & 51.3 \\
1925 & 83.3 & 87.6 & 116.1 & 87.9 & --- & --- & --- & --- & 53.8 & 36.2 & 52.1 \\
1926 & 71.8 & 91.3 & 95.9 & 83.7 & --- & --- & --- & --- & 57.7 & 39.6 & 56.0 \\
1927 & 76.5 & 93.0 & 106.0 & 87.2 & --- & --- & --- & --- & 54.1 & 15.7 & 50.5 \\
1928 & 64.3 & 93.7 & 96.2 & 82.0 & --- & --- & --- & --- & 55.6 & 51.9 & 55.2 \\
1929 & 78.2 & 87.4 & 94.6 & 84.2 & --- & --- & --- & --- & 60.7 & 51.8 & 59.8 \\
1930 & 100.1 & 92.9 & 85.5 & 95.3 & --- & --- & --- & --- & 63.4 & 46.1 & 61.7 \\
1931 & 79.3 & 96.9 & 80.3 & 88.6 & --- & --- & --- & --- & 62.1 & 39.1 & 59.9 \\
1932 & 79.8 & 96.2 & 111.8 & 90.6 & --- & --- & --- & --- & 55.0 & 37.2 & 53.3 \\
1933 & 92.5 & 94.9 & 111.8 & 95.1 & --- & --- & --- & --- & 61.1 & 36.8 & 58.8 \\
1934 & 104.8 & 93.2 & 100.1 & 98.4 & --- & --- & --- & --- & 75.2 & 38.7 & 71.7 \\
1935 & 90.2 & 93.9 & 98.7 & 92.7 & --- & --- & --- & --- & 78.8 & 55.7 & 76.6 \\
1936 & 87.2 & 89.7 & 100.0 & 89.4 & --- & --- & --- & --- & 98.5 & 72.1 & 96.0 \\
1937 & 101.7 & 100.2 & 98.6 & 100.7 & --- & --- & --- & --- & 104.4 & 102.6 & 104.2 \\
1938 & 100.0 & 100.0 & 100.0 & 100.0 & 100.0 & 100.0 & 100.0 & 100.0 & 100.0 & 100.0 & 100.0 \\
1939 & 101.5 & 99.7 & 96.7 & 100.2 & 104.4 & 110.4 & 96.7 & 89.1 & 107.1 & 111.7 & 107.6 \\
1940 & 104.4 & 85.3 & 83.4 & 92.9 & 101.2 & 113.1 & 85.1 & 81.4 & 108.3 & 102.5 & 107.7 \\
1941 & 76.9 & 62.6 & 82.4 & 68.8 & 78.9 & 64.0 & 75.9 & 67.3 & 83.8 & 88.9 & 86.3 \\
1942 & 77.0 & 48.8 & 81.5 & 61.6 & 61.7 & 22.5 & 66.7 & 60.7 & 64.4 & 80.3 & 69.9 \\
1943 & 64.2 & 47.1 & 80.5 & 55.1 & 55.7 & 30.1 & 65.4 & 57.9 & 63.4 & 69.8 & 66.1 \\
1944 & 57.4 & 39.1 & 78.4 & 47.8 & 42.1 & 20.1 & 65.3 & 34.5 & 15.5 & 19.0 & 16.7 \\
1945 & 78.9 & 44.0 & 78.6 & 59.9 & 45.5 & 32.7 & 68.9 & 51.6 & 61.2 & 87.3 & 69.7 \\
1946 & 86.4 & 47.1 & 90.9 & 65.1 & 67.4 & 47.8 & 84.8 & 58.7 & 78.5 & 95.9 & 84.6 \\
1947 & 118.4 & 49.9 & 104.2 & 81.0 & 78.4 & 73.5 & 99.1 & 80.8 & 99.1 & 117.6 & 105.8 \\
1948 & 131.5 & 55.9 & 102.2 & 89.9 & 104.6 & 96.9 & 112.5 & 103.9 & 122.9 & 136.2 & 128.2 \\
1949 & 117.2 & 59.2 & 107.4 & 85.5 & 117.1 & 132.4 & 127.8 & 141.1 & 150.9 & 162.4 & 156.0 \\
1950 & 113.6 & 62.0 & 111.5 & 85.5 & 130.3 & 172.0 & 138.7 & 182.7 & 179.4 & 187.0 & 183.7 \\
1951 & 106.9 & 62.3 & 114.9 & 82.7 & 150.8 & 232.3 & 140.4 & 223.6 & 210.6 & 209.2 & 212.5 \\
1952 & 105.3 & 66.3 & 107.5 & 83.9 & 163.7 & 269.8 & 136.7 & 266.1 & 232.9 & 231.4 & 235.0 \\
1953 & 97.3 & 64.5 & 110.8 & 79.5 & 179.9 & 304.2 & 155.0 & 297.3 & 261.0 & 252.5 & 261.3 \\
1954 & 88.7 & 69.0 & 131.7 & 78.6 & 203.7 & 326.3 & 170.5 & 322.2 & 285.2 & 269.3 & 283.6 \\
1955 & 86.0 & 67.5 & 165.5 & 77.4 & 219.0 & 379.2 & 207.0 & 380.0 & 332.5 & 320.9 & 332.6 \\
1956 & 109.8 & 70.5 & 177.9 & 89.8 & 229.7 & 382.3 & 223.2 & 456.6 & 364.0 & 356.1 & 365.6 \\
1957 & 107.5 & 79.5 & 174.2 & 93.4 & 236.9 & 466.5 & 236.3 & 493.7 & 401.1 & 387.9 & 401.5 \\
1958 & 116.4 & 82.3 & 177.4 & 99.0 & 257.9 & 527.8 & 251.6 & 571.7 & 448.6 & 436.4 & 449.9 \\
1959 & 127.6 & 92.7 & 186.8 & 109.6 & 295.9 & 555.7 & 275.5 & 637.1 & 492.5 & 482.4 & 494.8 \\
1960 & 138.7 & 102.7 & 178.9 & 119.5 & 337.3 & 583.0 & 285.2 & 710.2 & 533.4 & 516.7 & 534.2 \\
1961 & 134.2 & 99.3 & 165.6 & 115.5 & 345.5 & 597.2 & 297.9 & 845.3 & 582.0 & 580.1 & 587.8 \\
1962 & 109.1 & 99.1 & 193.1 & 104.9 & 363.4 & 626.0 & 324.9 & 913.4 & 623.3 & 630.4 & 632.2 \\
1963 & 132.2 & 97.6 & 242.1 & 115.5 & 424.3 & 634.5 & 350.5 & 992.8 & 672.8 & 676.7 & 681.3 \\
1964 & 152.7 & 98.4 & 273.0 & 125.7 & 467.1 & 661.2 & 370.8 & 1,073.3 & 720.3 & 727.4 & 730.3 \\
1965 & 154.7 & 105.6 & 280.2 & 130.6 & 511.8 & 695.4 & 379.8 & 1,200.5 & 778.6 & 794.5 & 791.9 \\
1966 & 144.7 & 113.3 & 304.1 & 130.8 & 522.7 & 690.4 & 403.8 & 1,241.7 & 801.5 & 841.0 & 822.2 \\
1967 & 158.5 & 116.6 & 316.7 & 138.9 & 540.4 & 795.2 & 429.4 & 1,335.8 & 867.1 & 886.0 & 882.3 \\
1968 & 168.6 & 119.6 & 413.9 & 147.3 & 597.6 & 815.8 & 476.6 & 1,395.7 & 921.8 & 927.5 & 933.6 \\
1969 & 150.4 & 121.5 & 415.3 & 140.3 & 637.5 & 867.5 & 470.0 & 1,492.9 & 968.2 & 949.4 & 973.1 \\
1970 & 162.1 & 124.7 & 423.0 & 147.3 & 689.9 & 878.4 & 471.0 & 1,598.2 & 1,013.5 & 991.8 & 1,018.1 \\
1971 & 165.6 & 131.5 & 453.8 & 153.2 & 752.5 & 929.2 & 481.0 & 1,735.2 & 1,083.2 & 1,047.6 & 1,084.3 \\
1972 & 137.9 & 134.6 & 462.1 & 142.8 & 773.4 & 944.2 & 481.2 & 1,722.9 & 1,088.3 & 1,043.3 & 1,086.7 \\
1973 & 159.4 & 130.8 & 495.3 & 151.1 & 834.8 & 915.4 & 465.3 & 1,777.4 & 1,105.3 & 1,069.6 & 1,106.7 \\
1974 & 152.6 & 144.3 & 534.8 & 156.2 & 888.2 & 825.3 & 499.5 & 1,778.9 & 1,111.6 & 1,097.6 & 1,119.6 \\
1975 & 151.9 & 149.0 & 595.6 & 159.7 & 941.5 & 815.6 & 519.9 & 1,829.1 & 1,144.7 & 1,148.4 & 1,158.4 \\
1976 & 147.7 & 156.6 & 640.8 & 163.0 & 1,031.9 & 826.5 & 540.0 & 1,894.7 & 1,195.8 & 1,165.5 & 1,199.7 \\
1977 & 140.6 & 161.4 & 601.1 & 161.4 & 928.8 & 837.5 & 550.1 & 1,893.5 & 1,177.3 & 1,165.8 & 1,186.7 \\
1978 & 106.4 & 160.5 & 548.7 & 144.5 & 914.3 & 854.7 & 565.4 & 1,895.1 & 1,184.5 & 1,153.0 & 1,188.0 \\
1979 & 127.2 & 157.6 & 531.0 & 151.8 & 908.0 & 872.9 & 578.4 & 1,870.2 & 1,185.5 & 1,100.2 & 1,172.8 \\
1980 & 143.6 & 165.5 & 569.6 & 164.1 & 916.9 & 869.1 & 596.8 & 1,920.3 & 1,208.4 & 1,130.2 & 1,198.1 \\
1981 & 107.0 & 164.0 & 564.0 & 147.0 & 932.7 & 873.3 & 616.5 & 1,831.1 & 1,196.0 & 1,107.2 & 1,182.3 \\
1982 & 130.0 & 149.7 & 555.9 & 149.5 & 898.0 & 885.1 & 604.3 & 1,846.3 & 1,189.7 & 1,083.9 & 1,170.8 \\
1983 & 132.6 & 165.2 & 550.9 & 158.6 & 932.7 & 913.7 & 636.1 & 1,883.7 & 1,228.2 & 1,126.7 & 1,211.1 \\
1984 & 142.6 & 174.5 & 595.5 & 169.0 & 966.5 & 922.2 & 677.2 & 1,863.0 & 1,249.2 & 1,137.9 & 1,229.4 \\
1985 & 112.6 & 179.9 & 571.4 & 158.0 & 941.5 & 952.3 & 686.8 & 1,799.3 & 1,236.1 & 1,117.8 & 1,214.0 \\
1986 & 142.8 & 181.1 & 583.9 & 172.3 & 945.4 & 972.9 & 699.8 & 1,813.3 & 1,251.1 & 1,131.3 & 1,228.7 \\
1987 & 112.0 & 184.6 & 575.9 & 160.3 & 941.9 & 987.5 & 699.2 & 1,815.4 & 1,254.0 & 1,147.2 & 1,235.5 \\
1988 & 99.2 & 183.2 & 586.9 & 154.1 & 957.1 & 986.6 & 735.0 & 1,820.6 & 1,273.5 & 1,153.5 & 1,251.3 \\
1989 & 132.9 & 183.1 & 553.6 & 168.2 & 921.5 & 974.6 & 699.2 & 1,783.2 & 1,236.9 & 1,108.8 & 1,211.9 \\
1990 & 110.2 & 176.6 & 482.3 & 153.0 & 828.2 & 938.6 & 608.8 & 1,602.8 & 1,118.3 & 978.4 & 1,088.4 \\
1991 & 110.9 & 153.0 & 424.4 & 139.5 & 751.7 & 870.6 & 537.1 & 1,489.0 & 1,022.7 & 899.6 & 996.8 \\
1992 & 91.9 & 120.7 & 192.5 & 108.6 & 464.8 & 510.6 & 332.8 & 772.2 & 584.2 & 483.8 & 560.4 \\
1993 & 90.6 & 100.7 & 193.1 & 97.6 & 378.7 & 239.9 & 217.3 & 451.9 & 363.2 & 336.9 & 359.2 \\
1994 & 81.6 & 88.1 & 210.8 & 87.4 & 374.0 & 274.6 & 264.3 & 430.2 & 383.1 & 355.4 & 378.9 \\
1995 & 81.0 & 81.3 & 218.8 & 83.7 & 369.2 & 282.3 & 205.8 & 467.2 & 370.3 & 341.9 & 365.8 \\
\end{longtable}
\endgroup


% tab:appsec2
\begingroup\setlength{\tabcolsep}{2.5pt}\scriptsize
\begin{longtable}{rrrrrrrrrrrrr}
\caption[Sectoral volume indices, Estonia 1919--95: services and total GDP]{\textsc{Sectoral volume indices, Estonia 1919--95: services and total GDP}}\label{tab:appsec2}\\
\toprule
Year & Transport, comm. & Trade & Commerce$^{*}$ & Housing & Finance & Administration & Defence & Health & Education & Other & Tertiary & GDP \\
\midrule
\endfirsthead
\multicolumn{13}{l}{\textit{Table~\ref{tab:appsec2} (continued)}}\\
\toprule
Year & Transport, comm. & Trade & Commerce$^{*}$ & Housing & Finance & Administration & Defence & Health & Education & Other & Tertiary & GDP \\
\midrule
\endhead
\midrule\multicolumn{13}{r}{\textit{continued}}\\
\endfoot
\bottomrule
\multicolumn{13}{p{0.96\linewidth}}{\scriptsize\textit{Sources:} see online appendix. \textit{Notes:} constant-price value-added volume indices, 1938 = 100; `---': not available. Commerce$^{*}$: interwar index of trade, finance and insurance combined.}\\
\endlastfoot
1919 & 30.3 & --- & 32.1 & 46.2 & --- & --- & 470.2 & --- & --- & 93.9 & 81.7 & 56.0 \\
1920 & 41.2 & --- & 39.1 & 46.4 & --- & --- & 118.7 & --- & --- & 94.3 & 57.6 & 51.8 \\
1921 & 50.1 & --- & 48.3 & 47.2 & --- & --- & 100.0 & --- & --- & 96.0 & 62.2 & 59.8 \\
1922 & 60.4 & --- & 55.0 & 47.9 & --- & --- & 100.0 & --- & --- & 97.3 & 66.8 & 63.7 \\
1923 & 63.9 & --- & 67.9 & 48.6 & --- & --- & 100.0 & --- & --- & 98.1 & 72.9 & 69.9 \\
1924 & 66.4 & --- & 64.4 & 49.4 & --- & --- & 100.0 & --- & --- & 98.5 & 72.1 & 69.5 \\
1925 & 65.8 & --- & 70.3 & 50.1 & --- & --- & 100.0 & --- & --- & 98.6 & 75.2 & 74.9 \\
1926 & 71.4 & --- & 71.3 & 50.8 & --- & --- & 100.0 & --- & --- & 98.6 & 76.2 & 74.5 \\
1927 & 70.3 & --- & 73.6 & 51.6 & --- & --- & 100.0 & --- & --- & 98.5 & 77.0 & 75.0 \\
1928 & 76.2 & --- & 78.9 & 53.4 & --- & --- & 100.0 & --- & --- & 98.5 & 79.6 & 75.1 \\
1929 & 83.3 & --- & 78.5 & 55.2 & --- & --- & 100.0 & --- & --- & 98.5 & 81.1 & 77.5 \\
1930 & 88.2 & --- & 80.1 & 69.0 & --- & --- & 100.0 & --- & --- & 98.5 & 85.1 & 83.8 \\
1931 & 86.2 & --- & 71.4 & 86.2 & --- & --- & 100.0 & --- & --- & 98.8 & 83.2 & 80.0 \\
1932 & 79.0 & --- & 62.3 & 86.2 & --- & --- & 100.0 & --- & --- & 99.0 & 78.9 & 77.7 \\
1933 & 71.6 & --- & 66.1 & 82.8 & --- & --- & 100.0 & --- & --- & 99.3 & 79.5 & 80.8 \\
1934 & 74.7 & --- & 79.1 & 79.3 & --- & --- & 100.0 & --- & --- & 99.5 & 85.0 & 87.2 \\
1935 & 84.9 & --- & 83.0 & 86.2 & --- & --- & 100.0 & --- & --- & 99.7 & 88.4 & 87.4 \\
1936 & 89.6 & --- & 92.0 & 93.1 & --- & --- & 100.0 & --- & --- & 99.8 & 93.8 & 92.6 \\
1937 & 95.4 & --- & 102.2 & 96.6 & --- & --- & 100.0 & --- & --- & 99.8 & 100.0 & 101.2 \\
1938 & 100.0 & 100.0 & 100.0 & 100.0 & 100.0 & 100.0 & 100.0 & 100.0 & 100.0 & 100.0 & 100.0 & 100.0 \\
1939 & 112.0 & 100.3 & 100.3 & 100.6 & 100.3 & 101.5 & 100.0 & 107.1 & 91.1 & 99.6 & 102.1 & 102.6 \\
1940 & 145.9 & 91.9 & 91.9 & 97.8 & 91.9 & 99.0 & 178.6 & 99.6 & 78.0 & 96.4 & 108.8 & 102.5 \\
1941 & 151.2 & 81.3 & --- & 47.1 & 81.6 & 113.9 & 180.9 & 73.3 & 81.5 & 96.4 & 96.2 & 82.7 \\
1942 & 130.8 & 79.2 & --- & 69.9 & 79.8 & 128.7 & 183.1 & 82.5 & 85.1 & 96.4 & 103.2 & 79.1 \\
1943 & 125.7 & 65.2 & --- & 66.7 & 77.9 & 143.6 & 185.4 & 88.8 & 88.6 & 96.4 & 103.4 & 75.4 \\
1944 & 101.5 & 52.6 & --- & 63.4 & 9.0 & 93.4 & 165.9 & 88.8 & 63.1 & 42.3 & 87.1 & 57.8 \\
1945 & 121.5 & 71.7 & --- & 60.2 & 66.7 & 173.3 & 189.8 & 95.0 & 92.2 & 96.4 & 105.5 & 79.1 \\
1946 & 183.1 & 81.3 & --- & 70.9 & 71.1 & 176.3 & 193.1 & 129.2 & 104.6 & 131.2 & 115.5 & 87.7 \\
1947 & 219.8 & 88.9 & --- & 81.6 & 73.7 & 179.2 & 196.4 & 151.9 & 117.0 & 166.1 & 124.0 & 102.0 \\
1948 & 211.1 & 96.6 & --- & 92.3 & 76.2 & 182.2 & 199.6 & 176.7 & 125.8 & 200.9 & 130.8 & 112.5 \\
1949 & 223.7 & 104.3 & --- & 103.0 & 78.8 & 185.2 & 202.9 & 196.4 & 134.7 & 235.8 & 138.1 & 117.6 \\
1950 & 252.5 & 114.9 & --- & 113.7 & 83.5 & 188.2 & 206.1 & 199.3 & 159.0 & 270.6 & 147.0 & 125.4 \\
1951 & 273.7 & 132.3 & --- & 122.9 & 103.6 & 177.8 & 194.7 & 226.8 & 162.0 & 301.5 & 150.0 & 129.9 \\
1952 & 301.3 & 147.0 & --- & 125.4 & 113.8 & 167.4 & 183.3 & 237.6 & 166.5 & 316.9 & 150.2 & 134.3 \\
1953 & 271.2 & 179.2 & --- & 130.4 & 138.1 & 157.0 & 171.9 & 256.0 & 170.2 & 320.0 & 151.1 & 136.8 \\
1954 & 297.1 & 209.5 & --- & 131.3 & 144.9 & 146.6 & 160.6 & 280.5 & 184.3 & 353.9 & 154.0 & 141.0 \\
1955 & 341.0 & 216.9 & --- & 137.1 & 170.5 & 136.2 & 149.2 & 302.0 & 187.3 & 370.1 & 155.7 & 149.1 \\
1956 & 372.2 & 233.4 & --- & 136.3 & 194.0 & 133.7 & 146.5 & 315.7 & 201.4 & 388.6 & 159.9 & 162.1 \\
1957 & 405.5 & 268.3 & --- & 142.1 & 205.6 & 131.8 & 144.4 & 341.8 & 209.6 & 428.7 & 167.3 & 172.4 \\
1958 & 431.9 & 284.0 & --- & 152.2 & 209.6 & 130.0 & 142.4 & 347.9 & 217.0 & 492.7 & 173.3 & 185.2 \\
1959 & 466.2 & 306.0 & --- & 143.8 & 209.6 & 128.1 & 140.3 & 364.8 & 223.7 & 545.1 & 175.7 & 198.6 \\
1960 & 517.8 & 341.9 & --- & 148.8 & 196.9 & 126.3 & 138.3 & 383.2 & 229.6 & 634.9 & 183.1 & 212.6 \\
1961 & 558.8 & 354.7 & --- & 160.5 & 198.9 & 129.0 & 141.3 & 403.1 & 244.5 & 670.8 & 192.4 & 222.8 \\
1962 & 617.8 & 368.5 & --- & 161.4 & 197.2 & 131.7 & 144.3 & 412.3 & 253.4 & 704.7 & 198.3 & 227.2 \\
1963 & 685.0 & 382.3 & --- & 163.0 & 200.2 & 134.4 & 147.3 & 416.9 & 260.1 & 731.7 & 204.2 & 242.5 \\
1964 & 777.1 & 402.5 & --- & 167.2 & 210.7 & 137.2 & 150.2 & 432.2 & 270.5 & 817.3 & 214.0 & 259.1 \\
1965 & 933.4 & 440.2 & --- & 173.9 & 225.4 & 139.9 & 153.2 & 439.9 & 285.4 & 834.2 & 226.6 & 276.2 \\
1966 & 981.3 & 485.2 & --- & 175.6 & 247.7 & 146.1 & 160.0 & 446.0 & 296.5 & 956.8 & 238.3 & 285.5 \\
1967 & 1,103.4 & 540.4 & --- & 183.1 & 278.5 & 154.7 & 169.5 & 446.0 & 311.4 & 1,009.3 & 254.2 & 305.1 \\
1968 & 1,169.2 & 601.9 & --- & 186.5 & 304.9 & 163.4 & 179.0 & 475.2 & 321.0 & 1,022.4 & 268.0 & 322.7 \\
1969 & 1,218.2 & 650.7 & --- & 189.8 & 328.1 & 170.8 & 187.1 & 489.0 & 334.4 & 1,075.6 & 280.2 & 330.3 \\
1970 & 1,303.4 & 704.0 & --- & 189.0 & 355.0 & 175.8 & 192.6 & 495.1 & 341.9 & 1,125.7 & 290.9 & 344.9 \\
1971 & 1,354.3 & 744.4 & --- & 205.7 & 378.7 & 182.0 & 199.3 & 505.8 & 352.3 & 1,134.9 & 304.1 & 363.5 \\
1972 & 1,379.9 & 786.7 & --- & 215.7 & 414.7 & 184.4 & 202.0 & 516.5 & 364.1 & 1,169.6 & 314.7 & 362.7 \\
1973 & 1,519.8 & 819.7 & --- & 224.1 & 435.0 & 189.4 & 207.5 & 528.8 & 369.3 & 1,220.5 & 327.6 & 374.8 \\
1974 & 1,683.6 & 851.0 & --- & 230.8 & 452.7 & 195.6 & 214.3 & 538.0 & 376.8 & 1,289.9 & 341.4 & 384.4 \\
1975 & 1,669.5 & 902.5 & --- & 238.3 & 500.4 & 204.3 & 223.7 & 545.7 & 382.7 & 1,373.2 & 353.5 & 397.0 \\
1976 & 1,730.6 & 961.3 & --- & 237.5 & 529.0 & 212.9 & 233.2 & 553.3 & 390.9 & 1,394.8 & 365.0 & 409.5 \\
1977 & 1,842.7 & 1,001.7 & --- & 237.5 & 577.2 & 221.6 & 242.7 & 556.4 & 413.2 & 1,388.6 & 378.2 & 411.5 \\
1978 & 1,901.7 & 1,035.7 & --- & 242.5 & 601.1 & 230.2 & 252.2 & 567.1 & 426.6 & 1,475.0 & 390.9 & 408.2 \\
1979 & 1,883.1 & 1,085.3 & --- & 250.0 & 556.5 & 238.9 & 261.7 & 585.5 & 437.7 & 1,521.3 & 400.8 & 412.8 \\
1980 & 1,880.2 & 1,159.8 & --- & 258.4 & 592.7 & 247.6 & 271.2 & 597.8 & 450.4 & 1,549.8 & 414.8 & 428.1 \\
1981 & 1,917.0 & 1,214.9 & --- & 263.4 & 628.0 & 252.5 & 276.6 & 609.6 & 459.6 & 1,572.9 & 425.9 & 421.3 \\
1982 & 1,956.3 & 1,184.6 & --- & 274.2 & 621.8 & 260.0 & 284.8 & 630.2 & 465.5 & 1,581.6 & 431.8 & 422.8 \\
1983 & 2,024.2 & 1,212.2 & --- & 287.6 & 634.4 & 257.5 & 282.1 & 644.9 & 468.8 & 1,595.2 & 439.8 & 436.8 \\
1984 & 2,099.9 & 1,258.1 & --- & 280.9 & 658.5 & 258.7 & 283.4 & 664.0 & 473.5 & 1,647.1 & 447.1 & 447.5 \\
1985 & 2,098.9 & 1,326.1 & --- & 283.4 & 688.1 & 230.2 & 252.2 & 687.6 & 478.7 & 1,876.5 & 448.7 & 440.2 \\
1986 & 2,233.4 & 1,352.6 & --- & 289.3 & 673.6 & 231.5 & 253.6 & 686.1 & 488.7 & 1,963.1 & 458.6 & 453.2 \\
1987 & 2,527.3 & 1,379.2 & --- & 286.8 & 719.3 & 231.5 & 253.6 & 695.0 & 501.2 & 1,819.6 & 468.9 & 452.3 \\
1988 & 2,893.7 & 1,472.0 & --- & 280.9 & 724.3 & 231.5 & 253.6 & 678.8 & 510.4 & 1,759.8 & 484.2 & 457.5 \\
1989 & 3,162.3 & 1,564.8 & --- & 280.1 & 817.3 & 212.9 & 233.2 & 730.3 & 519.0 & 1,868.6 & 499.4 & 463.4 \\
1990 & 2,988.0 & 1,670.9 & --- & 273.4 & 932.8 & 211.7 & 231.9 & 721.5 & 513.7 & 1,757.9 & 500.1 & 436.1 \\
1991 & 2,903.8 & 951.7 & --- & 255.0 & 581.9 & 160.4 & 220.3 & 731.7 & 502.4 & 1,697.2 & 423.6 & 386.5 \\
1992 & 2,346.3 & 407.3 & --- & 234.1 & 287.0 & 166.6 & 164.2 & 704.0 & 509.6 & 1,679.7 & 338.6 & 268.8 \\
1993 & 2,573.3 & 583.3 & --- & 231.6 & 420.9 & 179.4 & 111.3 & 689.4 & 514.8 & 1,676.5 & 346.3 & 233.3 \\
1994 & 2,714.6 & 554.3 & --- & 250.0 & 439.3 & 187.1 & 94.3 & 685.0 & 499.2 & 1,641.4 & 346.2 & 231.5 \\
1995 & 1,486.7 & 661.8 & --- & 269.2 & 499.9 & 183.5 & 86.4 & 532.0 & 578.0 & 1,842.1 & 327.4 & 220.7 \\
\end{longtable}
\endgroup



\end{document}

